% संपूर्ण मराठी ग्रंथातील प्रकरण 79: मूळ संकलनमार्गाबाहेरील विभाग
% हा संपादनयोग्य स्रोतखंड आहे. संकलनासाठी संपूर्ण स्रोतसंग्रहातील
% अक्षररूपे, पूर्वभाग, चिन्हव्याख्या आणि आकृती-स्रोत आवश्यक आहेत.
% मूळ: Open Logic Project; मजकूर आणि रूपांतर CC BY 4.0.
% यंत्रानुवाद, दुरुस्ती आणि तपासणी: OpenAI Codex —
% GPT-5.6 Sol आणि GPT-6 Sol, दोन्ही Ultra effort.
% दुरुस्ती-पुनरावलोकन: GPT-6.1 Sol, Ultra effort.
\chapter{मूळ संकलनमार्गाबाहेरील विभाग}\label{supp:chap}
या विभागांचे स्रोत-एककांमध्ये भाषांतर केलेले आहे. मुख्य संकलनमार्गात नसलेली सामग्री येथे विषयानुसार पूरक स्वरूपात दिली आहे.












\section{निष्पन्नता चे गुणधर्म}\label{fol:axd:prv-alt:sec}

\begin{prop}[एकस्वनिकता] जर $\Gamma \subseteq \Delta$ आणि $\Gamma \Proves \varphi$
असेल, तर $\Delta \Proves \varphi$ सुद्धा. \end{prop}

\begin{proof} कोणताही परिमित $\Gamma_0 \subseteq \Gamma$ हा~$\Delta$ चासुद्धा
परिमित उपसंच असतो. म्हणून $\Gamma_0 \Proves \varphi$ ची निष्पत्ती
$\Delta \Proves \varphi$ हेही दाखवते. \end{proof}

\begin{prop}\label{fol:axd:prv-alt:prop:provability} $\Gamma \Proves \varphi$ तेव्हा आणि केवळ
तेव्हाच, जेव्हा $\Gamma\cup \{\lnot\varphi \}$ विसंगत असतो. \end{prop}

\begin{prop} \begin{enumerate}
\item \label{fol:axd:prv-alt:prop:provability-contr} जर $\Gamma \Proves \varphi$ आणि $\Gamma
\cup \{ \varphi\} \Proves \lfalse$ असेल, तर $\Gamma$ विसंगत असतो.

\item \label{fol:axd:prv-alt:prop:provability-lnot} जर $\Gamma \cup \{\varphi\}
\Proves \lfalse$ असेल, तर $\Gamma \Proves \lnot \varphi$.

\item \label{fol:axd:prv-alt:prop:provability-exhaustive} जर $\Gamma \cup \{\varphi\}
\Proves \lfalse$ आणि $\Gamma \cup \{\lnot \varphi\} \Proves
\lfalse$ असेल, तर $\Gamma \Proves \lfalse$.

\item \label{fol:axd:prv-alt:prop:provability-lor-left} जर $\Gamma \cup \{\varphi\}
\Proves \lfalse$ आणि $\Gamma \cup \{\psi\} \Proves
\lfalse$ असेल, तर $\Gamma \cup \{\varphi \lor \psi\} \Proves \lfalse$.

\item \label{fol:axd:prv-alt:prop:provability-lor-right} जर $\Gamma
\Proves \varphi$ किंवा $\Gamma \Proves \psi$ असेल, तर $\Gamma
\Proves \varphi \lor \psi$.

\item \label{fol:axd:prv-alt:prop:provability-land-left} जर $\Gamma
\Proves \varphi \land \psi$ असेल, तर $\Gamma \Proves \varphi$ आणि
$\Gamma \Proves \psi$.

\item \label{fol:axd:prv-alt:prop:provability-land-right} जर $\Gamma
\Proves \varphi$ आणि $\Gamma \Proves \psi$ असेल, तर $\Gamma
\Proves \varphi \land \psi$.

\item \label{fol:axd:prv-alt:prop:provability-mp} जर $\Gamma \Proves
\varphi$ आणि $\Gamma \Proves \varphi \lif \psi$ असेल, तर $\Gamma
\Proves \psi$.

\item \label{fol:axd:prv-alt:prop:provability-lif} जर $\Gamma \Proves
\lnot \varphi$ किंवा $\Gamma \Proves \psi$ असेल, तर $\Gamma \Proves
\varphi \lif \psi$. \end{enumerate} \end{prop}

\begin{proof}
\begin{enumerate}
\item $\Gamma_0 \cup \{\varphi\} \Proves \lfalse$ असे गृहीत धरा.

\begin{derivation}
1. & $\Gamma_0 \cup \{\varphi\} \Proves \lfalse$ & गृहीत \\
2. & $\Gamma_1 \Proves \varphi$ & गृहीत\\
3. & $\Gamma_0 \cup \Gamma_1 \Proves \lfalse$ & कट नियम \\
\end{derivation}

कारण $\Gamma_0 \subseteq \Gamma$ आणि $\Gamma_1 \subseteq \Gamma$,
$\Gamma_0 \cup \Gamma_1 \subseteq \Gamma$; म्हणून $\Gamma \Proves \lfalse$.

\item $\Gamma_0 \cup\{\varphi\} \Proves \lfalse$ असे गृहीत धरा.

\begin{derivation}
1. & $\Gamma_0 \cup\{\varphi\} \Proves \lfalse$ & गृहीत \\
2. & $\Gamma_0 \Proves \varphi \lif \lfalse$ & निगमन प्रमेय\\
3. & $\Gamma_0 \Proves (\varphi \lif \lfalse) \lif \lnot \varphi$ & Ax0\\
4. & $\Gamma_0 \Proves \lnot \varphi$ & MP 2, 3 \\
\end{derivation}


\item $\Gamma_0 \cup \{ \varphi \} \Proves \lfalse$ आणि
$\Gamma_1 \cup \{\lnot \varphi \} \Proves \lfalse$ असे गृहीत धरा.


\begin{derivation}
1. & $\Gamma_0 \cup\{\varphi\} \Proves \lfalse$ & गृहीत \\
2. & $\Gamma_0 \Proves \varphi \lif \lfalse$ & निगमन प्रमेय \\
3. & $\Gamma_1 \cup \{\lnot \varphi \} \Proves \lfalse$ & गृहीत\\
4. & $\Gamma_0 \Proves (\varphi \lif \lfalse) \lif \lnot \varphi$ & Ax0\\
5. & $\Gamma_0 \Proves \lnot \varphi$ & MP 2, 4\\
6. & $\Gamma_0 \cup \Gamma_1 \Proves \lfalse$ & कट नियम 3, 5 \\
\end{derivation}

कारण $\Gamma_0 \subseteq \Gamma$ आणि $\Gamma_1 \subseteq \Gamma$,
$\Gamma_0 \cup \Gamma_1 \subseteq \Gamma$. म्हणून $\Gamma \Proves \lfalse$.


\item
सराव.


\item
$\Gamma_0 \Proves \varphi$ असे गृहीत धरा.

\begin{derivation}
1. & $\Gamma_0 \Proves \varphi$ & गृहीत \\
2. & $\Gamma_0 \Proves \varphi \lif (\varphi \lor \psi)$ & Ax0\\
3. & $\Gamma_0 \Proves \varphi \lor \psi$ & MP 1, 2 \\
\end{derivation}

म्हणून $\Gamma \Proves \varphi \lor \psi$. $\Gamma \Proves \psi$ असल्यास
सिद्धताही याचप्रमाणे मिळते.


\item
$\Gamma_0 \Proves \varphi \land \psi$ असे गृहीत धरा

\begin{derivation}
1. & $\Gamma_0 \Proves \varphi \land \psi$ & गृहीत \\
2. & $\Gamma_0 \Proves (\varphi\land \psi) \lif \varphi $ & Ax0\\
3. & $\Gamma_0 \Proves \varphi$ & MP 1, 2 \\
\end{derivation}


म्हणून $\Gamma \Proves \varphi$. याचप्रमाणे सुरू होणारे निष्पत्ती
$\Gamma \Proves \psi$ हे दाखवते.


\item सराव.


\item   $\Gamma_0 \Proves \varphi$ आणि $\Gamma_0 \Proves \varphi \lif \psi $ ही दोन्ही गृहीत धरा.

\begin{derivation}
1. & $\Gamma_0 \Proves \varphi$ & गृहीत \\
2. & $\Gamma_0 \Proves \varphi \lif \psi $ & गृहीत\\
3. & $\Gamma_0 \Proves \psi$ & MP 1, 2\\
\end{derivation}

कारण $\Gamma_0 \subseteq \Gamma$, यावरून $\Gamma
\Proves \psi$ मिळते.



\item प्रथम $\Gamma \Proves \lnot \varphi$ असे माना.

\begin{derivation}
1. & $\Gamma_0 \Proves \lnot \varphi$ & गृहीत \\
2. & $\Gamma_0 \Proves \lnot \varphi \lif (\varphi \lif \psi) $ & Ax0\\
3. & $\Gamma_0 \Proves \varphi \lif \psi$ & MP 1,
2\\ \end{derivation}

$\Gamma \Proves \psi$ यासाठीही सिद्धता याचप्रमाणे आहे:

\begin{derivation}
1. & $\Gamma_0 \Proves \psi$ & गृहीत \\
2. & $\Gamma_0 \Proves \psi \lif (\varphi \lif \psi) $ & Ax0\\
3. & $\Gamma_0 \Proves \varphi \lif \psi$ & MP 1, 2\\
\end{derivation}

\end{enumerate}
\end{proof}



\begin{thm}[व्यापकीकरण] \label{fol:axd:prv-alt:thm:Generalization} जर $\Gamma \Proves
\varphi$ असेल आणि $x$ हे $\Gamma$ मधील कोणत्याही सूत्रामध्ये मुक्त नसेल,
तर $\Gamma \Proves \lforall[x][\varphi]$.
\end{thm}

\begin{proof} $\mathsf{Thm}_\Gamma$ वर विगमन वापरा: जर $\varphi$ स्वयंसिद्धक
असेल, तर $\lforall[x][\varphi]$ हेसुद्धा स्वयंसिद्धक असते. जर $\varphi\in \Gamma$
असेल, तर $x$ हे $\varphi$ मध्ये मुक्त नसते. म्हणून $\varphi \lif \lforall[x][\varphi]$
हे स्वयंसिद्धक (\textbf{Ax3}) आहे आणि विध्यात्मक अनुमानाने $\Gamma
\Proves \lforall[x][\varphi]$. आता $\Gamma \Proves \psi$ आणि $\Gamma \Proves
\psi \lif \varphi$ यांमुळे विध्यात्मक अनुमानाने $\varphi$ मिळते असे माना. विगमन
गृहीतकानुसार $\Gamma \Proves \lforall[x][\psi]$ आणि $\Gamma \Proves
\lforall[x][( \psi \lif \varphi)]$. परंतु \textbf{Ax2} मुळे पुढीलही मिळते:
\[ \Gamma \Proves
\lforall[x][( \psi \lif \varphi)] \lif (\lforall[x][ \psi] \lif \lforall[x][\varphi]), \]
आणि विध्यात्मक अनुमान दोनदा वापरल्यावर अपेक्षेप्रमाणे $\Gamma \Proves
\lforall[x][\varphi]$ मिळते. \end{proof}

\begin{thm}\label{fol:axd:prv-alt:thm:WeakGen} (\emph{स्थिरांकांवरील दुर्बल व्यापकीकरण})
जर $\Gamma \Proves \varphi$ असेल आणि $c$ हा $\Gamma$ मध्ये न येणारा
स्थिरांक असेल, तर $\varphi$ मध्ये न येणारा असा चल $x$ असतो की
$\Gamma \Proves \lforall[x][\Subst{\varphi}{x}{c}]$; आणि सिद्धतेत $c$ येत नाही.
\end{thm}

\begin{proof} $\varphi_1,\ldots,\varphi_n$ ही $\Gamma$ पासून $\varphi$ ची सिद्धता
असू द्या; म्हणजे $\varphi=\varphi_n$. $\varphi_1,\ldots,\varphi_n$ यांमध्ये न येणारा चल $x$
निवडा आणि पुढील नवा अनुक्रम पाहा:
$\Subst{\varphi_1}{x}{c},\ldots,\Subst{\varphi_n}{x}{c}.$ हा अनुक्रम
$\Gamma$ पासून $\Subst{\varphi}{x}{c}$ ची सिद्धता असतो. खरेच, प्रत्येक
$i=1,\ldots,n$ साठी: \begin{itemize}
\item जर $\varphi_i$ स्वयंसिद्धक असेल, तर $\Subst{\varphi_i}{x}{c}$ हेसुद्धा स्वयंसिद्धक असते;
\item जर $\varphi_i \in \Gamma$ असेल, तर $c$ हे $\Gamma$ मध्ये नसल्यामुळे
$\Subst{\varphi_i}{x}{c} = \varphi_i \in \Gamma$;
\item जर $\varphi_i$ हे $\varphi_j$ आणि $\varphi_j \lif \varphi_i$ यांपासून मिळाले असेल,
तर $\Subst{\varphi_i}{x}{c}$ हे $\Subst{\varphi_j}{x}{c}$ आणि
$\Subst{(\varphi_j \lif \varphi_i)}{x}{c}
= \Subst{\varphi_j}{x}{c} \lif \Subst{\varphi_i}{x}{c}$ यांपासून विध्यात्मक
अनुमानाने मिळते. \end{itemize}
नव्या अनुक्रमात स्थिरांक $c$ उरलेला नाही, हे स्पष्ट आहे. आता
$\Subst{\varphi}{x}{c}$ च्या निष्पत्तीमध्ये येणाऱ्या $\Gamma$ मधील
सूत्रांचा संच $\Gamma'$ असू द्या. मग $\Gamma'$ मधील कोणत्याही
सूत्रामध्ये $x$ मुक्त नाही आणि $\Gamma' \Proves \Subst{\varphi}{x}{c}$.
व्यापकीकरणाने (\ref{fol:axd:prv-alt:thm:Generalization}) $\Gamma'
\Proves \lforall[x][\Subst{\varphi}{x}{c}]$; म्हणून एकस्वनिकतेने
अपेक्षेप्रमाणे $\Gamma \Proves
\lforall[x][\Subst{\varphi}{x}{c}]$. \end{proof}


\begin{explain}
\ref{fol:axd:prv-alt:thm:WeakGen} मधील $x$ हा $\varphi$ मध्ये (अजिबात) येत नाही, ही अट
कमकुवत करून $x$ हा $\varphi$ मध्ये मुक्त नाही आणि $\varphi$ मध्ये $c$ च्या जागी
बसवण्यासाठी आदेशनासाठी मुक्त आहे, इतक्याच अटी ठेवणे हे आपले उद्दिष्ट आहे.
बद्ध चल बदलल्यास हे करता येते; म्हणून प्रथम तो बदल सिद्ध करू.
\end{explain}

\begin{lem}
\label{fol:axd:prv-alt:lem:free-for}
जर $x$ हा $\varphi$ मध्ये $c$ च्या जागी बसवण्यासाठी मुक्त असेल आणि $y$
हा $\varphi$ मध्ये मुक्त नसेल, तर $x$ हा $\Subst{\varphi}{y}{c}$ मध्ये
$y$ च्या जागी बसवण्यासाठी आदेशनासाठी मुक्त असतो.
\end{lem}

\begin{proof} जर $x$ हा $\Subst{\varphi}{y}{c}$ मध्ये $y$ च्या जागी
बसवण्यासाठी आदेशनासाठी मुक्त नसेल, तर $\Subst{\varphi}{y}{c}$ मधील $y$ चा एखादा
मुक्त उद्भव $\lforall[x]$ या संख्यापकाच्या कार्यक्षेत्रात येतो. पण
गृहीतकानुसार $y$ हा $\varphi$ मध्ये मुक्त नाही; म्हणून असे सर्व उद्भव
$\Subst{}{y}{c}$ या आदेशनातूनच येतात. अशा रीतीने $c$ चा एखादा
उद्भव $\lforall[x]$ च्या कार्यक्षेत्रात येतो आणि $x$ हा $\varphi$ मध्ये
$c$ च्या जागी बसवण्यासाठी आदेशनासाठी मुक्त राहत नाही.
\end{proof}


\begin{lem}[बद्ध चराचा बदल]
\label{fol:axd:prv-alt:lem:ChangeBdVar}
जर $x$ आणि $y$ हे $\varphi$ मध्ये मुक्त नसतील आणि $\varphi$ मध्ये $c$ च्या जागी
बसवण्यासाठी दोन्ही आदेशनासाठी मुक्त असतील, तर
$\Proves \lforall[x][\Subst{\varphi}{x}{c}] \equiv
\lforall[y][\Subst{\varphi}{y}{c}]$ (आणि सिद्धतेत $c$ येत नाही).
\end{lem}

\begin{proof}
गृहीतकांनुसार $y$ हा $\varphi$ मध्ये $c$ च्या जागी बसवण्यासाठी आदेशनासाठी मुक्त
आहे आणि $x$ हा $\varphi$ मध्ये मुक्त नाही. म्हणून मागील
\ref{fol:axd:prv-alt:lem:free-for} वरून $y$ हा $\Subst{\varphi}{x}{c}$ मध्ये $x$ च्या
जागी बसवण्यासाठी आदेशनासाठी मुक्त आहे. त्यामुळे
$\lforall[x][\Subst{\varphi}{x}{c}] \lif \Subst{\varphi}{x}{c} \Subst{}{y}{x}$
हे स्वयंसिद्धक (\textbf{Ax1}) आहे. $x$ हा $\varphi$ मध्ये मुक्त नसल्यामुळे
$\Subst{\varphi}{x}{c} \Subst{}{y}{x} = \Subst{\varphi}{y}{c}$; म्हणून
\[
\Proves \lforall[x][\Subst{\varphi}{x}{c}] \lif \Subst{\varphi}{y}{c},
\]
आणि निगमन प्रमेयाने $\lforall[x][\Subst{\varphi}{x}{c}] \Proves
\Subst{\varphi}{y}{c}$. $y$ हा $\varphi$ मध्ये मुक्त नसल्यामुळे तो
$\lforall[x][\Subst{\varphi}{x}{c}]$ मध्येही मुक्त नाही. त्यामुळे
व्यापकीकरणाने $\lforall[x][\Subst{\varphi}{x}{c}] \Proves
\lforall[y][\Subst{\varphi}{y}{c}]$; पुन्हा निगमन प्रमेयाने
$\Proves \lforall[x][\Subst{\varphi}{x}{c}] \lif
\lforall[y][\Subst{\varphi}{y}{c}]$. उलट दिशेची
$\Proves \lforall[y][\Subst{\varphi}{y}{c}] \lif \lforall[x]
[\Subst{\varphi}{x}{c}]$ ही सिद्धता पूर्णपणे सममित आहे; त्यामुळे
निष्कर्ष Taut वापरून मिळतो.
\end{proof}

\begin{thm}[स्थिरांकांवरील प्रबळ व्यापकीकरण]
\label{fol:axd:prv-alt:thm:StrongGen}
जर $\Gamma \Proves \varphi$ असेल, स्थिरांक $c$ हा $\Gamma$ मध्ये येत
नसेल, आणि $x$ हा $\varphi$ मध्ये मुक्त नसेल पण $\varphi$ मध्ये $c$ च्या जागी
बसवण्यासाठी आदेशनासाठी मुक्त असेल, तर
$\Gamma \Proves \lforall[x][\Subst{\varphi}{x}{c}]$.
\end{thm}

\begin{proof}
$\Gamma \Proves \varphi$ आणि $c$ हा $\Gamma$ मध्ये नाही; म्हणून दुर्बल
व्यापकीकरणाने $\varphi$ मध्ये न येणारा चल $y$ असतो की
$\Gamma \Proves \lforall[y][\Subst{\varphi}{y}{c}]$. $y$ हा $\varphi$ मध्ये
(अजिबात) येत नसल्यामुळे तो $\varphi$ मध्ये मुक्तही नसतो आणि $\varphi$ मध्ये
$c$ च्या जागी बसवण्यासाठी आदेशनासाठी मुक्त असतो. गृहीतकानुसार $x$ हाही
$\varphi$ मध्ये मुक्त नाही आणि $\varphi$ मध्ये $c$ च्या जागी बसवण्यासाठी आदेशनासाठी मुक्त
आहे; म्हणून बद्ध चल बदलण्याच्या अटी पूर्ण होतात. त्यामुळे
$\Proves \lforall[x][\Subst{\varphi}{x}{c}] \equiv
\lforall[y][\Subst{\varphi}{y}{c}]$; म्हणून
$\Gamma \Proves \lforall[x][\Subst{\varphi}{x}{c}]$.
\end{proof}

\begin{thm}
\label{fol:axd:prv-alt:thm:provability-quantifiers}
\begin{enumerate}
\item जर $\Gamma \Proves \varphi(t)$ असेल, तर $\Gamma \Proves
  \lexists[x][\varphi(x)]$.
\item जर $\Gamma \Proves \lforall[x][\varphi(x)]$ असेल, तर $\Gamma
  \Proves \varphi(t)$.
\end{enumerate}
\end{thm}

\begin{proof}
\begin{enumerate}

\item सराव.
\item असे गृहीत धरा की  $\Gamma_0 \Proves \lforall[x][\varphi(x)]$.

\begin{derivation}
1. & $\Gamma_0 \Proves \lforall[x][\varphi(x)]$ & गृहीत \\
2. & $\Gamma_0 \Proves \lforall[x][\varphi(x)] \lif \Subst{\varphi}{t}{x}$ & Ax1 \\
3. & $\Gamma_0 \Proves \Subst{\varphi}{t}{x}$ & MP 1, 2 \\
\end{derivation}

\end{enumerate}
\end{proof}



\begin{quote}\small\textbf{स्रोतनोंद.} SOL6-A285: निषेधाच्या गुणधर्माची सिद्धता आता त्याच्या स्वतःच्या गृहीतकापासून not-A काढते. बद्ध चर बदलातील सार्विक संख्यापक फक्त अभिव्यंजनाच्या पूर्वांगावर आहे; मूळ संपूर्ण अभिव्यंजनावरील व्याप्ती योग्य आदेशन-स्वयंसिद्धक नव्हती.\end{quote}
\begin{quote}\small\textbf{स्रोतनोंद.} OLINC-321: मूळ हिशोब निषेधाच्या विधानाऐवजी आधीच्या समतुल्यतेची उलट दिशा सिद्ध करत होता. SOL6-A285 मध्ये A या अतिरिक्त गृहीतकापासून असत्य मिळते या आधारावर not-A काढला आहे.\end{quote}
\begin{quote}\small\textbf{स्रोतनोंद.} OLINC-326: पहिल्या गृहीतकातील आधारसंच पुढील हिशोबाप्रमाणे Gamma\_0 केला आहे.\end{quote}
\begin{quote}\small\textbf{स्रोतनोंद.} OLINC-322: A आणि B यांच्या संयुक्तातून A काढणाऱ्या हिशोबात स्रोताच्या तिसऱ्या ओळीत चुकीने A किंवा B आहे.\end{quote}
\begin{quote}\small\textbf{स्रोतनोंद.} OLINC-323: या हिशोबात Gamma\_1 परिभाषित नाही; वापरलेला परिमित आधारसंच Gamma\_0 आहे.\end{quote}
\begin{quote}\small\textbf{स्रोतनोंद.} OLINC-325: सिद्धता-अनुक्रमाच्या प्रत्येक टप्प्यातील स्वयंसिद्धक A\_i आहे; स्रोताने A म्हटले आहे.\end{quote}
\begin{quote}\small\textbf{स्रोतनोंद.} OLINC-324: पहिल्या आदेशनात स्रोताने c ऐवजी x लिहिला आहे; उरलेल्या विधानानुसार c केला आहे.\end{quote}











\section{वाक्य यांचे महत्तम सुसंगत संच}\label{fol:com:mcs:sec}

\begin{defn}[महत्तम सुसंगत संच]
\label{fol:com:mcs:mcs}
वाक्यांचा संच~$\Gamma$ \emph{महत्तम सुसंगत} तेव्हा आणि केवळ तेव्हाच असतो, जेव्हा
\begin{enumerate}
\item $\Gamma$ सुसंगत असतो, आणि
\item जर $\Gamma \subsetneq \Gamma'$ असेल, तर $\Gamma'$ विसंगत असतो.
\end{enumerate}
\end{defn}

\begin{explain}
वरील व्याख्येशी समतुल्य अशी पर्यायी व्याख्या ही: वाक्यांचा संच $\Gamma$
\emph{महत्तम सुसंगत} तेव्हा आणि केवळ तेव्हाच असतो, जेव्हा
\begin{enumerate}
\item $\Gamma$ सुसंगत असतो, आणि
\item जर $\Gamma \cup \{ \varphi \}$ सुसंगत असेल, तर $\varphi \in \Gamma$ असते.
\end{enumerate}
म्हणजेच~$\Gamma$ मध्ये आधीपासून नसलेली वाक्ये महत्तम सुसंगत
संच~$\Gamma$ मध्ये जोडली, तर मिळणारा मोठा संच विसंगत होतो.
\end{explain}

\begin{explain}
संपूर्णतेच्या सिद्धतेत महत्तम सुसंगत संच महत्त्वाचे आहेत. कारण प्रत्येक
सुसंगत वाक्यांचा संच~$\Gamma$ हा एखाद्या महत्तम सुसंगत
संचात~$\Gamma^*$ सामावतो, याची खात्री देता येते. शिवाय प्रत्येक
वाक्य~$\varphi$ साठी महत्तम सुसंगत संचात $\varphi$ किंवा त्याचा
निषेध $ \lnot \varphi$ यांपैकी एखादे वाक्य असते. विशेषतः अणुरूप वाक्ये साठी हे खरे
आहे. म्हणून स्थिरांक घालून योग्य रीतीने विस्तारलेल्या भाषेतील
महत्तम सुसंगत संचापासून आपण संरचना रचू शकतो: कोणती अणुरूप
वाक्ये~$\Gamma^*$ मध्ये आहेत यावरून विधेयांचे
अर्थनिर्वचन ठरवतो. मग या संरचनांमध्ये $\Gamma^*$ मधील
सर्व वाक्ये (आणि त्यामुळे $\Gamma$ मधीलसुद्धा) सत्य ठरतात,
हे दाखवता येते. या शेवटच्या विधानाच्या सिद्धतेसाठी
$\lnot \varphi \in \Gamma^*$ तेव्हा आणि केवळ तेव्हाच $\varphi \notin
\Gamma^*$, तसेच $(\varphi \lor \psi) \in \Gamma^*$ तेव्हा आणि केवळ तेव्हाच
$\varphi \in \Gamma^*$ किंवा $\psi \in \Gamma^*$, इत्यादी गुणधर्म लागतात.
\end{explain}

\begin{prop}
\label{fol:com:mcs:prop:mcs}
$\Gamma$ महत्तम सुसंगत आहे असे माना. मग:
\begin{enumerate}
\item \label{fol:com:mcs:prop:mcs-prov-in} जर $\Gamma \Proves \varphi$ असेल, तर $\varphi \in
  \Gamma$.

\item \label{fol:com:mcs:prop:mcs-either-or} कोणत्याही $\varphi$ साठी $\varphi \in
  \Gamma$ किंवा $\lnot \varphi \in \Gamma$.

\item \label{fol:com:mcs:prop:mcs-and} $(\varphi \land \psi) \in \Gamma$
  तेव्हा आणि केवळ तेव्हाच $\varphi \in \Gamma$ आणि $\psi \in \Gamma$ ही दोन्ही.

\item \label{fol:com:mcs:prop:mcs-or} $(\varphi \lor \psi) \in \Gamma$
  तेव्हा आणि केवळ तेव्हाच $\varphi \in \Gamma$ किंवा $\psi \in \Gamma$.

\item \label{fol:com:mcs:prop:mcs-if} $(\varphi \lif \psi) \in \Gamma$
  तेव्हा आणि केवळ तेव्हाच $\varphi \notin \Gamma$ किंवा $\psi \in \Gamma$.
\end{enumerate}
\end{prop}

\begin{proof}
पुढील सर्व भागांत $\Gamma$ महत्तम सुसंगत आहे असे मानू.
\begin{enumerate}
\item जर $\Gamma \Proves \varphi$ असेल, तर $\varphi \in \Gamma$.

$\Gamma \Proves \varphi$ असे माना. उलट $\varphi \notin \Gamma$ असेही
मानल्यास, $\Gamma$ महत्तम सुसंगत असल्यामुळे $\Gamma \cup \{\varphi\}$
विसंगत असतो; म्हणजे $\Gamma \cup \{\varphi\} \Proves \lfalse$.
\ref{fol:seq:prv:prop:provability-contr}
वरून $\Gamma$ विसंगत ठरतो. पण $\Gamma$ सुसंगत आहे या गृहीतकाला
हा व्याघात आहे. म्हणून $\varphi \notin \Gamma$ शक्य नाही; म्हणजे
$\varphi \in \Gamma$.

\item कोणत्याही $\varphi$ साठी $\varphi \in \Gamma$ किंवा $\lnot \varphi \in \Gamma$.

उलट, एखाद्या~$\varphi$ साठी $\varphi \notin \Gamma$ आणि
$\lnot \varphi \notin \Gamma$ असे दोन्ही आहेत असे माना. $\Gamma$ महत्तम
सुसंगत असल्यामुळे $\Gamma \cup \{\varphi\}$ आणि $\Gamma \cup \{\lnot \varphi\}$
दोन्ही विसंगत आहेत. म्हणून $\Gamma \cup \{\varphi\} \Proves \lfalse$
आणि $\Gamma \cup \{\lnot \varphi\} \Proves \lfalse$.
\ref{fol:seq:prv:prop:provability-exhaustive}
वरून $\Gamma$ विसंगत ठरतो, हा व्याघात. म्हणून असे वाक्य~$\varphi$
असू शकत नाही. प्रत्येक $\varphi$ साठी $\varphi \in \Gamma$ किंवा
$\lnot \varphi \in \Gamma$.

\item

$(\varphi \land \psi) \in \Gamma$ तेव्हा आणि केवळ तेव्हाच $\varphi \in \Gamma$
आणि $\psi \in \Gamma$ ही दोन्ही:

पुढील दिशेसाठी $(\varphi \land \psi) \in \Gamma$ असे माना. मग
$\Gamma \Proves \varphi \land \psi$. आता
\ref{fol:seq:ppr:prop:provability-land-left}
वरून $\Gamma \Proves \varphi$ आणि $\Gamma \Proves \psi$. तसेच
\ref{fol:com:mcs:prop:mcs-prov-in} वरून अपेक्षेप्रमाणे $\varphi \in \Gamma$
आणि $\psi \in \Gamma$.

उलट दिशेसाठी $\varphi \in \Gamma$ आणि $\psi \in \Gamma$ असे माना.
मग $\Gamma \Proves \varphi$ आणि $\Gamma \Proves \psi$. आता
\ref{fol:seq:ppr:prop:provability-land-right}
वरून $\Gamma \Proves \varphi \land \psi$. \ref{fol:com:mcs:prop:mcs-prov-in}
वरून $(\varphi \land \psi) \in \Gamma$.

\item

$(\varphi \lor \psi) \in \Gamma$ तेव्हा आणि केवळ तेव्हाच $\varphi \in \Gamma$
किंवा $\psi \in \Gamma$.

पुढील दिशेचा प्रतिव्यत्यास सिद्ध करण्यासाठी $\varphi \notin \Gamma$
आणि $\psi \notin \Gamma$ असे माना. आपल्याला $(\varphi \lor \psi)
\notin \Gamma$ दाखवायचे आहे. $\Gamma$ महत्तम सुसंगत असल्यामुळे
$\Gamma \cup \{\varphi\} \Proves \lfalse$ आणि
$\Gamma \cup \{\psi\} \Proves \lfalse$. आता
\ref{fol:axd:prv-alt:prop:provability-lor-left}
वरून $\Gamma \cup \{(\varphi \lor \psi)\}$ विसंगत ठरतो. म्हणून
अपेक्षेप्रमाणे $(\varphi \lor \psi) \notin \Gamma$.

उलट दिशेसाठी $\varphi \in \Gamma$ किंवा $\psi \in \Gamma$ असे
माना. मग $\Gamma \Proves \varphi$ किंवा $\Gamma \Proves \psi$. आता
\ref{fol:axd:prv-alt:prop:provability-lor-right}
वरून $\Gamma \Proves \varphi \lor \psi$. \ref{fol:com:mcs:prop:mcs-prov-in}
वरून अपेक्षेप्रमाणे $(\varphi \lor \psi) \in \Gamma$.

\item

$(\varphi \lif \psi) \in \Gamma$ तेव्हा आणि केवळ तेव्हाच
$\varphi \notin \Gamma$ किंवा $\psi \in \Gamma$:

पुढील दिशेसाठी $(\varphi \lif \psi) \in \Gamma$ असे माना.
उलट $\varphi \in \Gamma$ आणि $\psi \notin \Gamma$ असेही माना.
या गृहीतकांवरून $\Gamma \Proves \varphi \lif \psi$ आणि
$\Gamma \Proves \varphi$. आता
\ref{fol:axd:prv-alt:prop:provability-mp}
वरून $\Gamma \Proves \psi$. परंतु \ref{fol:com:mcs:prop:mcs-prov-in} वरून
$\psi \in \Gamma$, हा $\psi \notin \Gamma$ या गृहीतकाला व्याघात.

उलट दिशेसाठी प्रथम $\varphi \notin \Gamma$ हे प्रकरण पाहू.
\ref{fol:com:mcs:prop:mcs-either-or} वरून $\lnot \varphi \in \Gamma$ आणि म्हणून
$\Gamma \Proves \lnot \varphi$. आता
\ref{fol:seq:ppr:prop:provability-lif}
वरून $\Gamma \Proves \varphi \lif \psi$.
पुन्हा \ref{fol:com:mcs:prop:mcs-prov-in} वरून अपेक्षेप्रमाणे
$(\varphi \lif \psi) \in \Gamma$.

आता $\psi \in \Gamma$ हे प्रकरण पाहू. मग $\Gamma \Proves \psi$,
आणि
\ref{fol:seq:ppr:prop:provability-lif}
वरून $\Gamma \Proves \varphi \lif \psi$. \ref{fol:com:mcs:prop:mcs-prov-in}
वरून $(\varphi \lif \psi) \in \Gamma$.
\end{enumerate}
\end{proof}

\begin{prob}
\ref{fol:com:mcs:prop:mcs} ची सिद्धता पूर्ण करा.
\end{prob}












\section{परिचय}\label{fol:syn:int:sec}

प्रथम-क्रम तर्कशास्त्राची उपपत्ती आणि अधिसिद्धांत विकसित करण्यासाठी,
आपण प्रथम त्याच्या व्यंजकांची विन्यासमीमांसा आणि चिन्हार्थमीमांसा
परिभाषित केली पाहिजे. प्रथम-क्रम तर्कशास्त्राची व्यंजके म्हणजे पदे आणि
सूत्रे. पदे चले, स्थिरांक आणि फलने
यांपासून तयार होतात. सूत्रे मात्र विधेये आणि पदे
यांपासून तयार होतात (यांतून सर्वांत लहान, ``आण्विक'' सूत्रे
तयार होतात); मग तार्किक संयोजके आणि संख्यापक वापरून आण्विक
सूत्रांपासून अधिक जटिल सूत्रे तयार करता येतात. रचनानियम
मांडण्याचे अनेक निरनिराळे मार्ग आहेत; आपण त्यांपैकी केवळ एक शक्य
मार्ग देतो. इतर पद्धती वेगळी चिन्हे निवडतील, संयोजकांचे वेगळे संच
आदिम म्हणून निवडतील आणि कंसांचा वेगळ्या रीतीने वापर करतील (किंवा
तथाकथित पोलिश संकेतपद्धतीप्रमाणे कंस अजिबात वापरणार नाहीत). तरी सर्व
दृष्टिकोनांत सामाईक बाब ही की रचनानियम पदांचा आणि सूत्रांचा
संच \emph{विगमनाने} परिभाषित करतात. हे योग्य रीतीने केल्यास,
रचनानियमांनुसार प्रत्येक व्यंजक मूलतः एकाच रीतीने तयार होऊ शकते.
त्यामुळे तयार होणारी व्यंजके \emph{एकमेव रीतीने वाचनीय} असतात. परिणामी,
त्यांची ही विगमनाधारित व्याख्या आपल्याला त्याच पद्धतीने---विगमनाधारित
व्याख्येने---त्या व्यंजकांना अर्थ देता येतो, याची खात्री देते.

व्यंजकांना अर्थ देणे हे चिन्हार्थमीमांसेचे क्षेत्र आहे. चिन्हार्थमीमांसेतील
मध्यवर्ती संकल्पना म्हणजे संरचना मधील पूर्ति. संरचना
भाषेच्या मूलभूत घटकांना अर्थ देते: प्रांत म्हणजे वस्तूंचा अरिक्त
संच. संख्यापक या प्रांतातील घटकांवर फिरतात असा त्यांचा अर्थ लावला जातो,
स्थिरांक ना प्रांतातील घटक नेमून दिले जातात, फलने ना
प्रांत पासून त्याच प्रांताकडे जाणारी फलने नेमून दिली जातात आणि
विधेये ना प्रांत वरील संबंध नेमून दिले जातात. प्रांत
आणि मूलभूत शब्दसंग्रहातील चिन्हांना केलेल्या नेमणुका मिळून
संरचना बनते.
चले ही सूत्रे मध्ये येऊ शकतात; त्यांना चिन्हार्थ देण्यासाठी
आपल्याला प्रांत मधील घटक ही नेमून द्यावी लागतात---हेच
चर-मूल्यांकन होय. शेवटी, पूर्ति-संबंध या सर्व गोष्टी एकत्र आणतो.
सूत्राची संरचना~$\Struct{M}$ मध्ये चर-मूल्यांकन~$s$ च्या सापेक्ष पूर्ति होऊ शकते; हे $\Sat{M}{\varphi}[s]$ असे
लिहितात. हा संबंधदेखील~$\varphi$ च्या संरचनेवरील विगमनाने परिभाषित केला जातो.
उदाहरणार्थ, $\varphi \land \psi$ ची पूर्ति परिभाषित करताना तार्किक संयोजकांचे
सत्यता-कोष्टक वापरून ती $\varphi$ आणि~$\psi$ यांची पूर्ति होते (किंवा होत नाही)
यांच्या आधारे सांगितली जाते. त्यानंतर असे निष्पन्न होते की सूत्र~$\varphi$
हे वाक्य, म्हणजे त्यात मुक्त चरे नसतील, तर चराचे
मूल्यांकन असंबद्ध ठरते; त्यामुळे संरचनांमध्ये वाक्यांची
सरळ पूर्ति होते (किंवा होत नाही) असे आपण म्हणू शकतो.

वाक्यांसाठीच्या पूर्ति-संबंध~$\Sat{M}{\varphi}$ च्या आधारे आपण वैधता,
तार्किक निष्पन्नता आणि पूर्ततायोग्यता या मूलभूत चिन्हार्थविषयक संकल्पना
परिभाषित करू शकतो. एखाद्या वाक्याची प्रत्येक संरचनेत पूर्ति
होत असेल, तर ते वैध, $\Entails \varphi$, असते. वाक्यांच्या एखाद्या
संचातून ते तार्किकरीत्या निष्पन्न होते, $\Gamma \Entails \varphi$, जर~$\Gamma$ मधील
सर्व वाक्यांची पूर्ति करणारे प्रत्येक संरचना हे~$\varphi$ चीही
पूर्ति करत असेल. आणि वाक्यांच्या एखाद्या संचातील सर्व
वाक्यांची एकाच वेळी पूर्ति करणारे एखादे संरचना असेल, तर तो
संच पूर्ततायोग्य असतो. सूत्रे विगमनाधारित रीतीने परिभाषित केली आहेत
आणि पूर्ति हीही सूत्रांच्या संरचनेवरील विगमनाने परिभाषित केली जाते;
म्हणून आपल्या चिन्हार्थमीमांसेचे गुणधर्म सिद्ध करण्यासाठी आणि परिभाषित
चिन्हार्थविषयक संकल्पनांचा परस्परसंबंध लावण्यासाठी विगमन वापरता येते.












\section{$C$ मधील फलने~$\Th{Q}$ मध्ये निरूपणीय असतात}\label{inc:req:cre:sec}

$C$ मधील प्रत्येक फलन~$\Th{Q}$ मध्ये निरूपणीय आहे, हे दाखवायचे
आहे. अखेरीस $C$ मधील प्रत्येक $k$-आदानी फलन
$f(x_0,\dots,x_{k-1})$ ला त्याचे निरूपण करणारे सूत्र
$\varphi_f(x_0,\dots,x_{k-1},y)$ कसे द्यायचे, हे दाखवावे लागेल.

\begin{lem}
$n$ आणि $m$ या नैसर्गिक संख्या असू द्या. $n \neq m$ असेल,
तर $Q \Proves \num n \neq \num m$.
\end{lem}

\begin{proof}
$n$ वर विगमन वापरून प्रत्येक $m$ साठी $n \neq m$ असेल, तर
$Q \Proves \eq/[\num n][\num m]$, हे दाखवू.

आधार टप्प्यात $n = 0$. जर $m$ हा $0$ च्या बरोबर नसेल,
तर एखाद्या नैसर्गिक संख्या $k$ साठी $m = k +
1$. आपल्याकडे $\lforall[x][\eq/[0][x']]$ असे सांगणारे स्वयंसिद्धक
आहे. संख्यापकाच्या स्वयंसिद्धकाने $x$ च्या जागी
$\num k$ ठेवून $\eq/[0][\num k']$ हा निष्कर्ष मिळतो. पण
$\num k'$ म्हणजेच $\num m$.

विगमन टप्प्यात विधान $n$ साठी खरे आहे असे मानून $n+1$
पाहू. $m$ ही कोणतीही नैसर्गिक संख्या असू द्या. दोन शक्यता आहेत:
$m = 0$, किंवा एखाद्या $k$ साठी $m = k+1$. पहिली बाब
वरप्रमाणे हाताळता येते. दुसऱ्या बाबतीत $n+1 \neq
k+1$ असे समजा. मग $n \neq k$. $n$ साठीच्या विगमन गृहीतकाने
$\Th{Q} \Proves \eq/[\num n][\num k]$. आपल्याकडे
$\lforall[x][\lforall[y][\eq[x'][y'] \lif \eq[x][y]]]$
असे सांगणारे स्वयंसिद्धक आहे. संख्यापकाच्या स्वयंसिद्धकाने
$\eq[\num n'][\num k'] \lif \eq[\num n][\num
  k]$ मिळते. विधानीय तर्कशास्त्र वापरून $\Th{Q}$ मध्ये
$\eq/[\num n][\num k] \lif \eq/[\num n'][\num k']$
हे निष्पन्न होते. विध्यात्मक अनुमानाने
$\eq/[\num n'][\num k']$ मिळते; आणि $\num k'$ म्हणजे
$\num m$ असल्यामुळे हेच अपेक्षित होते.
\end{proof}

या पूर्वप्रमेयाचा आशय मर्यादित आहे: वेगवेगळे संख्यांक
वेगवेगळ्या वस्तू दर्शवतात, हे $\Th{Q}$ सिद्ध करू शकते. उदाहरणार्थ,
$\Th{Q}$ हे $0'' \neq 0'''$ सिद्ध करते. पण हे सर्वसाधारणपणे
दाखवण्यासाठी थोडी काळजी घ्यावी लागते. येथे वापरलेले विगमन
$\Th{Q}$ च्या \emph{बाहेर} आहे, हेही लक्षात घ्या.

शून्य, उत्तरवर्ती, बेरीज, गुणाकार, समतेचे जातिबोधक फल
आणि प्रक्षेपण यांचे निरूपण करता येईल. प्रत्येक बाबतीत योग्य
निरूपक सूत्र सहज दिसते. उदाहरणार्थ, शून्याचे निरूपण
\[
y = 0,
\]
उत्तरवर्तीचे निरूपण
\[
x_0' = y,
\]
आणि बेरीजेचे निरूपण
\[
x_0 + x_1 = y.
\]
खरे काम म्हणजे संबंधित विधाने $\Th{Q}$ सिद्ध करू शकते, हे
दाखवणे. उदाहरणार्थ, वरील सूत्र बेरीजेचे निरूपण करते असे
म्हणण्यासाठी नैसर्गिक संख्यांच्या प्रत्येक जोडी $m$ आणि $n$
साठी $\Th{Q}$ पुढील सिद्ध करते, हे दाखवावे लागते:
\[
\eq[\num n + \num m][\num {n+m}]
\]
आणि
\[
\lforall[y][(\eq[(\num n + \num m)][y] \lif \eq[y][\num {n+m}])].
\]

संयोजनाचे काय? $h$ पुढीलप्रमाणे परिभाषित आहे असे माना:
\[
h(x_0,\dots,x_{l-1}) = f(g_0(x_0,\dots,x_{l-1}), \dots,
g_{k-1}(x_0,\dots,x_{l-1})).
\]
येथे $f,g_0,\dots,g_{k-1}$ यांचे अनुक्रमे निरूपण करणारी
सूत्रे $\varphi_f, \varphi_{g_0}, \dots,
\varphi_{g_{k-1}}$ आपल्याला आधीच मिळाली आहेत. आता $h$ चे निरूपण
करणारे सूत्र $\varphi_h$ परिभाषित करता येते. त्यासाठी
$\varphi_h(x_0,\dots,x_{l-1},y)$ पुढील असे घेऊ:
\begin{multline*}
  \lexists[z_0\dots][\lexists[z_{k-1}][(\varphi_{g_0}(x_0,\dots,x_{l-1},z_0) \land
  \dots \land \varphi_{g_{k-1}}(x_0,\dots,x_{l-1},z_{k-1}) \land\\
  \varphi_f(z_0,\dots,z_{k-1},y))]].
\end{multline*}


शेवटी अपरिबद्ध शोध पाहू. $g(x,\vec z)$ नियमित आहे आणि
$\Th{Q}$ मध्ये, समजा $\varphi_g(x,\vec
z,y)$ या सूत्राने, निरूपणीय आहे असे माना. $f$ ची व्याख्या
$f(\vec z) = \mu x \; g(x,\vec z)$ अशी असू द्या. $f$ चे
निरूपण करणारे सूत्र $\varphi_f(\vec z,y)$ शोधायचे आहे. नैसर्गिक
निवड अशी:
\[
\varphi_f(\vec z,y) \ident \varphi_g(y,\vec z,0) \land \lforall[w][(w < y
\lif \lnot \varphi_g(w,\vec z,0))].
\]

आणखी काही पूर्वप्रमेये वापरून हे सूत्र योग्य आहे असे दाखवता येते; उदाहरणार्थ:

\begin{lem}
  प्रत्येक चल $x$ आणि प्रत्येक नैसर्गिक संख्या $n$ साठी $\Th{Q}$
  हे $\eq[(x'
  + \num n)][(x + \num n)']$ सिद्ध करते.
\end{lem}

पुन्हा लक्षात घ्या: यावरून $\Th{Q}$ हे
$\lforall[x][\lforall[y][(x' + y) = (x +y)']]$
सिद्ध करते असा सिद्धतायोग्यतेचा दावा चुकीचा आहे.

\begin{proof}
नेहमीप्रमाणे $n$ वर विगमन वापरू. आधार टप्प्यात $n =
0$ असेल, तर $\Th{Q}$ हे $\eq[(x'+0)][(x+0)']$
सिद्ध करते, हे दाखवायचे आहे. पुढील विधाने मिळतात:
\begin{eqnarray*}
& & \eq[(x' + 0)][x'] \quad \text{स्वयंसिद्धक 4 वरून} \\
& & \eq[(x + 0)][x] \quad \text{स्वयंसिद्धक 4 वरून} \\
& & \eq[(x+0)'][x'] \quad \text{ओळ 2 वरून} \\
& & \eq[(x' + 0)][(x + 0)'] \quad \text{ओळी 1 आणि 3 वरून}
\end{eqnarray*}
विगमन टप्प्यात $\Th{Q}$ मध्ये $\eq[(x' +
  \num n)][(x + \num n)']$ आधी सिद्ध झाले आहे असे
माना. $\num{n+1}$ म्हणजे $\num
n'$ असल्यामुळे $\Th{Q}$ हे
$\eq[(x' + \num n')][(x + \num
n')']$ सिद्ध करते, हे दाखवायचे आहे. पुढील समताशृंखला
$\Th{Q}$ मध्ये सिद्ध करता येते:
\begin{eqnarray*}
(x' + \num n') & \eq & (x' + \num n)' \quad \text{स्वयंसिद्धक 5} \\
& \eq & (x + \num n')' \quad \text{विगमन गृहीतकावरून}
\end{eqnarray*}
\end{proof}

\begin{lem}
\begin{enumerate}
\item $\Th{Q}$ हे $\lnot (x < \num 0)$ सिद्ध करते.
\item प्रत्येक नैसर्गिक संख्या $n$ साठी $\Th{Q}$ पुढील सिद्ध करते:
\[
x < \num {n+1} \lif (\eq[x][\num 0] \lor \dots \lor \eq[x][\num n]).
\]
\end{enumerate}
\end{lem}

\begin{proof}
१ आणि २ चा काही भाग अनौपचारिकपणे पाहू; म्हणजे औपचारिक
निष्पत्ती कसे रचावे यासाठी फक्त संकेत देऊ.

भाग १ साठी $<$ च्या व्याख्येनुसार $\Th{Q}$ मध्ये
$\lnot \lexists[y][\eq[(y'+x)][0]]$ सिद्ध करायचे आहे. प्रथम-क्रम
तर्कशास्त्राच्या स्वयंसिद्धकांनी आणि नियमांनी हे
$\lforall[y][\eq/[(y' +
    x)][0]]$ याच्याशी समतुल्य आहे. कल्पना अशी: $\eq[(y'+x)][0]$
असे माना. $x$ हा 0 असेल, तर $\eq[(y' + 0)][0]$.
पण $\Th{Q}$ च्या स्वयंसिद्धक 4 वरून
$\eq[(y' + 0)][y']$, आणि स्वयंसिद्धक 2 वरून
$\eq/[y'][0]$; हा व्याघात. म्हणून
$\lforall[y][\eq/[(y' +x)][0]]$. $x$ हा $0$ नसेल,
तर स्वयंसिद्धक 3 वरून असा $z$ असतो की
$\eq[x][z']$. मग $\eq[(y' + z')][0]$.
स्वयंसिद्धक~5 वरून $\eq[(y' + z)'][0]$, हेही
स्वयंसिद्धक~2 शी विसंगत आहे.

भाग २ साठी $n$ वर विगमन वापरा. प्रथम $n =0$
हा आधार टप्पा पाहू. त्या वेळी
$x < \num 1 \lif \eq[x][\num
  0]$ दाखवायचे आहे. $x < \num 1$ असे माना.
$<$ च्या परिभाषक स्वयंसिद्धकाने
$\lexists[y][\eq[(y'+x)][0']]$ मिळते. $y$ कडे हा
गुणधर्म आहे असे माना; म्हणजे $\eq[y'+x][0']$.

$\eq[x][0]$ दाखवायचे आहे. स्वयंसिद्धक 3 वरून
$x$ हा 0 नसेल, तर एखाद्या $z$ साठी तो $z'$ च्या
बरोबर असतो. मग $\eq[(y' + z')][0']$.
$\Th{Q}$ च्या स्वयंसिद्धक~5 वरून
$\eq[(y'+z)'][0']$. स्वयंसिद्धक 1 वरून
$\eq[(y' +
    z)][0]$. पण व्याख्येने याचा अर्थ
$z < 0$; हे भाग~१ शी विसंगत आहे.
\end{proof}

\begin{explain}
संगणनक्षम फलनांचा संच म्हणजेच $\Th{Q}$ मध्ये निरूपणीय
फलनांचा संच, असे लक्षणन आपण दाखवले आहे. प्रत्यक्षात हा
पुरावा अधिक व्यापक आहे. निरूपणीयतेच्या व्याख्येवरून
$\Th{Q}$ चा विस्तार असलेल्या (किंवा ज्यात $\Th{Q}$ चे
अर्थनिर्धारण करता येते अशा) कोणत्याही उपपत्तीत
संगणनक्षम फलनांचे निरूपण करता येते, हे सहज दिसते.
उलट दिशेसाठी पूर्वीच्या निरूपणीयतेच्या विभागातील अटी
आवश्यक आहेत: निष्पत्ती व्यवस्थेतील उपपत्ती सुसंगत
असावी, भिन्न संख्यांकांची असमता सिद्ध व्हावी आणि
सिद्धतांची तपासणी संगणनक्षम असावी. या अटींखाली
प्रत्येक निरूपणीय फलन संगणनक्षम असते. उदाहरणार्थ, संगणनक्षम फलनांचा संच
म्हणजे पेआनो अंकगणितात किंवा अगदी झर्मेलो--फ्रेंकेल संच
उपपत्तीत निरूपणीय फलनांचा संच, असे त्या उपपत्ती सुसंगत
आहेत या गृहीतकाखाली लक्षणन करता येते.
गोडेल यांनी नोंदवल्याप्रमाणे हे काहीसे आश्चर्यकारक आहे.
सिद्धतायोग्यतेचे प्रश्न कोणती उपपत्ती विचारात घेतो यावर
फार अवलंबून असतात, हे आपण पाहू: साधारणपणे स्वयंसिद्धके
जितकी बलवान, तितके अधिक सिद्ध करता येते. या अटी पूर्ण करणाऱ्या आणि संगणनक्षम फलनांचे निरूपण
करणाऱ्या स्वयंसिद्धकीय उपपत्त्यांत निरूपणीय फलने
नेमकी संगणनक्षम फलनेच असतात.

\end{explain}



\begin{quote}\small\textbf{स्रोतनोंद.} SOL6-A289: संयोजनातील सर्व मध्यवर्ती मूल्ये आणि अंतिम फलसूत्र यांना अस्तित्व-संख्यापकांची एकच व्याप्ती दिली आहे. निरूपणीय फलनांच्या उलट दाव्यास सुसंगतता, वेगळ्या संख्यांकांची सिद्ध असमता आणि संगणनक्षम सिद्धता-तपासणीच्या अटी आवश्यक आहेत. मूळ एक-दिशीय आधार टप्प्यात उलट दिशा गहाळ असल्याची पूर्वीची नोंद चुकीची होती; ती मागे घेतली आहे.\end{quote}
\begin{quote}\small\textbf{स्रोतनोंद.} OLINC-328: मूळ संयोजन-सूत्रातील संख्यापक आणि संयुक्त-विधानाचे कंस SOL6-A289 मध्ये दुरुस्त केले आहेत. मध्यवर्ती z चलांच्या अस्तित्वाची व्याप्ती सर्व g सूत्रे आणि अंतिम f सूत्र व्यापते.\end{quote}
\begin{quote}\small\textbf{स्रोतनोंद.} OLINC-327: SOL6-A289 मध्ये पूर्वीच्या OLINC-035 प्रमाणे सुसंगतता, भिन्न संख्यांकांची सिद्ध असमता आणि सिद्धतांची संगणनक्षम तपासणी या अटी उलट दाव्यात स्पष्ट घातल्या आहेत.\end{quote}









\section{फलनांचा संच~$C$}\label{inc:req:cee:sec}

$C$ हा पुढील फलने समाविष्ट करणारा सर्वात लहान फलनसंच असू द्या:
\begin{enumerate}
\item $0$,
\item उत्तरवर्ती फलन,
\item बेरीज,
\item गुणाकार,
\item प्रक्षेपण फलने, आणि
\item समतेचे जातिबोधक फल, $\Char{=}$;
\end{enumerate}
आणि जो पुढील क्रियांखाली बंद आहे:
\begin{enumerate}
\item संयोजन, आणि
\item नियमित फलनांवर लावलेला अपरिबद्ध शोध.
\end{enumerate}
ही शेवटची मर्यादा म्हणजे $\mu$ क्रिया फक्त तिचा परिणाम
सर्वत्र परिभाषित असेल तेव्हाच वापरता येते, एवढेच. याची
\emph{सामान्य पुनरावर्ती} फलनांच्या व्याख्येशी तुलना करा:
येथे बेरीज, गुणाकार आणि $\Char{=}$ समाविष्ट केले आहेत,
पण आदिम पुनरावर्तन वगळले आहे.

$C$ मधील प्रत्येक फलन पुनरावर्ती आहे, हे स्पष्ट आहे; कारण
बेरीज, गुणाकार आणि $\Char{=}$ ही फलने पुनरावर्ती आहेत.
याचा उलट दावा खरा आहे, हेही दाखवू. म्हणजे $C$ मधील
इतर क्रिया वापरून आदिम पुनरावर्तन करता येते.












\section{विधानसंच}\label{int:sem:pro:sec}

कधीकधी संबंधात्मक प्रतिमान~$\mModel{M}$ मधील ज्या
जगतांत सूत्र~$\varphi$ सत्य असते, त्या जगांच्या संचाला नाव देणे
उपयुक्त ठरते. त्याला $\varphi$ ने परिभाषित केलेला
\emph{विधानसंच} म्हणू आणि $\Prop{M}{\varphi}$ असे लिहू.

\begin{defn}
  $V$ या फलनाचा विस्तार $\Prop{M}{\varphi} \colon \Frm
  \to \Pow{W}$ अशा फलनात पुढीलप्रमाणे विगमनाने करू:
  \begin{enumerate}
  \item $\Prop{M}{\lfalse} = \emptyset$
  \item $\Prop{M}{p} = V(p)$
  \item $\Prop{M}{\lnot \psi} = {}$
    \[
      \Setabs{w \in W}{\text{$Rww'$ असेल अशा प्रत्येक $w'$ साठी $w'
          \notin \Prop{M}{\psi}$}}.
    \]
  \item $\Prop{M}{\psi \land \chi} = \Prop{M}{\psi} \cap \Prop{M}{\chi}$.
  \item $\Prop{M}{\psi \lor \chi} = \Prop{M}{\psi} \cup \Prop{M}{\chi}$.
  \item $\Prop{M}{\psi \lif \chi} = {}$
    \[
      \Setabs{w \in W}{\text{$Rww'$ असेल अशा प्रत्येक $w'$ साठी,
          $w' \in \Prop{M}{\psi}$ असेल तर $w' \in \Prop{M}{\chi}$}}.
    \]
  \end{enumerate}
\end{defn}


\begin{prop}
  \label{int:sem:pro:prop:propositions-true}
  $\mSat{M}{\varphi}[w]$ तेव्हा आणि केवळ तेव्हाच $w \in \Prop{M}{\varphi}$.
\end{prop}


\begin{proof}
  सराव.
\end{proof}

\begin{prob}
  \ref{int:sem:pro:prop:propositions-true} सिद्ध करा.
\end{prob}



\begin{quote}\small\textbf{स्रोतनोंद.} OLINC-331: या चिन्हात A हे बदलते सूत्र आहे; त्याला त्या सूत्राने परिभाषित केलेल्या जगांचा संच जोडणारे फलन असा आशय आहे. ही चिन्हवाचनाची स्पष्टता आहे, सिद्धतेतील उणीव किंवा निश्चित नवा स्रोतदोष नव्हे.\end{quote}
\begin{quote}\small\textbf{स्रोतनोंद.} OLINC-330: स्रोताने पूर्ति-संबंधाच्या जागी फक्त प्रतिमानाचे नाव छापणारा mModel मॅक्रो वापरला आहे; येथे mSat केला आहे.\end{quote}










\section{याद्या}\label{lam:ldf:lst:sec}

यादी $[M_1, M_2, \ldots, M_n]$ पुढीलप्रमाणे परिभाषित करता येते.
येथे $s$ आणि $z$ ही परस्पर भिन्न आणि कोणत्याही $M_i$ मध्ये
मुक्त नसलेली चरे ही पदे गुंफण्यापूर्वी निवडली आहेत:
\[
  [M_1, M_2, \ldots, M_n] \ident \lambd[sz][s M_1 (s M_2 (\ldots (s M_n z)))]
\]
अनौपचारिकपणे, यादीचे निरूपण करणारे पद हे त्या यादीचे
“संचायक” फलन असते. हे फलन दोन पदे स्वीकारते: पहिले
असे फलन असते, जे आतापर्यंत साचलेले मूल्य आणि नवा घटक
घेऊन नवे संचयित मूल्य देते; दुसरे म्हणजे संचयाचे आरंभीचे
मूल्य. ते एकेका घटकाचा संचय करते आणि शेवटी परिणाम देते.

\begin{digress}
  फलनीय प्रोग्रामलेखनाशी परिचय असेल, तर ही व्याख्या
  \emph{उजवी घडी} (right fold) यासारखी वाटेल:
  यादीतील घटकांवर उजवी घडी घालणारे फलन म्हणूनच
  यादी परिभाषित होते.
\end{digress}

या संकेतावर आधारलेली काही उपयुक्त फलने सहज परिभाषित
करता येतात:
\begin{align*}
  \fn{Sum} & \ident
  \lambd[l][l \, (\lambd[xa][\fn{Add} \, x \, a])\,  \num{0}]\\
  \fn{Len} & \ident
  \lambd[l][l \, (\lambd[xa][\fn{Add} \, a \, \num{1}]) \num{0}]
\end{align*}

$\fn{Sum}$ हे चर्च अंकांच्या यादीतील बेरजेची गणना
करते. ते यादीवर संचय करते: आरंभीचे मूल्य
$\num{0}$ असते आणि प्रत्येक घटकासाठी
$\fn{Add} \, x\, a$ हे नवे मूल्य देते
(येथे $x$ हा चालू घटक आणि $a$ हे संचयित मूल्य आहे).
परिणाम म्हणजे सर्व घटकांची बेरीज.

\begin{prob}
  $\fn{Len}$ काय करते? स्पष्ट करा.
\end{prob}

\begin{prob}
  यादीचा क्रम उलटवणारे फलन परिभाषित करा.
\end{prob}


\begin{quote}\small\textbf{स्रोतनोंद.} SOL6-A292: यादीतील पदांच्या मुक्त चरांना बाहेरच्या lambda ने पकडू नये म्हणून संचायक आणि आरंभीच्या मूल्याची चरे यादी गुंफण्यापूर्वी निवडली आहेत. ही अट सर्वसाधारण पदांसाठी आवश्यक आहे; सर्व बंद पदांसाठी ती आपोआप पूर्ण होते.\end{quote}









\section{परिवर्तन आणि न्यूनीकरण}\label{lam:int:con:sec}

लॅम्डा कलनात परिवर्तनाचे किंवा न्यूनीकरणाचे तीन प्रकार आहेत.
“रूपांतरण” एकदिश आहे की द्विदिश, यावर परिवर्तन आणि
न्यूनीकरण यांतील फरक अवलंबून असतो; हे आपण पुढे पाहू.

$\alpha$-परिवर्तनात चरांची नावे बदलली जातात.
$\lambda x. x$ आणि $\lambda y.y$ ही दोन्ही पदे एकच फलन
(ओळख फलन) दर्शवतात असे मानणे स्वाभाविक आहे.
ही कल्पना आपण पुढीलप्रमाणे औपचारिक करू:




\begin{quote}\small\textbf{स्रोतनोंद.} OLINC-332: गोठवलेला स्रोत येथेच थांबतो; आश्वासित औपचारिक व्याख्या किंवा उर्वरित प्रकार या विभागात दिलेले नाहीत.\end{quote}










\section{सिद्धता-उपपत्तीय संकल्पना}\label{mvl:seq:prf:sec}


\begin{explain}
जशा आपण अनेक महत्त्वाच्या चिन्हार्थविषयक संकल्पना (वैधता, तार्किक निष्पन्नता,
पूर्ततायोग्यता) परिभाषित केल्या आहेत, तशाच त्यांना अनुरूप
\emph{सिद्धता-उपपत्तीय संकल्पना} आता परिभाषित करू. संरचनांमध्ये
वाक्यांची पूर्ति होते की नाही, याचा आधार घेऊन या संकल्पना परिभाषित
केल्या जात नाहीत; त्याऐवजी ठरावीक क्रमवर्तींची निष्पन्नता किंवा
अनिष्पन्नता यांचा आधार घेतला जातो. या दोन प्रकारच्या संकल्पना
एकमेकांशी जुळतात, हा एक महत्त्वाचा शोध होता. त्या जुळतात, हाच
\emph{निर्दोषता प्रमेय} आणि \emph{संपूर्णता प्रमेय} यांचा आशय आहे.
\end{explain}

\begin{defn}[प्रमेये]
वाक्य~$\varphi$ हे \emph{प्रमेय} असते, जर $\Log{LK}$ मध्ये
$\quad \Sequent \varphi$ या क्रमवर्तीची निष्पत्ती असेल. $\varphi$ हे प्रमेय
असेल तर आपण $\Proves \varphi$ लिहितो आणि ते प्रमेय नसेल तर $\Proves/ \varphi$ लिहितो.
\end{defn}

\begin{defn}[निष्पन्नता]
वाक्य~$\varphi$ हे वाक्यांच्या संच~$\Gamma$ \emph{पासून
निष्पन्न करता येण्याजोगे} असते, म्हणजे $\Gamma \Proves \varphi$, तेव्हा आणि केवळ तेव्हाच,
जेव्हा एखादा सांत उपसंच~$\Gamma_0 \subseteq \Gamma$ आणि $\Gamma_0$ मधील
वाक्यांची एखादी क्रमिका $\Gamma_0'$ अशी असते की $\Log{LK}$
$\Gamma_0' \Sequent \varphi$ हे निष्पन्न करते. $\varphi$ हे $\Gamma$ पासून
निष्पन्न करता येण्याजोगे नसेल, तर आपण $\Gamma \Proves/ \varphi$ लिहितो.
\end{defn}

आकुंचन, दुर्बलीकरण आणि अदलाबदल नियमांमुळे $\Gamma_0'$ मधील
वाक्यांचा क्रम आणि त्यांची पुनरावृत्ती कितीदा होते हे महत्त्वाचे नसते:
$\Gamma_0' \Sequent \varphi$ ही क्रमवर्ती निष्पन्न करता येण्याजोगे असेल, तर
$\Gamma_0'$ मध्ये असलेली तीच वाक्ये असणाऱ्या कोणत्याही
$\Gamma_0''$ साठी $\Gamma_0'' \Sequent \varphi$ ही क्रमवर्तीसुद्धा निष्पन्न करता
येते. उदाहरणार्थ, $\Gamma_0 = \{\psi, \chi\}$ असेल, तर
$\Gamma_0' = \tuple{\psi, \psi, \chi}$ आणि $\Gamma_0'' =
\tuple{\chi, \chi, \psi}$ या दोन्ही क्रमिकांमध्ये $\Gamma_0$ मधीलच
वाक्ये आहेत. यांपैकी एक क्रमिका असणारी क्रमवर्ती निष्पन्न करता येण्याजोगे असेल,
तर दुसरी असणारी क्रमवर्तीसुद्धा निष्पन्न करता येते; उदा.:
\begin{prooftree}
  \AxiomC{}
  \Deduce$\psi, \psi, \chi \fCenter \varphi$
  \RightLabel{\LeftR{\Contraction}}
  \UnaryInf$\psi, \chi \fCenter \varphi$
  \RightLabel{\LeftR{\Exchange}}
  \UnaryInf$\chi, \psi \fCenter \varphi$
  \RightLabel{\LeftR{\Weakening}}
  \UnaryInf$\chi, \chi, \psi \fCenter \varphi$
\end{prooftree}
यापुढे $\Gamma_0$ हा वाक्यांचा सांत संच असेल, तर
$\Gamma_0 \Sequent \varphi$ याचा अर्थ असा कोणताही क्रमवर्ती घेऊ की ज्याच्या
पूर्वांगात $\Gamma_0$ मधील वाक्यांची क्रमिका आहे; तसेच आवश्यक असल्यास
आकुंचन, अदलाबदल आणि दुर्बलीकरणे गृहीत धरू.

\begin{defn}[सुसंगतता]
वाक्यांचा संच~$\Gamma$ हा \emph{विसंगत} असतो तेव्हा आणि केवळ तेव्हाच, जेव्हा
एखादा सांत उपसंच~$\Gamma_0 \subseteq \Gamma$ असा असतो की $\Log{LK}$
$\Gamma_0 \Sequent \quad$ हे निष्पन्न करते. $\Gamma$ विसंगत नसेल,
म्हणजे प्रत्येक सांत $\Gamma_0 \subseteq \Gamma$ साठी $\Log{LK}$
$\Gamma_0 \Sequent \quad$ हे निष्पन्न करत नसेल, तर त्याला
\emph{सुसंगत} म्हणतो.
\end{defn}

\begin{prop}[परावर्तिता]
\label{mvl:seq:prf:prop:reflexivity}
$\varphi \in \Gamma$ असेल, तर $\Gamma \Proves \varphi$.
\end{prop}

\begin{proof}
$\varphi \Sequent \varphi$ ही आरंभीची क्रमवर्ती निष्पन्न करता येण्याजोगे आहे आणि $\{\varphi\}
\subseteq \Gamma$.
\end{proof}

\begin{prop}[एकस्वनिकता]
\label{mvl:seq:prf:prop:monotonicity}
$\Gamma \subseteq \Delta$ आणि $\Gamma \Proves \varphi$ असेल, तर $\Delta
\Proves \varphi$.
\end{prop}

\begin{proof}
$\Gamma \Proves \varphi$ असे समजा; म्हणजे, एखादा सांत $\Gamma_0
\subseteq \Gamma$ असा आहे की $\Gamma_0 \Sequent \varphi$ ही क्रमवर्ती
निष्पन्न करता येण्याजोगे आहे. $\Gamma \subseteq \Delta$ असल्यामुळे $\Gamma_0$ हा
$\Delta$ चासुद्धा सांत उपसंच आहे. म्हणून $\Gamma_0 \Sequent \varphi$ ची
निष्पत्ती ही $\Delta \Proves \varphi$ हेसुद्धा दाखवते.
\end{proof}

\begin{prop}[संक्रमकता]
\label{mvl:seq:prf:prop:transitivity}
$\Gamma \Proves \varphi$ आणि $\{\varphi\} \cup \Delta \Proves
\psi$ असेल, तर $\Gamma \cup \Delta \Proves \psi$.
\end{prop}

\begin{proof}
$\Gamma \Proves \varphi$ असेल, तर एखादा सांत $\Gamma_0 \subseteq \Gamma$
आणि $\Gamma_0 \Sequent \varphi$ ची निष्पत्ती~$\pi_0$ असते.
$\{\varphi\} \cup \Delta \Proves \psi$ असेल, तर एखाद्या सांत उपसंचासाठी
$\Delta_0 \subseteq \Delta$, $\varphi, \Delta_0 \Sequent \psi$ ची
निष्पत्ती~$\pi_1$ असते. पुढील निष्पत्ती विचारात घ्या:
\begin{prooftree}
  \AxiomC{}
  \RightLabel{$\pi_0$}
  \Deduce$\Gamma_0 \fCenter \varphi$
  \AxiomC{}
  \RightLabel{$\pi_1$}
  \Deduce$\varphi, \Delta_0 \fCenter \psi$
  \RightLabel{\Cut}
  \BinaryInf$\Gamma_0, \Delta_0 \fCenter \psi$
\end{prooftree}
$\Gamma_0 \cup \Delta_0 \subseteq \Gamma \cup \Delta$ असल्यामुळे यावरून
$\Gamma \cup \Delta \Proves \psi$ हे दिसते.
\end{proof}

विशेषतः, याचा अर्थ असा की $\Gamma \Proves \varphi$ आणि $\varphi
\Proves \psi$ असेल, तर $\Gamma \Proves \psi$. तसेच $\varphi_1,
\dots, \varphi_n \Proves \psi$ असेल आणि प्रत्येक~$i$ साठी $\Gamma \Proves \varphi_i$
असेल, तर $\Gamma \Proves \psi$, हेसुद्धा यावरून मिळते.

\begin{prop}
\label{mvl:seq:prf:prop:incons}
$\Gamma$ विसंगत असतो तेव्हा आणि केवळ तेव्हाच, जेव्हा प्रत्येक
  वाक्य~$\varphi$ साठी $\Gamma \Proves {\varphi}$.
\end{prop}

\begin{proof}
सराव.
\end{proof}

\begin{prob}
\ref{fol:seq:ptn:prop:incons} सिद्ध करा.
\end{prob}

\begin{prop}[संहतता]
  \label{mvl:seq:prf:prop:proves-compact}
  \begin{enumerate}
  \item $\Gamma \Proves \varphi$ असेल, तर एखादा सांत उपसंच $\Gamma_0
    \subseteq \Gamma$ असा आहे की $\Gamma_0 \Proves \varphi$.
  \item $\Gamma$ चा प्रत्येक सांत उपसंच सुसंगत असेल, तर $\Gamma$ सुसंगत
    असतो.
  \end{enumerate}
\end{prop}

\begin{proof}
  \begin{enumerate}
    \item $\Gamma \Proves \varphi$ असेल, तर एखादा सांत उपसंच
      $\Gamma_0 \subseteq \Gamma$ असा आहे की $\Gamma_0
      \Sequent \varphi$ या क्रमवर्तीची निष्पत्ती आहे. परिणामी,
      $\Gamma_0 \Proves \varphi$.
    \item $\Gamma$ विसंगत असेल, तर एखादा सांत
      उपसंच~$\Gamma_0 \subseteq \Gamma$ असा आहे की $\Log{LK}$
      $\Gamma_0 \Sequent \quad$ हे निष्पन्न करते. पण मग $\Gamma_0$ हा
      $\Gamma$ चा विसंगत सांत उपसंच आहे.
  \end{enumerate}
\end{proof}



\begin{quote}\small\textbf{स्रोतनोंद.} OLINC-333: हा विभाग अभिजात LK आणि द्विस्थानी क्रमवर्ती वापरतो; आधीच्या बहुमूल्य विभागातील n-स्थानी कलनाची येथे व्याख्या केलेली नाही.\end{quote}










\section{अंकगणिताची अमानक प्रतिमाने}\label{mod:bas:nsa:sec}

\begin{defn}
$\Lang{L_N}$ ही अंकगणिताची भाषा असू द्या. तिच्यात
स्थिरांक $\Obj{0}$, द्विस्थानी विधेय $<$,
एकस्थानी फलन $\prime$, आणि द्विस्थानी फलने
$+$ व $\times$ आहेत.
\begin{enumerate}
\item अंकगणिताचे \emph{मानक प्रतिमान} म्हणजे $\Lang{L_N}$ साठीची
  संरचना $\Struct{N}$: तिचा प्रांत $\Nat = \{ 0, 1, 2, \dots\}$
  असतो; $\Obj{0}$ चा अर्थ $0$, $<$ चा अर्थ $\Nat$ वरील लहान-असण्याचा
  संबंध, तर $\prime$, $+$ आणि $\times$ यांचे अर्थ अनुक्रमे
  $\Nat$ वरील उत्तरवर्ती, बेरीज आणि गुणाकार असे असतात.
\item \emph{सत्य अंकगणित} म्हणजे $\Theory{N}$ ही उपपत्ती.
\end{enumerate}
\end{defn}

$\Lang{L_N}$ मध्ये काम करताना $\Obj{0}$ वर उत्तरवर्ती फलन
$n$ वेळा लावून मिळणाऱ्या $\Obj{0}^{\prime\dots\prime}$
या रूपातील पदाचे संक्षिप्त रूप~$\num{n}$ लिहितो.

\begin{defn}
$\Lang{L_N}$ साठीची संरचना~$\Struct{M}$ ही
$\Struct{N} \iso \Struct{M}$ असेल तेव्हा आणि केवळ तेव्हाच
\emph{मानक} असते.
\end{defn}

\begin{thm}\label{mod:bas:nsa:thm:non-std}
सत्य अंकगणिताची अमानक गणनीय प्रतिमाने आहेत.
\end{thm}

\begin{proof}
$\Lang{L_N}$ मध्ये नवा स्थिरांक~$c$ घालून तिचा विस्तार करा
आणि पुढील उपपत्ती विचारात घ्या:
\[
\Theory{N} \cup \Setabs{\num{n} < c}{n \in \Nat}.
\]
ही उपपत्ती सांत पूर्ततायोग्य आहे. त्यामुळे संहततेनुसार तिला
$\Struct{M}$ हे प्रतिमान आहे; अधोगामी लोव्हेनहाइम--स्कोलेम
प्रमेयानुसार ते गणनीय घेता येते. $\Domain{M}$ हा
$\Struct{M}$ चा प्रांत असताना, $\Lang{L_N}$ मधील अतार्किक
स्थिरांकांचे $\Struct{M}$ मधील अर्थनिर्धारण असे लिहू:
$\mathbf{z} = \Assign{\Obj{0}}{M} \in \Domain{M}$,
${\prec} = \Assign{<}{M} \subseteq \Domain{M}^2$,
$* = \Assign{\prime}{M} \colon \Domain{M} \to \Domain{M}$,
आणि $\oplus = \Assign{+}{M}, \otimes =
\Assign{\times}{M} \colon \Domain{M}^2 \to \Domain{M}$.
प्रत्येक $x \in \Domain{M}$ साठी, $x$ वर~$*$ लावल्याने
मिळणाऱ्या $\Domain{M}$ मधील घटकाला $x^*$ लिहू.


आता, $h$ हे $\Struct{N}$ आणि $\Struct{M}$ यांचे
समरूपण असेल, तर काही $n \in \Nat$ साठी
$h(n) = \Assign{c}{M}$ असते. म्हणून $\Struct{N}$ मध्ये
$s(x) = n$ अशी कोणतीही चलनियुक्ती $s$ घ्या. मग
$\Sat{N}{\eq[\num{n}][x]}[s]$; तसेच
\ref{mod:bas:iso:thm:isom} च्या सिद्धतेनुसार
$\Sat{M}{\eq[\num{n}][x]}[h \circ s]$.
त्यामुळे $\Assign{c}{M} = \mathbf{z}^{*\cdots *}$
(येथे $*$ ही क्रिया $n$ वेळा केली आहे).
पण गृहीतकानुसार $\Sat{M}{\num{n} < c}$ आहे आणि
$\prec$ हा संबंध अपरावर्ती आहे; म्हणून हे अशक्य आहे.
त्यामुळे $\Struct{M}$ अमानक आहे.
\end{proof}

\begin{prob}
संच $X$ वरील संबंध $R$ हा तेव्हा आणि केवळ तेव्हाच
\emph{सुस्थापित} असतो, जेव्हा $R$ मध्ये अनंत अवरोही साखळी
नसते; म्हणजे $\dots x_2Rx_1Rx_0$ अशी अट पूर्ण करणारे
$X$ मधील $x_0$, $x_1$, $x_2$,~\dots नसतात.
झेर्मेलो--फ्रेंकेल संच उपपत्ती~$ZF$ सुसंगत आहे असे मानून,
$ZF$ ची असुस्थापित प्रतिमाने आहेत हे दाखवा; म्हणजे
$\dots x_2 \in x_1 \in x_0$ असणारी प्रतिमाने
$\mathfrak{M}$ आहेत हे सिद्ध करा.
\end{prob}

\ref{mod:bas:nsa:thm:non-std} मधील अमानक प्रतिमान
$\Struct{M}$ हे मानक प्रतिमानाशी प्राथमिक सममूल्य असल्यामुळे
$\Struct{M}$ चे अनेक गुणधर्म काढता येतात. या विभागाचा
उर्वरित भाग अशा गुणधर्मांच्या अभ्यासासाठी आहे; त्यातून
$\Theory{N}$ च्या अमानक गणनीय प्रतिमानांचे
अचूक लक्षण मिळेल.

\begin{enumerate}
\item $\Domain{M}$ मधील कोणताही घटक स्वतःहून $\prec$-लहान
  नसतो: $\lforall[x][\lnot x < x]$ हे वाक्य $\Struct{N}$
  मध्ये सत्य असल्यामुळे $\Struct{M}$ मध्येही सत्य आहे.
\item याच रीतीने $\prec$ हा $\Domain{M}$ वरील
  \emph{रेषीय क्रम} आहे, म्हणजे तो $\Domain{M}$ वरील
  पूर्ण, अपरावर्ती आणि संक्रमक संबंध आहे.
\item $\mathbf{z}$ हा $\Domain{M}$ मधील $\prec$-लघुतम घटक आहे.
\item $\Domain{M}$ मधील प्रत्येक घटक त्याच्या $*$-उत्तरवर्तीपेक्षा
  $\prec$-लहान असतो; आणि $x^*$ हा $x$ पेक्षा मोठ्या
  $\Domain{M}$ मधील घटकांपैकी $\prec$-लघुतम असतो.
\item $\Struct{M}$ मध्ये $\Nat$ शी समरूप असा $\prec$ चा
  आरंभीचा खंड आहे: $\mathbf{z}, \mathbf{z}^*,
  \mathbf{z}^{**}, \dots$. त्याला $\Domain{M}$ चा
  \emph{मानक भाग} म्हणतो. $\Domain{M}$ मधील इतर प्रत्येक
  घटक \emph{अमानक} असतो. $\Domain{M}$ मध्ये अमानक
  घटक असलेच पाहिजेत;
  अन्यथा \ref{mod:bas:nsa:thm:non-std} च्या सिद्धतेतील $h$ हे
  समरूपण ठरेल. $n, m, \dots$ ही चले आपण
  $\Struct{M}$ च्या या मानक भागावर चालणारी मानतो.
\item प्रत्येक अमानक घटक प्रत्येक मानक घटकापेक्षा मोठा
  असतो; कारण प्रत्येक $n \in \Nat$ साठी
  \[
  \Sat{N}{\lforall[z][(\lnot(\eq[z][\Obj{0}] \lor
  \dots \lor \eq[z][\num{n}]) \lif \num{n} < z)]},
  \]
  म्हणून $z \in \Domain{M}$ सर्व मानक घटकांहून
  वेगळा असेल, तर तो त्या सर्वांपेक्षा \emph{मोठा} असतो.
\item $\mathbf{z}$ सोडून $\Domain{M}$ मधील प्रत्येक
  घटक हा एखाद्या एकमेव $\Domain{M}$ मधील घटकाचा
  $*$-उत्तरवर्ती असतो; त्या पूर्ववर्तीला $^*x$ लिहितो.
  $x = y^*$ असेल, तर $x$ आणि $y$ यांपैकी एक मानक
  असल्यास दोन्ही मानक असतात (आणि एक अमानक असल्यास
  दोन्ही अमानक असतात).
\item $\Domain{M}$ वर $\approx$ हा सममूल्यता संबंध
  परिभाषित करा: एखाद्या \emph{मानक} $n$ साठी
  $x \oplus n = y$ किंवा $y \oplus n =x$
  असेल तेव्हा आणि केवळ तेव्हाच $x \approx y$.
  दुसऱ्या शब्दांत, $x$ आणि $y$ यांच्यात सांत अंतर
  असेल तेव्हा आणि केवळ तेव्हाच $x \approx y$.
  $n$ आणि $m$ मानक असतील, तर $n \approx m$.
  $x$ चा \emph{खंड} म्हणजे सममूल्यता वर्ग
  $[x] = \Setabs{y}{x \approx y}$.
\item $x \prec y$ पण $x \not\approx y$ असे समजा.
  $\Sat{N}{\lforall[x][\lforall[y][(x < y \lif (x' < y \lor x' =
      y))]]}$ असल्यामुळे, $x^* \prec y$ किंवा
  $x^* = y$. दुसरी शक्यता $x \approx y$ सूचित
  करत असल्यामुळे अशक्य आहे; म्हणून
  $x^* \prec y$. त्याचप्रमाणे $x \prec y$ आणि
  $x \not\approx y$ असतील, तर $x \prec {^*y}$.
  म्हणून $x \prec y$ आणि $x \not\approx y$ असतील,
  तर प्रत्येक $w \approx x$ हा प्रत्येक $v \approx y$
  पेक्षा $\prec$-लहान असतो. त्यानुसार, प्रत्येक
  अमानक खंड $[x]$ पुढील द्विदिश अनंत साखळी बनवतो:
  \[
  \cdots \prec  {^{**}x} \prec {^*}x \prec x \prec x^* \prec x^{**}
  \prec \cdots
  \]
  पूर्णांकांच्या क्रमप्रकारासारखा असल्यामुळे तिला
  $Z$-साखळी म्हणतात.


\item $\prec$ हा क्रम खंडांवरही नेऊ शकतो: $x \prec y$
  आणि ते भिन्न खंडांत असतील, तर $x$ चा खंड
  $y$ च्या खंडापेक्षा लहान असतो. अमानक घटक असणारा
  खंड \emph{अमानक} असतो. मानक खंड हा लघुतम खंड आहे.

\item लघुतम अमानक खंड नाही: $y$ अमानक असेल,
  तर $x \prec y$ आणि $x \not\approx y$ अशी अट
  पूर्ण करणारा अमानक $x$ असतो. सिद्धता: मानक
  प्रतिमान $\Struct{N}$ मध्ये प्रत्येक संख्येला दोनने
  भाग जातो, कदाचित बाकी एक उरते; म्हणजे
  $\Sat{N}{\lforall[y][\lexists[x][(y = x + x \lor y = x + x +
        \Obj{0}')]]}$. प्राथमिक सममूल्यतेमुळे
  प्रत्येक $y \in \Domain{M}$ साठी असा
  $x \in \Domain{M}$ असतो की $x \oplus x = y$
  किंवा $x \oplus x \oplus \mathbf{z}^*= y$.
  $x$ मानक असेल, तर $y$ सुद्धा मानक असेल;
  म्हणून $x$ अमानक आहे. शिवाय $x$ आणि $y$
  वेगळ्या खंडांत येतात, म्हणजे $x \not\approx y$.
  तसे नसल्यास, एखाद्या मानक $n$ साठी
  $x \oplus n = y$ असेल. $y = x \oplus x$
  असेल, तर $x \oplus n = x \oplus x$; बेरीज
  रद्द करण्याच्या नियमानुसार $x = n$ येईल.
  हा नियम $\Struct{N}$ मध्ये आणि म्हणून
  $\Struct{M}$ मध्येही चालतो; पण मग $x$
  मानक ठरेल. $y = x \oplus x \oplus \mathbf{z}^*$
  या प्रसंगातही तसेच होते.

\item याच प्रकारच्या युक्तिवादाने सर्वांत मोठा खंडही नाही.
\item खंडांचा क्रम सघन आहे: $[x]$ हा $[y]$
  पेक्षा लहान असेल ($x \not\approx y$), तर
  त्यांच्या मध्ये दोन्हींपेक्षा वेगळा खंड $[z]$
  असतो. $x \prec y$ असे समजा. आधीप्रमाणे
  $x \oplus y$ ला दोनने भाग जातो (कदाचित
  बाकी उरते); म्हणून काही $u \in \Domain{M}$
  साठी $x \oplus y = u \oplus u$ किंवा
  $x \oplus y = u \oplus u \oplus \mathbf{z}^*$.
  $u$ हा $x$ आणि $y$ यांचा मध्य आहे आणि
  त्यांच्यामध्ये येतो. $x \oplus y = u \oplus u$
  असे समजा (दुसरा प्रसंगही तसाच): $u \approx x$
  असेल, तर एखाद्या मानक $n$ साठी
  \[
  x \oplus y = x \oplus n \oplus x \oplus n,
  \]
  म्हणून $y = x \oplus n \oplus n$; यातून
  $x \approx y$ येते, जे गृहीतकाविरुद्ध आहे.
  म्हणून $u \not\approx x$. त्याच युक्तिवादाने
  $u \not\approx y$.
\end{enumerate}
त्यामुळे अमानक खंडांचा क्रम परिमेय संख्यांप्रमाणे
आहे: तो अंत्यबिंदू नसलेला गणनीय रेषीय क्रम आहे.
यावरून सत्य अंकगणिताच्या कोणत्याही दोन
गणनीय अमानक प्रतिमानांचे,
$\Struct{M}_1$ आणि $\Struct{M_2}$ यांचे,
फक्त $<$ आणि $=$ असणाऱ्या भाषेतील
न्यूनीकृत प्रतिमान समरूप असतात.
खरोखर, समरूपण $h$ असे बनवता येते:
$\Struct{M_1}$ आणि $\Struct{M_2}$ यांचे
मानक भाग मानक प्रतिमान $\Struct{N}$ शी,
म्हणून एकमेकांशी, समरूप असतात. अमानक भाग
बनवणारे खंड परिमेयांप्रमाणे क्रमबद्ध असतात;
म्हणून \ref{mod:bas:dlo:thm:cantorQ} नुसार ते
समरूप असतात. प्रत्येक खंडातील कोणतेही
घटक जुळवून खंडांचे समरूपण खंडांच्या
\emph{आत} विस्तारता येते; मग
$\Struct{M_1}$ मधील $x$ च्या उत्तरवर्तीची
प्रतिमा ही $\Struct{M_2}$ मधील $x$ च्या
प्रतिमेचा उत्तरवर्ती ठरवता येते.
यावरून मात्र $\mathfrak{M}_1$ आणि
$\mathfrak{M}_2$ ही संपूर्ण अंकगणितीय
भाषेत समरूप आहेत असे \emph{निघत नाही}
(कारण समरूपता नेहमी चिन्हसंचाच्या संदर्भातच
असते). $\Domain{M_1}$ आणि $\Domain{M_2}$
वर बेरीज व गुणाकार परिभाषित करण्याचे
असमरूप मार्ग आहेत. एका संपूर्ण उपपत्तीच्या
गणनीय प्रतिमानांची संख्या नेमकी~2
असू शकत नाही, या वॉट यांच्या प्रसिद्ध
प्रमेयावरूनही हे निघते.

\begin{prob}
अंकगणिताच्या अमानक प्रतिमानात सर्वांत मोठा
खंड असू शकत नाही हे दाखवा.
\end{prob}

\begin{prob}
$\Lang{L}$ ही $<$ हा एकच विधेय
($\eq$ व्यतिरिक्त) असणारी प्रथम-क्रमाची
भाषा असू द्या आणि
$\Struct{N} = (\Nat, <)$ असू द्या.
$\Struct{N}$ चे सर्व सांत किंवा सहसांत
उपसंच परिभाषणीय आहेत. \emph{हेच एकमेव}
$\Struct{N}$ चे परिभाषणीय उपसंच आहेत हे दाखवा.

(सूचना: प्रथम, $\Obj{prc}(x,y)$ हे
$\Lang{L}$ मधील “$x$ हा $y$ चा लगतचा
पूर्ववर्ती आहे” याचे संक्षिप्त रूप असलेले
सूत्र असू द्या:
\[
x<y \land \lnot \lexists[z][(x<z \land z < y)].
\]
आता $\Struct{N}$ च्या कोणत्याही परिभाषणीय
उपसंचाला $\Lang{L}$ मधील सूत्र $\varphi(x)$
अनुरूप असते. अशा प्रत्येक $\varphi$ साठी
पुढील वाक्य $\theta$ विचारात घ्या:
\[
\lexists[x][\lforall[y][\lforall[z][((x<y \land x<z \land \Obj{prc}(y,z)
  \land \varphi(y)) \lif \varphi(z))]]].
\]
$\Sat{N}{\theta}$ असेल तेव्हा आणि केवळ
तेव्हाच $\varphi$ ने परिभाषित केलेला
$\Struct{N}$ चा उपसंच सांत किंवा
सहसांत असतो हे दाखवा.

आता $\Struct{M}$ हे $\Struct{N}$
शी प्राथमिक सममूल्य अमानक प्रतिमान असू द्या.
$a \in \Domain{M}$ अमानक असेल, तर
$b, c \in \Domain{M}$ हे $a$ पेक्षा
मोठे घ्या आणि $b$ हा~$c$ चा लगतचा
पूर्ववर्ती असू द्या. मग $\Domain{M}$
चे असे स्वसमरूपण $h$ आहे की $h(b)=c$
(का?). त्यामुळे $b$ ने $\varphi$ ची पूर्ति
केली, तर $c$ नेही ती केली पाहिजे (का?).
म्हणून $\theta$ हे $\Struct{M}$ मध्ये सत्य
आहे आणि त्यामुळे $\Struct{N}$ मध्येही.
पण याचा अर्थ असा की $\varphi$ ने परिभाषित
केलेला $\Struct{N}$ चा उपसंच सांत
किंवा सहसांत आहे.
\end{prob}



\begin{quote}\small\textbf{स्रोतनोंद.} SOL6-A293: अमानक प्रतिमानाच्या भाषेसाठी L\_N आणि संबंधाच्या प्रांतासाठी Domain(M) वापरले आहेत. वेगळ्या खंडांतील घटकांच्या क्रमात x चा उत्तरवर्तीही y पेक्षा लहान असतो, हा मिळालेला निष्कर्ष आता स्पष्ट आहे.\end{quote}
\begin{quote}\small\textbf{स्रोतनोंद.} OLINC-334: SOL6-A293 मध्ये आधी निश्चित केलेली अंकगणिताची भाषा L\_N आणि प्रतिमानाचा प्रांत Domain(M) स्पष्ट वापरला आहे; ऐतिहासिक विसंगत संकेत स्रोतामध्ये जतन आहेत.\end{quote}
\begin{quote}\small\textbf{स्रोतनोंद.} OLINC-336: गोठवलेल्या स्रोताचा 'प्रत्येक खंड' हा दावा मानक खंडासाठी चुकीचा आहे; द्विदिश अनंत साखळी फक्त अमानक खंडांना लागू होते.\end{quote}
\begin{quote}\small\textbf{स्रोतनोंद.} OLINC-337: SOL6-A293 मध्ये या प्रसंगाचा प्रत्यक्ष निष्कर्ष x चा उत्तरवर्ती y पेक्षा लहान आहे असा केला आहे; आधीचेच गृहीतक पुन्हा सांगणारे मूळ वाक्य दुरुस्त केले आहे.\end{quote}
\begin{quote}\small\textbf{स्रोतनोंद.} OLINC-338: भिन्न खंडांसाठी x आणि y असण्याची अट स्रोताच्या पहिल्या वाक्यात गाळलेली आहे; एकाच खंडातील x<y त्यांचा खंड स्वतःहून लहान करत नाही.\end{quote}
\begin{quote}\small\textbf{स्रोतनोंद.} OLINC-335: स्रोताच्या अर्धे करण्याच्या सूत्रात प्रत्येक y साठी 'प्रत्येक x' असे चुकीने आहे; अस्तित्वदर्शक x आवश्यक आहे आणि येथे दुरुस्त केला आहे.\end{quote}










\section{सामान्य मोडल तर्कशास्त्रे}\label{nml:syn:nor:sec}

मोडल तर्कशास्त्र म्हणजे योग्य अटी पूर्ण करणारा
सूत्रांचा संच असे मानता येते. काही मोडल
तर्कशास्त्रे विशेष महत्त्वाची आहेत: त्यांची अर्थलक्षणे
आणि उपयोग रंजक आहेत, आणि विशेषतः त्यांचा
सैद्धान्तिक अभ्यास चांगला झालेला आहे. यांनाच
सामान्य मोडल तर्कशास्त्रे म्हणतात.

\begin{defn}\label{nml:syn:nor:defn:modal-logic}
\emph{मोडल तर्कशास्त्र}~$\Log L$ म्हणजे मूलभूत
मोडल भाषेतील सूत्रांचा असा संच की
\begin{enumerate}
\item त्यात सर्व विधानात्मक सर्वत:सत्य सूत्रे असतात;
\item तो एकरूप आदेशनाखाली बंद असतो; आणि
\item तो मोडस पोनेन्सखाली बंद असतो; म्हणजे $\varphi$ आणि
  $(\varphi \lif \psi) \in \Log L$ असतील, तर
  $\psi \in \Log L$ सुद्धा असते.
\end{enumerate}
\end{defn}

\begin{ex}
फारसे रंजक नसले तरी ही व्याख्या पूर्ण करणारे,
म्हणून मोडल तर्कशास्त्रे असणारे, सूत्रे
चे काही संच असे आहेत:
\begin{enumerate}
\item सर्व सर्वत:सत्य सूत्रांचा संच;
\item सर्व सूत्रांचा संच (विसंगत मोडल तर्कशास्त्र).
\end{enumerate}
\end{ex}

\begin{defn}\label{defn:normal-modal-logic}
मोडल तर्कशास्त्र~$\Log L$ हे तेव्हा आणि केवळ
तेव्हाच \emph{सामान्य} असते, जेव्हा
\begin{enumerate}
\item त्यात पुढील दोन्ही असतात:
\begin{align}
\tag{K} \Box(p \lif q) & \lif (\Box p \lif \Box q) \\
\tag{Dual} \Diamond p & \liff \lnot \Box \lnot p
\end{align}
\item आणि ते अनिवार्यीकरणाखाली बंद असते; म्हणजे
  $\varphi \in \Log L$ असेल, तर
  $\Box \varphi \in \Log L$ सुद्धा असते.
\end{enumerate}
\end{defn}

मोडल तर्कशास्त्रे परिभाषित करण्याचे दोन प्रमुख
मार्ग आपण पाहू. पहिला सिद्धता-उपपत्तीय: विशिष्ट
निष्पत्ती-पद्धतींमध्ये निष्पन्न करता येणाऱ्या
सूत्रांचा संच. दुसरा प्रतिमान-उपपत्तीय:
प्रतिमानांच्या विशिष्ट वर्गांत वैध असणाऱ्या
सूत्रांचे संच. साध्या विधानात्मक तर्कशास्त्रात
(किंवा प्रथम-क्रम तर्कशास्त्रात) हाच प्रकार दिसतो.
तेथेही सूत्रांचा संच दोन प्रकारे ठरवता येतो:
सर्व सत्यतामूल्य-नियुक्त्यांत सत्य असलेल्या
सर्वत:सत्य सूत्रांचा संच म्हणून, किंवा एखाद्या
निष्पत्ती-पद्धतीतील सर्व निष्पन्न करता येण्याजोगे
सूत्रांचा संच म्हणून. सामान्य मोडल
तर्कशास्त्रांतही संबंधीय प्रतिमाने आणि स्वयंसिद्धकी
(हिल्बर्ट-प्रकारच्या) सिद्धता पद्धती वापरून हे
दुहेरी वर्णन शक्य होते.












\section{क्रमवर्ती-रूप नैसर्गिक निगमन}\label{int:pty:snd:sec}

$\Gamma$ ही मुक्त न केलेली गृहीतके आणि~$\varphi$ हा
निष्कर्ष असलेली नैसर्गिक निगमनाची निष्पत्ती असेल,
तर $\Gamma \Sequent \varphi$ असे लिहू;~$\Gamma$ रिकामे
असेल, तर $\Sequent \varphi$ असे लिहू.

$\Gamma \cup \{\varphi_1, \dots, \varphi_n\}$ साठी
$\Gamma, \varphi_1, \dots, \varphi_n$ आणि $\Gamma \cup \Delta$
साठी $\Gamma, \Delta$ असे लिहू.

$\Gamma \Sequent \varphi \land \psi$ असेल, म्हणजे~$\Gamma$
ही मुक्त न केलेली गृहीतके आणि $\varphi \land \psi$ हे
अंतिम सूत्र असलेली निष्पत्ती असेल,
तर तळाशी $\Elim{\land}$ लावून तीच
मुक्त न केलेली गृहीतके आणि~$\varphi$ हा निष्कर्ष
असलेली निष्पत्ती मिळते. दुसऱ्या शब्दांत,
$\Gamma \Sequent \varphi \land \psi$ असेल, तर
$\Gamma \Sequent \varphi$.
\begin{prooftree}
  \Axiom$\Gamma \fCenter \varphi \land \psi$
  \RightLabel{$\Elim{\land}$}
  \UnaryInf$\Gamma \fCenter \varphi$
  \DisplayProof\qquad\bottomAlignProof
  \Axiom$\Gamma \fCenter \varphi \land \psi$
  \RightLabel{$\Elim{\land}$}
  \UnaryInf$\Gamma \fCenter \psi$
\end{prooftree}
$\Elim{\land}$ ही खूण नैसर्गिक निगमनातील त्याच
नावाच्या नियमाशी असलेला संबंध सूचित करते.

त्याचप्रमाणे $\Gamma, \varphi \Sequent \psi$ असेल,
म्हणजे $\Gamma, \varphi$ ही मुक्त न केलेली गृहीतके
आणि~$\psi$ हे अंतिम सूत्र असलेली
निष्पत्ती असेल, तर $\Intro{\lif}$ नियम लावून
$\Gamma$ ही मुक्त न केलेली गृहीतके आणि
$\varphi \lif \psi$ हे अंतिम सूत्र असलेली
निष्पत्ती मिळते; म्हणजे
$\Gamma \Sequent \varphi \lif \psi$.
यात गृहीतकांचे मुक्त करणे अधिक स्पष्ट झाले आहे.
\begin{prooftree}
  \Axiom $\Gamma, \varphi \fCenter \psi$
  \RightLabel{$\Intro{\lif}$}
  \UnaryInf $\Gamma \fCenter \varphi \lif \psi$
\end{prooftree}

इतर नियमांवरूनही अशाच प्रकारे निष्कर्ष काढता
येतात. ते पुढीलप्रमाणे स्पष्ट केले आहेत:
\begin{gather*}
  \Axiom $\Gamma \fCenter \varphi$
  \Axiom $\Delta \fCenter \psi$
  \RightLabel{$\Intro{\land}$}
  \BinaryInf $\Gamma,\Delta \fCenter \varphi \land \psi$
  \DisplayProof\\
  \Axiom $\Gamma \fCenter \varphi \land \psi$
  \RightLabel{$\Elim{\land}_1$}
  \UnaryInf $\Gamma \fCenter \varphi$
  \DisplayProof
  \qquad
  \Axiom $\Gamma \fCenter \varphi \land \psi$
  \RightLabel{$\Elim{\land}_2$}
  \UnaryInf $\Gamma \fCenter \psi$
  \DisplayProof
  \\
  \Axiom $\Gamma \fCenter \varphi$
  \RightLabel{$\Intro{\lor}_1$}
  \UnaryInf $\Gamma \fCenter \varphi \lor \psi$
  \DisplayProof\qquad
  \Axiom $\Gamma \fCenter \psi$
  \RightLabel{$\Intro{\lor}_2$}
  \UnaryInf $\Gamma \fCenter \varphi \lor \psi$
  \DisplayProof\\
  \Axiom $\Gamma \fCenter \varphi \lor \psi$
  \Axiom $\Delta, \varphi \fCenter \chi$
  \Axiom $\Delta', \psi \fCenter \chi$
  \RightLabel{$\Elim{\lor}$}
  \TrinaryInf $\Gamma, \Delta, \Delta' \fCenter \chi$
  \DisplayProof
  \\
  \Axiom $\Gamma, \varphi \fCenter \psi$
  \RightLabel{$\Intro{\lif}$}
  \UnaryInf $\Gamma \fCenter \varphi \lif \psi$
  \DisplayProof
  \qquad
  \Axiom $\Delta \fCenter \varphi \lif \psi$
  \Axiom $\Gamma \fCenter \varphi$
  \RightLabel{$\Elim{\lif}$}
  \BinaryInf $\Gamma, \Delta \fCenter \psi$
  \DisplayProof
  \\
  \Axiom $\Gamma \fCenter \lfalse$
  \RightLabel{$\FalseInt$}
  \UnaryInf $\Gamma \fCenter \varphi$
  \DisplayProof
\end{gather*}

कोणतेही गृहीतक स्वतःच~$\varphi$ वरून~$\varphi$ ची
निष्पत्ती असते; म्हणजे $\varphi \Sequent \varphi$
हे नेहमीच असते.

\begin{prooftree}
  \AxiomC{$ $}
  \UnaryInfC{$\varphi \fCenter \varphi$}
\end{prooftree}

हे नियम एकत्रितपणे कोणत्या नैसर्गिक निगमनाच्या
निष्पत्त्या अस्तित्वात आहेत यांचे कलन म्हणून
घेता येतात. त्यांना नैसर्गिक निगमनाचे चिन्हांकनात्मक
रूपांतरही मानता येते; प्रत्येक टप्प्यावर कोणते
सूत्र निष्पन्न झाले आहे एवढेच नव्हे,
तर ते कोणत्या मुक्त न केलेली गृहीतकांवरून
निष्पन्न झाले आहे हेही त्यात नोंदते.

\begin{prooftree}
  \Axiom$\varphi \fCenter \varphi$
  \UnaryInf $\varphi \fCenter \varphi \lor (\varphi \lif \lfalse)$
  \Axiom $\psi \fCenter \psi$
  \BinaryInf$\varphi, \psi \fCenter \lfalse$
  \UnaryInf$\psi \fCenter \varphi \lif \lfalse$
  \UnaryInf$\psi \fCenter \varphi \lor (\varphi \lif \lfalse)$
  \Axiom$\psi \fCenter \psi$
  \BinaryInf$\psi \fCenter \lfalse$
  \UnaryInf$\fCenter \psi \lif \lfalse$
\end{prooftree}
येथे~$\psi$ हे $(\varphi \lor (\varphi \lif \lfalse))\lif \lfalse$
याचे संक्षिप्त रूप आहे.



\begin{quote}\small\textbf{स्रोतनोंद.} SOL6-A325: अंतिम उदाहरणातील संदर्भात चुकून आलेले अपूर्ण अभिव्यंजन-चिन्ह व कंस काढले; दिलेला B चा अर्थ आणि सिद्धतेचा निष्कर्ष जपला.\end{quote}











\section{विधानात्मक तर्कशास्त्राची संपूर्णता}\label{pl:prp:com:sec}

\begin{defn}
\label{pl:prp:com:def:MaxCon}
सूत्रांचा संच~$\Gamma$ \emph{महत्तम सुसंगत} असतो, जर तो
सुसंगत असेल आणि $\Gamma \subseteq \Delta$ असा कोणताही सुसंगत
संच~$\Delta$ घेतल्यास $\Gamma = \Delta$ असेल.
\end{defn}

\begin{lem}[सत्यता पूर्वप्रमेय]
  \label{pl:prp:com:lem:truth}
$\Gamma$ महत्तम सुसंगत असू द्या; मग:
\begin{enumerate}
\item \label{pl:prp:com:lem:truth-1} $\varphi \in \Gamma$ तेव्हा आणि केवळ तेव्हा $\lnot \varphi
  \notin \Gamma$;
\item \label{pl:prp:com:lem:truth-2} $\Gamma \Proves \varphi$ तेव्हा आणि केवळ तेव्हा $\varphi \in \Gamma$;
\item \label{pl:prp:com:lem:truth-3} $\varphi \lif \psi \in \Gamma$ तेव्हा आणि केवळ तेव्हा
  $\varphi \notin \Gamma$ किंवा $\psi \in \Gamma$.
\end{enumerate}
\end{lem}

\begin{proof}
\begin{enumerate}
\item $\varphi \in \Gamma$ असे माना. $\lnot \varphi \in \Gamma$ सुद्धा असेल,
  तर $\Gamma$ विसंगत ठरतो. $\varphi$ आणि $\lnot\varphi$ यांपैकी एकही
  $\Gamma$ मध्ये नसेल, तर \readerexternalref{pl:prp:axd:prop:phi} नुसार
  $\Gamma\cup\{\varphi\}$ आणि $\Gamma \cup \{\lnot \varphi\}$ यांपैकी
  एक संच सुसंगत असतो; म्हणून $\Gamma$ महत्तम नसतो. उलटी दिशाही
  याचप्रमाणे सिद्ध होते.
\item $\varphi \in \Gamma$ असेल, तर अर्थातच $\Gamma \Proves \varphi$.
  उलट, $\Gamma \Proves \varphi$ असे माना. $\varphi \notin \Gamma$ असेल,
  तर \ref{pl:prp:com:lem:truth-1} नुसार $\lnot \varphi \in \Gamma$. त्यामुळे
  $\Gamma \Proves \varphi$ आणि $\Gamma \Proves \lnot \varphi$, हा व्याघात.
\item सरावासाठी.
\end{enumerate}
\end{proof}

\ref{pl:prp:com:lem:truth}\ref{pl:prp:com:lem:truth-2} ची डावीकडून उजवीकडील
दिशा दाखवते की $\Gamma$ \emph{निगमनतः बंद} आहे: म्हणजे
कोणत्याही~$\varphi$ साठी $\Gamma \Proves \varphi$ असेल, तर $\varphi
\in \Gamma$.

\begin{prob}
  \ref{pl:prp:com:lem:truth} ची सिद्धता पूर्ण करा.
\end{prob}

\begin{thm}[संपूर्णता]
\label{pl:prp:com:thm:completeness}
$\Gamma$ सुसंगत असेल, तर $\Gamma$ पूर्ततायोग्य असतो.
\end{thm}

\begin{proof}
$\varphi_0$, $\varphi_1$,~\dots ही भाषेतील सर्व सूत्रांची
पूर्ण यादी असू द्या. पुढीलप्रमाणे सूत्रांच्या संचांची
$\Gamma_0$, $\Gamma_1$,~\dots ही वाढती क्रमिका
आवर्ती रीतीने परिभाषित करू:
\begin{align*}
\Gamma_0 & = \Gamma\\
\Gamma_{n+1} &  =
\begin{cases}
  \Gamma_n \cup \{\varphi_n\} & \text{जर $\Gamma_n \cup \{\varphi_n\}$ सुसंगत असेल;}\\
  \Gamma_n \cup \{\lnot \varphi_n\} & \text{अन्यथा}.
\end{cases}
\end{align*}
मग पुढील व्याख्या करू:
\[
\Gamma^* = \bigcup_{0\le n}\Gamma_n.
\]
आता क्रमाने पुढील तथ्ये सिद्ध करायची आहेत:
\begin{enumerate}
\item प्रत्येक $n$ साठी $\Gamma_n$ हा संच सुसंगत आहे
  ($n$ वर विगमनाने, \readerexternalref{pl:prp:axd:prop:phi} वापरून);
\item $\Gamma^*$ सुसंगत आहे;
\item $\Gamma^*$ महत्तम आहे.
\end{enumerate}
मग $p_i \in \Gamma^*$ असेल तेव्हाच $v(p_i) = \True$
असे ठरवून सत्य-मूल्यांकन~$v$ परिभाषित करू. $\varphi$ वर विगमन
केल्यास अधिक गुंतागुंतीच्या वाक्ये साठीही
$\Gamma^*$ मधील सदस्यत्व आणि $v$ नुसार सत्यता जुळतात,
हे दिसते: $\overline{v}(\varphi) = \True$ तेव्हा आणि केवळ तेव्हा
$\varphi \in \Gamma^*$. विशेषतः, $v$ हे $\Gamma^*$ ची पूर्तता
करते. $\Gamma \subseteq \Gamma^*$ असल्याने $\Gamma$
सुद्धा पूर्ततायोग्य आहे, हेच सिद्ध करायचे होते.
\end{proof}

\begin{cor}
$\Gamma \Entails \varphi$ असेल, तर $\Gamma \Proves \varphi$.
\end{cor}

\begin{proof}
$\Gamma\Proves/ \varphi$ असेल, तर \readerexternalref{pl:prp:axd:prop:prov-incons} नुसार
$\Gamma \cup \{\lnot \varphi\}$ सुसंगत आहे. म्हणून संपूर्णता
प्रमेयानुसार $\Gamma \cup \{\lnot \varphi\}$ पूर्ततायोग्य आहे;
त्यामुळे $\Gamma \Entails/ \varphi$.
\end{proof}

\begin{prop}[संहतता प्रमेय]
\label{pl:prp:com:prop:compactness}
$\Gamma$ पूर्ततायोग्य असेल तेव्हा आणि केवळ तेव्हा $\Gamma$ चा प्रत्येक
\emph{सान्त} उपसंच~$\Gamma_0$ पूर्ततायोग्य असेल.
\end{prop}

\begin{proof}
$\Gamma$ अपूर्ततायोग्य असेल तेव्हा आणि केवळ तेव्हा तो विसंगत असेल;
तेव्हा आणि केवळ तेव्हा $\Gamma$ चा एखादा सान्त उपसंच~$\Gamma_0$
विसंगत असेल; आणि तेव्हा आणि केवळ तेव्हा $\Gamma$ चा एखादा सान्त
उपसंच~$\Gamma_0$ अपूर्ततायोग्य असेल.
\end{proof}

\begin{cor}
$\Gamma \Entails \varphi$ असेल तेव्हा आणि केवळ तेव्हा $\Gamma$ च्या
एखाद्या सान्त उपसंच~$\Gamma_0$ साठी $\Gamma_0 \Entails \varphi$ असेल.
\end{cor}



\begin{quote}\small\textbf{स्रोतनोंद.} SOL6-A338: द्विदिशा दर्शवणारा वाक्प्रयोग आणि विगमनाची संज्ञा सुसंगत केली; संक्षिप्त सिद्धता आणि सराव जपले.\end{quote}











\section{विधानात्मक तर्कशास्त्राची निर्दोषता}\label{pl:prp:snd:sec}

\begin{thm}[निर्दोषता]
\label{pl:prp:snd:thm:soundness}
$\Gamma\Proves\varphi$ असेल, तर
$\Gamma \Entails \varphi$.
\end{thm}

\begin{proof}
प्रमेयांवर विगमन करू. $\varphi$ स्वयंसिद्धक असेल, तर $\Entails
\varphi$; म्हणून विशेषतः $\Gamma \Entails\varphi$. तसेच $\varphi \in \Gamma$
असेल, तर $\Gamma \Entails\varphi$. विगमनात्मक पायरीसाठी $\chi$
आणि $\chi\lif\varphi$ यांपासून \emph{विध्यात्मक अनुमानाने} $\varphi$
मिळाले असे माना. तसेच विगमनाच्या गृहीतकानुसार प्रमेय
$\chi\lif \varphi$ आणि $\chi$ यांसाठी खरे आहे असे माना.
\ref{pl:syn:sem:prop:semanticalfacts} मधील तिसऱ्या मुद्द्यावरून
अपेक्षित निष्कर्ष मिळतो.
\end{proof}

\begin{cor}
$\Gamma$ पूर्ततायोग्य असेल, तर $\Gamma$ सुसंगत असतो. म्हणून
विधानात्मक तर्कशास्त्र सुसंगत आहे.
\end{cor}



\begin{quote}\small\textbf{स्रोतनोंद.} SOL6-A339: निर्दोषतेच्या सिद्धतेतील विगमनाची संज्ञा सुसंगत केली; सर्व गणिती आधारप्रसंग व पायरी जपली.\end{quote}










\section{द्वितीय-क्रम तर्कशास्त्राची भाषा}\label{sol:syn:lan:sec}

प्रथम-क्रम तर्कशास्त्राप्रमाणेच द्वितीय-क्रम तर्कशास्त्रातील
अभिव्यक्ती मूलभूत शब्दसंग्रहापासून बनतात. त्यात \emph{चले},
\emph{स्थिरांक}, \emph{विधेये} आणि कधी कधी
\emph{फलने} असतात. यांच्याबरोबर तार्किक संकारके,
संख्यक आणि कंस व स्वल्पविरामांसारखी विरामचिन्हे वापरून
\emph{पदे} आणि \emph{सूत्रे} तयार होतात. फरक असा की
वस्तूंसाठीच्या चलांव्यतिरिक्त द्वितीय-क्रम तर्कशास्त्रात संबंध
आणि फलने यांसाठीही चले असतात आणि त्यांच्यावर संख्यापन
करता येते. म्हणून द्वितीय-क्रम तर्कशास्त्राची तार्किक चिन्हे
म्हणजे प्रथम-क्रम तर्कशास्त्राची चिन्हे, तसेच पुढील चिन्हे:

\begin{enumerate}
\item प्रत्येक स्थानसंख्या~$n$ साठी द्वितीय-क्रम संबंध चलांचा
  गणनीय अनंत संच: $\Obj V_0^n$, $\Obj V_1^n$, $\Obj V_2^n$, \dots
\item द्वितीय-क्रम फलन चलांचा गणनीय अनंत संच:
  $\Obj u_0^n$, $\Obj u_1^n$, $\Obj u_2^n$, \dots
\end{enumerate}

प्रथम-क्रम तर्कशास्त्रात एकरूपता~$\eq$ सहसा समाविष्ट
असते. प्रथम-क्रम तर्कशास्त्रात $\Lang{L}$ या भाषेतील
अतार्किक चिन्हे आपल्याला काही अर्थपूर्ण गोष्ट व्यक्त करता
यावी यासाठी अत्यावश्यक असतात. अतार्किक चिन्हे नसलेली
वाक्ये नक्कीच आहेत, पण केवळ~$\eq$ वापरून
फारसे अर्थपूर्ण काही सांगता येत नाही. द्वितीय-क्रम
तर्कशास्त्रात मात्र संबंधचल आणि फलनचल अमर्याद प्रमाणात
उपलब्ध असल्याने, अतार्किक चिन्हांचा वेगळा साठा नसतानाही
प्रथम-क्रम भाषेत सांगता येणारी कोणतीही गोष्ट सांगता येते.












\section{समरूपता}\label{sfr:fun:iso:sec}

\begin{explain}
\emph{समरूपता} म्हणजे संबंधित संचांची रचना जपणारे एकास-एक
आच्छादन. ही रचना संचांतील घटक मधील संबंधांवर अवलंबून
असते. पुढील दोन संच पाहा: $X=\{1,2,3\}$ आणि $Y=\{4,5,6\}$.
दोन्ही संचांवर उत्तरवर्ती, लहान असणे आणि मोठे असणे हे संबंध
आहेत. या दोन संचांमधील समरूपता म्हणजे या रचना जपणारे
एकास-एक आच्छादन. म्हणून $f \colon X
\to Y$ हे एकास-एक व आच्छादक फलन समरूपता असेल, जर
$i<j$ तेव्हा आणि केवळ तेव्हा $f(i)<f(j)$, $i>j$ तेव्हा आणि केवळ तेव्हा $f(i)>f(j)$, आणि $j$ हा $i$ चा उत्तरवर्ती असेल
तेव्हा आणि केवळ तेव्हा $f(j)$ हा $f(i)$ चा उत्तरवर्ती असेल.
\end{explain}

\begin{defn}[समरूपता]
$U$ ही जोडी $\langle X, R\rangle$ आणि $V$ ही जोडी $\langle
Y, S\rangle$ असो; येथे $X$ आणि $Y$ हे संच, तर $R$ आणि $S$ हे
अनुक्रमे $X$ आणि $Y$ वरील संबंध आहेत. $X$ पासून $Y$ कडे जाणारे
एकास-एक आच्छादन~$f$ हे $U$ पासून $V$ कडे जाणारी \emph{समरूपता}
असेल तेव्हा आणि केवळ तेव्हा ते संबंधात्मक रचना जपते; म्हणजे
$X$ मधील कोणत्याही $x_{1}$ आणि $x_{2}$ साठी
$\tuple{x_1,x_2} \in R$ तेव्हा आणि केवळ तेव्हा
$\tuple{f(x_1),f(x_2)} \in S$.
\end{defn}

\begin{ex}
$X=\{1,2,3\}$ आणि $Y=\{4,5,6\}$ हे दोन संच, तसेच
लहान असणे आणि मोठे असणे हे संबंध पाहा. $f(x) = 7-x$ असे
असलेले $f\colon X \to
Y$ हे फलन $\tuple{X,<}$ आणि $\tuple{Y,>}$ यांमधील
समरूपता आहे.
\end{ex}



\begin{quote}\small\textbf{स्रोतनोंद.} SOL6-A340: समरूपतेच्या व्याख्येतील द्विदिशा स्पष्ट केली; सर्व संबंध आणि उदाहरणे जपली.\end{quote}











\section{परिचय}\label{sfr:ind:int:sec}

\begin{explain}
विगमन ही तर्कशास्त्रातील गणितात वारंवार वापरली जाणारी
सिद्धतेची पद्धत आहे. नैसर्गिक संख्यांवर केलेले विगमन
तुम्हाला गणितातून परिचित असेल. प्रत्यक्षात अनेक प्रकारचे
गणिती व तार्किक वस्तुसंच विगमनाला अनुकूल असतात.

सर्वसाधारणपणे, विगमनाने सार्वत्रिक दावा सिद्ध करता येतो;
म्हणजे एखाद्या प्रकारच्या प्रत्येक वस्तूला एक गुणधर्म आहे हे
दाखवता येते. विशेषतः, “सार्वत्रिक परिचय” या सिद्धता-पद्धतीने
सिद्धता करणे फार गुंतागुंतीचे असेल तेव्हा विगमन उपयुक्त ठरते.
विशिष्ट रीतीने बनवलेल्या गणिती वस्तूंवरच विगमन चालते:
संचातील प्रत्येक घटक मूलभूत असेल किंवा मूलभूत घटकांपासून
बनलेला असेल, तर त्या संचावर विगमन करता येते. अधिक नेमकेपणाने:
\end{explain}

\begin{defn}
एखादा संच~$S$ हा फलन~$f$ च्या खाली \emph{बंद} असतो तेव्हा
आणि केवळ तेव्हा $x \in
S$ असताना $f(x) \in S$ असते.
\end{defn}

\begin{explain}
उदाहरणार्थ, नैसर्गिक संख्यांचा संच ($\mathbb{N}$) उत्तरवर्ती
फलन $f(x) = x+1$ च्या खाली बंद आहे: $x$ नैसर्गिक संख्या
असेल, तर $x+1$ सुद्धा नैसर्गिक संख्या असते. मात्र नैसर्गिक
संख्यांचा संच पूर्ववर्ती फलन $f(x) = x-1$ च्या खाली बंद
नाही; कारण $0$ नैसर्गिक संख्या आहे, पण $-1$ नाही.
\end{explain}

\begin{defn}
फलन~$f(x)$ च्या प्रांतातील कोणत्याही वस्तू~$a$ साठी $P(a)$
असल्यास $P(f(a))$ असेल (म्हणजे $a$ ला गुणधर्म~$P$ असेल,
तर $f(a)$ लाही असेल), तर गुणधर्म~$P$ फलन~$f$ च्या खाली
\emph{जपला जातो} असे म्हणतात.
\end{defn}

\begin{explain}
“सम संख्या असणे” हा गुणधर्म $f(x) = x+2$ या फलनाखाली
जपला जातो; कारण $x$ सम असेल, तर $x+2$ सुद्धा सम असतो.
मात्र हा गुणधर्म उत्तरवर्ती फलनाखाली जपला जात नाही;
कारण $x$ सम असेल, तर $x+1$ विषम असतो.

नैसर्गिक संख्यांवर विगमन चालते, कारण प्रत्येक नैसर्गिक
संख्येपर्यंत 0 पासून उत्तरवर्ती फलन सान्त वेळा लावून
पोहोचता येते. ज्या संचावर विगमन करता येते, त्या प्रत्येक
संचाचे हे वैशिष्ट्य असते.
\end{explain}

\begin{defn}[$\Nat$ वरील विगमन]
0 ला गुणधर्म~$P$ असेल आणि $P$ हा गुणधर्म उत्तरवर्ती
फलनाखाली जपला जात असेल, तर प्रत्येक नैसर्गिक संख्येला
गुणधर्म~$P$ असतो.
\end{defn}

\begin{ex} 0 पासून $n$ पर्यंतच्या नैसर्गिक
संख्यांची बेरीज $\frac{n(n+1)}{2}$ असते. \\

\textbf{विगमनाचा आधार.} 0 पासून 0 पर्यंतच्या नैसर्गिक
संख्यांची बेरीज $0=
\frac{0\cdot 1}{2}$ असते.\\

\textbf{विगमनाची पायरी.} गुणधर्म $k$ साठी खरा आहे असे
माना. तो $k+1$ साठीही खरा आहे हे दाखवू. म्हणजे $P(k)$
खरे आहे असे माना; अर्थात,
\[ 0 + 1 + \ldots + k = \frac{k(k+1)}{2} \]
$P(k+1)$ खरे आहे, म्हणजे
\[ 0 + 1 + \ldots + k + (k+1) = \frac{(k+1)(k+2)}{2} \]
हे $P(k)$ या गृहीतकाचा उपयोग करून दाखवू.

\begin{align*}
0 + 1 + \ldots + k &= \frac{k(k+1)}{2} \tag{गृहीतक} \\
0 + 1 + \ldots + k + (k+1) &= \frac{k(k+1)}{2} + (k+1)
\tag{दोन्ही बाजूंना $k+1$ मिळवा} \\
&= \frac{k(k+1) + 2(k+1)}{2} \\
&= \frac{k^2 + k + 2k +2}{2} \\
&=\frac{(k+1)(k+2)}{2}
\end{align*}
अशा रीतीने विगमनाची पायरी सिद्ध होते आणि दावा सिद्ध झाला.
\end{ex}

\begin{explain}
सूत्रे बाबतीत मूलभूत घटक अण्वीय सूत्रे असतात.
कमी गुंतागुंतीची सूत्रे संकारकांनी जोडून अधिक गुंतागुंतीची
सूत्रे मिळतात. प्रत्येक संकारकाला सूत्रे वरचे
असे एक फलन अनुरूप मानता येते की ते संबंधित संकारक वापरून
एका किंवा अधिक सूत्रांपासून नवे सूत्र देते. उदाहरणार्थ,
$\varphi$ आणि $\psi$ ही दोन सूत्रे जोडून संयुती
बनवणारे फलन $\mathcal E _\land (\varphi,
\psi) = \varphi \land \psi$ असे लिहिता येते. यांना
\emph{सूत्ररचना-फलने} म्हणू. सूत्रांचा संच या
फलनांच्या खाली बंद आहे: $\varphi$ आणि $\psi$ ही
सूत्रे असतील, तर $\lnot \varphi$,
$\varphi \land
\psi$ इत्यादीही तशीच असतात.
\end{explain}

\begin{defn}[सूत्रे वरील विगमन]
प्रत्येक अण्वीय सूत्राला गुणधर्म~$P$ असेल आणि
सूत्र-रचना-फलनांच्या खाली $P$ जपला जात असेल, तर प्रत्येक
सूत्राला गुणधर्म~$P$ असतो.
\end{defn}

\begin{explain}
या व्याख्येतून सूत्रे वरील विगमनात्मक सिद्धतेची
“कृती” सुचते. प्रत्येक सूत्राला गुणधर्म~$P$
आहे हे दाखवायचे असल्यास पुढील पायऱ्या पाळा:\\

\textbf{विगमनाचा आधार.} $\varphi$ हे अण्वीय
सूत्र असो. [\ldots] म्हणून $\varphi$ ला गुणधर्म~$P$ असतो.

\textbf{विगमनाची पायरी.} $\varphi$ आणि $\psi$ ही दोन्ही
गुणधर्म~$P$ असलेली सूत्रे असोत.

प्रकरण $\lnot$: [\ldots] म्हणून $\lnot \varphi$ ला गुणधर्म~$P$ असतो.

प्रकरण $\land$: [\ldots] म्हणून $\varphi \land \psi$ ला गुणधर्म~$P$ असतो.

प्रकरण $\lor$: [\ldots] म्हणून $\varphi \lor \psi$ ला गुणधर्म~$P$ असतो.

प्रकरण $\lif$: [\ldots] म्हणून $\varphi \lif \psi$ ला गुणधर्म~$P$ असतो.

प्रकरण $\liff$: [\ldots] म्हणून $\varphi \liff \psi$ ला गुणधर्म~$P$ असतो.

प्रकरण $\lforall[]$: [\ldots] म्हणून $\lforall[x] \varphi$ ला गुणधर्म~$P$ असतो.

प्रकरण $\lexists[]$: [\ldots] म्हणून $\lexists[x] \varphi$ ला गुणधर्म~$P$ असतो. \\

म्हणून प्रत्येक सूत्राला गुणधर्म~$P$ असतो.
\end{explain}



\begin{quote}\small\textbf{स्रोतनोंद.} SOL6-A341: सूत्ररचना नेहमी दोनच सूत्रांपासून होत नाही हे दुरुस्त केले; बेरीजच्या आधारप्रसंगाची मर्यादा आधीच्या विधानाशी जुळवली आणि परिचय-नियमाची संज्ञा सुसंगत केली.\end{quote}










\section{संचांबद्दलच्या सिद्धता}\label{sfr:set:prf:sec}

\begin{editorial}
हा विभाग \verb|content/method/proofs| ने अधिक्रमित केला आहे आणि
या प्रकरणातून पूर्वनिर्धारितपणे काढला आहे.
\end{editorial}

\begin{explain}
आतापर्यंत ओळख करून दिलेले संच आणि संकेतलेखन
संच निर्दिष्ट करण्यासाठी व त्यांच्यातील संबंध व्यक्त
करण्यासाठी सोयीची लघुरूपे देतात. अशा संबंधांबद्दलचे
दावेही अनेकदा सिद्ध करावे लागतात. गणिती सिद्धतांची
ओळख नसेल, तर हे तुमच्यासाठी नवीन असेल. म्हणून एक
सोपे उदाहरण टप्प्याटप्प्याने पाहू. कोणत्याही संच~$X$
आणि $Y$ साठी $X \cap
(X \cup Y) = X$ हे नेहमी खरे असते, हे सिद्ध करू.
अशी संच-समानता कशी सिद्ध करायची? संच केवळ त्यांच्या
घटक वरून ठरतात, म्हणजे त्यांत तेच घटक
असतील तेव्हाच ते समान असतात, हे आठवा. म्हणून येथे
(a) $X \cap (X \cup
Y)$ मधील प्रत्येक घटक हा $X$ चाही घटक
आहे आणि, उलट, (b) $X$ मधील प्रत्येक घटक
हा $X \cap (X \cup Y)$ चाही घटक आहे, हे
सिद्ध करायचे आहे. दुसऱ्या शब्दांत, (a) $X \cap (X \cup Y) \subseteq X$
आणि (b) $X \subseteq X \cap (X \cup Y)$ हे दोन्ही दाखवतो.

“$X \cap (X
\cup Y)$ मधील प्रत्येक घटक~$z$ हा $X$ चाही घटक
आहे” असा सर्वसाधारण दावा सिद्ध करताना प्रथम कोणताही
$z \in X \cap (X \cup Y)$ दिला आहे असे मानतो आणि या
गृहीतकापासून $z \in X$ सिद्ध करतो. ही पद्धत “सर्वसाधारण
सोपाधिक सिद्धता” म्हणून परिचित असेल. या सिद्धतेत
गृहीतक आणि निष्कर्ष यांतील व्याख्यांचाही उपयोग करावा
लागेल; येथे उदाहरणार्थ “$\cap$” आणि “$\cup$” यांच्या
व्याख्यांचा. म्हणून (a) साठीचा युक्तिवाद असा:


\begin{quote}
(a) प्रथम $X \cap (X \cup Y) \subseteq X$ दाखवायचे आहे;
म्हणजे $\subseteq$ च्या व्याख्येनुसार कोणत्याही~$z$ साठी
$z \in X \cap (X \cup Y)$ असेल, तर $z \in X$ हे दाखवायचे.
म्हणून $z \in X \cap (X \cup
Y)$ असे माना. $z$ हा दोन संचांच्या छेदाचा घटक
तेव्हा आणि केवळ तेव्हा असतो जेव्हा तो दोन्ही संचांचा
घटक असतो. म्हणून $z \in X$ आणि $z \in X \cup Y$
असे निष्पन्न होते. विशेषतः $z \in X$, हेच दाखवायचे होते.
\end{quote}

याने सिद्धतेचा पहिला भाग पूर्ण झाला. शेवटच्या
पायरीत आपण एक तथ्य वापरले: एखाद्या गृहीतकापासून
संयुती ($z \in X$ आणि $z \in X
\cup Y$) निष्पन्न झाली, तर तिचा प्रत्येक संयुतीघटक
त्याच गृहीतकापासून निष्पन्न होतो. हा नियम “संयुती-निरसन”
किंवा $\land$Elim म्हणून परिचित असेल. आता (b) सिद्ध करू:

\begin{quote}
(b) आता $X \subseteq X \cap (X \cup Y)$ सिद्ध करू; म्हणजे
$\subseteq$ च्या व्याख्येनुसार कोणत्याही~$z$ साठी
$z \in X$ असेल, तर $z \in X \cap
(X \cup Y)$ हे दाखवू. $z \in X$ असे माना. $z \in X
\cap (X \cup Y)$ दाखवण्यासाठी “$\cap$” च्या व्याख्येनुसार
(i) $z \in X$ आणि (ii) $z \in X \cup Y$ हे दाखवावे
लागेल. (i) हे आपले गृहीतकच आहे; म्हणून अधिक
सिद्ध करण्याची गरज नाही. (ii) साठी आठवा की $z$
हा संचांच्या संयोगाचा घटक तेव्हा आणि केवळ तेव्हा
असतो जेव्हा तो त्या संचांपैकी किमान एकाचा घटक
असतो. $z \in X$ आणि $X \cup Y$ हा $X$ व $Y$ यांचा
संयोग असल्याने येथे ही अट पूर्ण होते. म्हणून
$z \in X \cup Y$. आपण (i) $z \in X$ आणि
(ii) $z \in X \cup Y$ ही दोन्ही दाखवली; त्यामुळे
“$\cap$” च्या व्याख्येनुसार $z \in X \cap (X \cup Y)$.
\end{quote}

हा युक्तिवाद काहीसा लांबलचक होता; पण संच आणि
त्यांच्यातील संबंधांविषयी आपण कसा विचार करतो हे
तो दाखवतो. सहसा आपण इतक्या स्पष्टपणे सर्व काही
लिहीत नाही; विशेषतः प्रत्येक व्याख्या पुन्हा सांगत
नाही. अधिक प्रगत पाठ्यात याच निकालाची सिद्धता
खूप संक्षिप्त असेल. ती पुढीलप्रमाणे दिसू शकते.
\end{explain}

\begin{prop}[शोषण]
सर्व संच $X$, $Y$ साठी,
\[
X \cap (X \cup Y) = X
\]
\end{prop}

\begin{proof}
(a) $z \in X \cap (X \cup Y)$ असे माना. मग $z \in X$;
म्हणून $X \cap (X
  \cup Y) \subseteq X$.

(b) आता $z \in X$ असे माना. मग $z \in X \cup Y$ सुद्धा;
त्यामुळे $z \in X \cap (X \cup Y)$ सुद्धा. म्हणून
$X \subseteq X \cap (X \cup
  Y)$.
\end{proof}

\begin{prob}
$X \cup (X \cap Y) = X$ हे तपशीलवार सिद्ध करा.
मग त्याची संक्षिप्त सिद्धता द्या. (सूचना:
$X \cup (X \cap Y) \subseteq X$ या दिशेसाठी
प्रकरणांनुसार सिद्धता, म्हणजे $\lor$Elim, लागेल.)
\end{prob}



\begin{quote}\small\textbf{स्रोतनोंद.} SOL6-A342: छेद आणि संयोगातील सदस्यत्वासाठी द्विदिशेचे शब्दरूप स्पष्ट केले; दोन्ही सिद्धता आणि सराव जपले.\end{quote}

\section{अंतःप्रज्ञावादी कट-निर्मूलनातील पूरक प्रसंग}
$A \ident (B \lif C)$ असेल, तर
$\pi$ चा शेवट असा असतो:
\[
\AxiomC{}
\RightLabel{$\pi_1'$}
\Deduce$B, \Gamma \fCenter \Delta, C$
\RightLabel{\RightR{\lif}}
\UnaryInf$\Gamma \fCenter \Delta, B \lif C$
\LeftSubproofLabel{$\pi_1$}
\AxiomC{}
\RightLabel{$\pi_2'$}
\Deduce$B \lif C, \Pi \fCenter \Lambda, B$
\AxiomC{}
\RightLabel{$\pi_2''$}
\Deduce$C, \Pi \fCenter \Lambda$
\RightLabel{\LeftR{\lif}}
\BinaryInf$B \lif C, \Pi \fCenter \Lambda$
\RightSubproofLabel{$\pi_2$}
\RightLabel{\Cut}
\BinaryInf$\Gamma, \Pi \fCenter \Delta, \Lambda$
\DisplayProof
\]
\Cut ने संपणारी सिद्धता पुढीलप्रमाणे:
\[
\AxiomC{}
\RightLabel{$\pi_1'$}
\Deduce$B, \Gamma \fCenter \Delta, C$
\RightLabel{\RightR{\lif}}
\UnaryInf$\Gamma \fCenter \Delta, B \lif C$
\LeftSubproofLabel{$\pi_1$}
\AxiomC{}
\RightLabel{$\pi_2'$}
\Deduce$B \lif C, \Pi \fCenter \Lambda, B$
\RightLabel{\Cut}
\BinaryInf$\Gamma, \Pi \fCenter \Delta, \Lambda, B$
\DisplayProof
\]
तिची कटची उंची
$\pi$ च्या कट-उंचीपेक्षा कमी आहे.
विगमन गृहीतकाने
$\Gamma, \Pi \fCenter \Delta, \Lambda, B$
याची सिद्धता~$\pi_1''$ मिळते.
\Cut ने संपणारी दुसरी
सिद्धता अशी आहे:
\[
\AxiomC{}
\RightLabel{$\pi_1'$}
\Deduce$B, \Gamma \fCenter \Delta, C$
\AxiomC{}
\RightLabel{$\pi_2''$}
\Deduce$C, \Pi \fCenter \Lambda$
\RightLabel{\Cut}
\BinaryInf$B, \Gamma, \Pi \fCenter \Delta, \Lambda$
\DisplayProof
\]
तिची कट-कोटी आणि
कटची उंची या दोन्ही
$\pi$ च्या कट-कोटी आणि
कटच्या उंची पेक्षा कमी आहेत.
म्हणून विगमन गृहीतक आणि
दुर्बलीकरण वापरून
$B, \Gamma, \Pi, \Pi \fCenter \Delta, \Lambda,
\Lambda$ याची
~\Log{G3c} मधील सिद्धता~$\pi_2'''$
मिळते. आता पुढील
\Cut ने संपणाऱ्या
सिद्धता वर विगमन
गृहीतक लावा:
\[
\AxiomC{}
\RightLabel{$\pi_1''$}
\Deduce$\Gamma, \Pi \fCenter \Delta, \Lambda, B$
\AxiomC{}
\RightLabel{$\pi_2'''$}
\Deduce$B, \Gamma, \Pi, \Pi \fCenter \Delta, \Lambda,
\Lambda$
\RightLabel{\Cut}
\BinaryInf$\Gamma, \Gamma, \Pi, \Pi, \Pi \fCenter \Delta, \Delta, \Lambda, \Lambda, \Lambda$
\DisplayProof
\]
(हिची कट-कोटी
$< \cutr{\pi}$ आहे.)
त्यातून
$\Gamma,
\Gamma, \Pi, \Pi, \Pi \fCenter \Delta, \Delta, \Lambda, \Lambda, \Lambda$
याची सिद्धता मिळते.
मग \ref{pt:seq:inv:prop:G3c-cont-adm}
लावून
$\Gamma, \Pi \fCenter \Delta, \Lambda$
याची सिद्धता~$\pi'$ मिळते.

\section{पूरक rules-G1i नियमसारणी}








\begin{table}
\begin{tabular}{@{}c@{}c@{}}
\multicolumn{2}{@{}c@{}}{स्वयंसिद्धके}\\[2ex]
$A \Sequent A$ & $\lfalse \Sequent \quad$
\\[2ex]
\multicolumn{2}{@{}c@{}}{संरचनात्मक नियम}\\[2ex]
\Axiom$ \Gamma \fCenter \Delta $
\RightLabel{\LeftR{\Weakening}}
\UnaryInf$ A, \Gamma \fCenter \Delta$
\DisplayProof
&
\Axiom$ \Gamma \fCenter $
\RightLabel{\RightR{\Weakening}}
\UnaryInf$ \Gamma \fCenter A$
\DisplayProof
\\[4ex]
\Axiom$ A, A, \Gamma \fCenter \Delta $
\RightLabel{\LeftR{\Contraction}}
\UnaryInf$ A, \Gamma \fCenter \Delta$
\DisplayProof
&
\\[4ex]
\multicolumn{2}{@{}c@{}}{तार्किक नियम}\\[2ex]
\Axiom$ \Gamma \fCenter A $
\RightLabel{\LeftR{\lnot}}
\UnaryInf$ \lnot A, \Gamma \fCenter $
\DisplayProof
&
\Axiom$A, \Gamma \fCenter $
\RightLabel{\RightR{\lnot}}
\UnaryInf$ \Gamma \fCenter \lnot A $
\DisplayProof
\\[4ex]\begin{tabular}[t]{@{}l@{}}
\Axiom$ A, \Gamma \fCenter \Delta$
\RightLabel{\LeftR{\land}}
\UnaryInf$ A \land B, \Gamma \fCenter \Delta$
\DisplayProof
\\[3ex]
\Axiom$B, \Gamma \fCenter \Delta$
\RightLabel{\LeftR{\land}}
\UnaryInf$A \land B, \Gamma \fCenter \Delta$
\DisplayProof\\[3ex]
\end{tabular}
&
\Axiom$\Gamma \fCenter A$
\Axiom$ \Gamma \fCenter B$
\RightLabel{\RightR{\land}}
\BinaryInf$ \Gamma \fCenter A \land B $
\DisplayProof
\\[4ex]\Axiom$A, \Gamma \fCenter \Delta$
\Axiom$B, \Gamma \fCenter \Delta$
\RightLabel{\LeftR{\lor}}
\BinaryInf$A \lor B, \Gamma \fCenter \Delta$
\DisplayProof
&
\begin{tabular}[t]{@{}r@{}}
\Axiom$\Gamma \fCenter A$
\RightLabel{\RightR{\lor}}
\UnaryInf$ \Gamma \fCenter A \lor B$
\DisplayProof
\\[3ex]
\Axiom$ \Gamma \fCenter B$
\RightLabel{\RightR{\lor}}
\UnaryInf$ \Gamma \fCenter A \lor B$
\DisplayProof\\[3ex]
\end{tabular}
\\[4ex]
\Axiom$ \Gamma \fCenter A$
\Axiom$ B, \Gamma \fCenter \Delta$
\RightLabel{\LeftR{\lif}}
\BinaryInf$ A \lif B, \Gamma \fCenter \Delta$
\DisplayProof
&
\Axiom$ A, \Gamma \fCenter B$
\RightLabel{\RightR{\lif}}
\UnaryInf$ \Gamma \fCenter A \lif B $
\DisplayProof
\\[4ex]
\Axiom$ A(t), \Gamma \fCenter \Delta$
\RightLabel{\LeftR{\lforall}}
\UnaryInf$ \lforall[x][A(x)],\Gamma \fCenter \Delta$
\DisplayProof
&
\Axiom$ \Gamma \fCenter A(a) $
\RightLabel{\RightR{\lforall}}
\UnaryInf$ \Gamma \fCenter \lforall[x][A(x)]$
\DisplayProof
\\[4ex]
\Axiom$ A(a), \Gamma \fCenter \Delta $
\RightLabel{\LeftR{\lexists}}
\UnaryInf$ \lexists[x][A(x)], \Gamma \fCenter \Delta$
\DisplayProof
&
\Axiom$ \Gamma \fCenter A(t) $
\RightLabel{\RightR{\lexists}}
\UnaryInf$ \Gamma \fCenter \lexists[x][A(x)]$
\DisplayProof
\end{tabular}
\caption{अंतःप्रज्ञावादी तर्कशास्त्रासाठी \Log{G1i} चे नियम.
$\RightR{\forall}$ आणि $\LeftR{\exists}$ मध्ये स्थिरांक~$a$ हा
निष्कर्ष-क्रमवर्तीमध्ये येऊ नये. $\Gamma$ हा सूत्रांचा
बहुसंच असून $\Delta$ रिकामा असू शकतो किंवा त्यात एकच
सूत्र असू शकते. \RightR{\Contraction} नियम नाही. \Log{G1m}
म्हणजे~\RightR{\Weakening} आणि तक्त्यातील
$\lfalse \Sequent \quad$ हे असत्याचे स्वयंसिद्धक
वगळलेली \Log{G1i}.}
\label{pt:seq:supp:tab:G1i}
\end{table}



\section{पूरक rules-G3i नियमसारणी}








\begin{table}
\begin{tabular}{@{}c@{}c@{}}
\multicolumn{2}{@{}c@{}}{स्वयंसिद्धके}\\[2ex]
$C, \Gamma \Sequent C$ & $\lfalse, \Gamma \Sequent \Delta$
\\[2ex]
\multicolumn{2}{@{}c@{}}{तार्किक नियम}\\[2ex]
\Axiom$\lnot A, \Gamma \fCenter A $
\RightLabel{\LeftR{\lnot}}
\UnaryInf$\lnot A, \Gamma \fCenter \Delta$
\DisplayProof
&
\Axiom$A, \Gamma \fCenter$
\RightLabel{\RightR{\lnot}}
\UnaryInf$ \Gamma \fCenter \lnot A $
\DisplayProof
\\[4ex]
\Axiom$ A, B, \Gamma \fCenter \Delta$
\RightLabel{\LeftR{\land}}
\UnaryInf$ A \land B, \Gamma \fCenter \Delta$
\DisplayProof
&
\Axiom$\Gamma \fCenter A$
\Axiom$ \Gamma \fCenter B$
\RightLabel{\RightR{\land}}
\BinaryInf$ \Gamma \fCenter A \land B $
\DisplayProof
\\[4ex]\Axiom$A, \Gamma \fCenter \Delta$
\Axiom$B, \Gamma \fCenter \Delta$
\RightLabel{\LeftR{\lor}}
\BinaryInf$A \lor B, \Gamma \fCenter \Delta$
\DisplayProof
&
\begin{tabular}[t]{@{}r@{}}
\Axiom$\Gamma \fCenter A$
\RightLabel{\RightR{\lor}}
\UnaryInf$ \Gamma \fCenter A \lor B$
\DisplayProof
\\[3ex]
\Axiom$ \Gamma \fCenter B$
\RightLabel{\RightR{\lor}}
\UnaryInf$ \Gamma \fCenter A \lor B$
\DisplayProof\\[3ex]
\end{tabular}
\\[4ex]
\Axiom$A \lif B, \Gamma \fCenter A$
\Axiom$B, \Gamma \fCenter \Delta$
\RightLabel{\LeftR{\lif}}
\BinaryInf$A \lif B, \Gamma \fCenter \Delta$
\DisplayProof
&
\Axiom$A, \Gamma \fCenter B$
\RightLabel{\RightR{\lif}}
\UnaryInf$ \Gamma \fCenter A \lif B $
\DisplayProof
\\[4ex]
\Axiom$ A(t), \lforall[x][A(x)], \Gamma \fCenter \Delta$
\RightLabel{\LeftR{\lforall}}
\UnaryInf$ \lforall[x][A(x)],\Gamma \fCenter \Delta$
\DisplayProof
&
\Axiom$ \Gamma \fCenter A(a) $
\RightLabel{\RightR{\lforall}}
\UnaryInf$ \Gamma \fCenter \lforall[x][A(x)]$
\DisplayProof
\\[4ex]
\Axiom$ A(a), \Gamma \fCenter \Delta $
\RightLabel{\LeftR{\lexists}}
\UnaryInf$ \lexists[x][A(x)], \Gamma \fCenter \Delta$
\DisplayProof
&
\Axiom$ \Gamma \fCenter A(t) $
\RightLabel{\RightR{\lexists}}
\UnaryInf$ \Gamma \fCenter \lexists[x][A(x)]$
\DisplayProof
\end{tabular}
\caption{अंतःप्रज्ञावादी तर्कशास्त्रासाठी \Log{G3i} चे नियम.
$\RightR{\forall}$ आणि $\LeftR{\exists}$ मध्ये स्थिरांक~$a$ हा
निष्कर्ष-क्रमवर्तीमध्ये येऊ नये. $\Gamma$ हा सूत्रांचा
बहुसंच असून $\Delta$ रिकामा असू शकतो किंवा त्यात एकच
सूत्र असू शकते. स्वयंसिद्धकांत सूत्र~$C$ अण्वीय
असावे. \Log{G1m} म्हणजे \Log{G1i} मधून
\RightR{\Weakening} आणि $\lfalse \Sequent \quad$
हे स्वयंसिद्धक वगळून मिळणारी पद्धत.}
\label{pt:seq:supp:tab:G3i}
\end{table}



\section{पूरक rules-LK नियमसारणी}







\begin{table}
\begin{tabular}{@{}c@{}c@{}}
\multicolumn{2}{@{}c@{}}{स्वयंसिद्धके}\\[2ex]
$A \Sequent A$ & $\lfalse \Sequent \quad$
\\[2ex]
\multicolumn{2}{@{}c@{}}{संरचनात्मक नियम}\\[2ex]
\Axiom$ \Gamma \fCenter \Delta $
\RightLabel{\LeftR{\Weakening}}
\UnaryInf$ A, \Gamma \fCenter \Delta$
\DisplayProof
&
\Axiom$ \Gamma \fCenter \Delta$
\RightLabel{\RightR{\Weakening}}
\UnaryInf$ \Gamma \fCenter \Delta, A$
\DisplayProof
\\[4ex]
\Axiom$ A, A, \Gamma \fCenter \Delta $
\RightLabel{\LeftR{\Contraction}}
\UnaryInf$ A, \Gamma \fCenter \Delta$
\DisplayProof
&
\Axiom$ \Gamma \fCenter \Delta, A, A$
\RightLabel{\RightR{\Contraction}}
\UnaryInf$ \Gamma \fCenter \Delta, A$
\DisplayProof
\\[4ex]
\Axiom$ \Gamma, A, B, \Pi \fCenter \Delta $
\RightLabel{\LeftR{\Exchange}}
\UnaryInf$\Gamma, B, A, \Pi \fCenter \Delta$
\DisplayProof
&
\Axiom$ \Gamma \fCenter \Delta, A, B, \Lambda$
\RightLabel{\RightR{\Exchange}}
\UnaryInf$ \Gamma \fCenter \Delta, B, A, \Lambda$
\DisplayProof
\\[4ex]
\multicolumn{2}{@{}c@{}}{तार्किक नियम}\\[2ex]
\Axiom$ \Gamma \fCenter \Delta, A $
\RightLabel{\LeftR{\lnot}}
\UnaryInf$ \lnot A, \Gamma \fCenter \Delta$
\DisplayProof
&
\Axiom$A, \Gamma \fCenter \Delta$
\RightLabel{\RightR{\lnot}}
\UnaryInf$ \Gamma \fCenter \Delta, \lnot A $
\DisplayProof
\\[4ex]\begin{tabular}[t]{@{}l@{}}
\Axiom$ A, \Gamma \fCenter \Delta$
\RightLabel{\LeftR{\land}}
\UnaryInf$ A \land B, \Gamma \fCenter \Delta$
\DisplayProof
\\[3ex]
\Axiom$B, \Gamma \fCenter \Delta$
\RightLabel{\LeftR{\land}}
\UnaryInf$A \land B, \Gamma \fCenter \Delta$
\DisplayProof\\[3ex]
\end{tabular}
&
\Axiom$\Gamma \fCenter \Delta, A$
\Axiom$ \Gamma \fCenter \Delta, B$
\RightLabel{\RightR{\land}}
\BinaryInf$ \Gamma \fCenter \Delta, A \land B $
\DisplayProof
\\[4ex]\Axiom$A, \Gamma \fCenter \Delta$
\Axiom$B, \Gamma \fCenter \Delta$
\RightLabel{\LeftR{\lor}}
\BinaryInf$A \lor B, \Gamma \fCenter \Delta$
\DisplayProof
&
\begin{tabular}[t]{@{}r@{}}
\Axiom$\Gamma \fCenter \Delta, A$
\RightLabel{\RightR{\lor}}
\UnaryInf$ \Gamma \fCenter \Delta, A \lor B$
\DisplayProof
\\[3ex]
\Axiom$ \Gamma \fCenter \Delta, B$
\RightLabel{\RightR{\lor}}
\UnaryInf$ \Gamma \fCenter \Delta, A \lor B$
\DisplayProof\\[3ex]
\end{tabular}
\\[4ex]
\Axiom$ \Gamma \fCenter \Delta, A$
\Axiom$ B, \Gamma \fCenter \Delta$
\RightLabel{\LeftR{\lif}}
\BinaryInf$ A \lif B, \Gamma \fCenter \Delta$
\DisplayProof
&
\Axiom$ A, \Gamma \fCenter \Delta, B$
\RightLabel{\RightR{\lif}}
\UnaryInf$ \Gamma \fCenter \Delta, A \lif B $
\DisplayProof
\\[4ex]
\Axiom$ A(t), \Gamma \fCenter \Delta$
\RightLabel{\LeftR{\lforall}}
\UnaryInf$ \lforall[x][A(x)],\Gamma \fCenter \Delta$
\DisplayProof
&
\Axiom$ \Gamma \fCenter \Delta, A(a) $
\RightLabel{\RightR{\lforall}}
\UnaryInf$ \Gamma \fCenter \Delta, \lforall[x][A(x)]$
\DisplayProof
\\[4ex]
\Axiom$ A(a), \Gamma \fCenter \Delta $
\RightLabel{\LeftR{\lexists}}
\UnaryInf$ \lexists[x][A(x)], \Gamma \fCenter \Delta$
\DisplayProof
&
\Axiom$ \Gamma \fCenter \Delta, A(t) $
\RightLabel{\RightR{\lexists}}
\UnaryInf$ \Gamma \fCenter \Delta, \lexists[x][A(x)]$
\DisplayProof
\end{tabular}
\caption{\Log{LK} चे नियम. $\RightR{\forall}$ आणि
$\LeftR{\exists}$ मध्ये स्थिरांक~$a$ हा
निष्कर्ष-क्रमवर्तीमध्ये येऊ नये. $\Gamma$, $\Delta$, $\Pi$ आणि~$\Lambda$
या सूत्रांच्या क्रमिका आहेत.}
\label{pt:seq:supp:tab:LK}
\end{table}



\section{पूरक rules-mG3i नियमसारणी}








\begin{table}
\begin{tabular}{@{}c@{}c@{}}
\multicolumn{2}{@{}c@{}}{स्वयंसिद्धके}\\[2ex]
$C, \Gamma \Sequent \Delta, C$ & $\lfalse, \Gamma \Sequent \Delta$
\\[2ex]
\multicolumn{2}{@{}c@{}}{तार्किक नियम}\\[2ex]
\Axiom$\lnot A, \Gamma \fCenter \Delta, A $
\RightLabel{\LeftR{\lnot}}
\UnaryInf$\lnot A, \Gamma \fCenter \Delta$
\DisplayProof
&
\Axiom$A, \Gamma \fCenter$
\RightLabel{\RightR{\lnot}}
\UnaryInf$ \Gamma \fCenter \Delta, \lnot A $
\DisplayProof
\\[4ex]
\Axiom$ A, B, \Gamma \fCenter \Delta$
\RightLabel{\LeftR{\land}}
\UnaryInf$ A \land B, \Gamma \fCenter \Delta$
\DisplayProof
&
\Axiom$\Gamma \fCenter \Delta, A$
\Axiom$ \Gamma \fCenter \Delta, B$
\RightLabel{\RightR{\land}}
\BinaryInf$ \Gamma \fCenter \Delta, A \land B $
\DisplayProof
\\[4ex]\Axiom$A, \Gamma \fCenter \Delta$
\Axiom$B, \Gamma \fCenter \Delta$
\RightLabel{\LeftR{\lor}}
\BinaryInf$A \lor B, \Gamma \fCenter \Delta$
\DisplayProof
&
\Axiom$\Gamma \fCenter \Delta, A, B$
\RightLabel{\RightR{\lor}}
\UnaryInf$ \Gamma \fCenter \Delta, A \lor B$
\DisplayProof
\\[4ex]
\Axiom$A \lif B, \Gamma \fCenter \Delta, A$
\Axiom$ B, \Gamma \fCenter \Delta$
\RightLabel{\LeftR{\lif}}
\BinaryInf$A \lif B, \Gamma \fCenter \Delta$
\DisplayProof
&
\Axiom$ A, \Gamma \fCenter B$
\RightLabel{\RightR{\lif}}
\UnaryInf$ \Gamma \fCenter \Delta, A \lif B $
\DisplayProof
\\[4ex]
\Axiom$ A(t), \lforall[x][A(x)], \Gamma \fCenter \Delta$
\RightLabel{\LeftR{\lforall}}
\UnaryInf$ \lforall[x][A(x)],\Gamma \fCenter \Delta$
\DisplayProof
&
\Axiom$ \Gamma \fCenter A(a) $
\RightLabel{\RightR{\lforall}}
\UnaryInf$ \Gamma \fCenter \lforall[x][A(x)]$
\DisplayProof
\\[4ex]
\Axiom$ A(a), \Gamma \fCenter \Delta $
\RightLabel{\LeftR{\lexists}}
\UnaryInf$ \lexists[x][A(x)], \Gamma \fCenter \Delta$
\DisplayProof
&
\Axiom$ \Gamma \fCenter \Delta, \lexists[x][A(x)], A(t) $
\RightLabel{\RightR{\lexists}}
\UnaryInf$ \Gamma \fCenter \Delta, \lexists[x][A(x)]$
\DisplayProof
\end{tabular}
\caption{अंतःप्रज्ञावादी तर्कशास्त्रासाठी बहुनिष्कर्षी \Log{mG3i} चे नियम.
$\RightR{\forall}$ आणि $\LeftR{\exists}$ मध्ये स्थिरांक~$a$ हा
निष्कर्ष-क्रमवर्तीमध्ये येऊ नये. $\Gamma$ आणि $\Delta$ हे सूत्रांचे
बहुसंच आहेत. स्वयंसिद्धकांत सूत्र~$C$ अण्वीय असावे.}
\label{pt:seq:supp:tab:mG3i}
\end{table}



\section{मूळ संपूर्ण-सामग्री चालकाची संपादकीय नोंद}
खालील नोंद गोठवलेल्या मूळ संपूर्ण-सामग्री चालकाची आहे; प्रस्तुत ग्रंथाच्या संपादकीय स्थितीचा ती नवा दावा करत नाही.
\begin{editorial}
ही फाइल ओपन लॉजिक प्रकल्पातील सर्व सामग्री लोड करते. अशा
संपादकीय नोंदी दिसत असल्यास, ही सामग्री कशी सादर केली जाणार
याचा विचार न करता फाइल संकलित केली गेली आहे, हे त्यावरून
समजते. तुम्हाला हे वाचता येत असेल, तर या PDF वरून शिकवणे
किंवा अभ्यास करणे बहुधा \emph{योग्य नाही}.

अध्यापनासाठी किंवा स्व-अध्ययनासाठी अधिक योग्य मजकूर तयार
करता यावा यासाठी ओपन लॉजिक प्रकल्पात अनेक सोयी आहेत.
उदाहरणार्थ, नेहमीच्या मांडणीत सर्व तार्किक संयोजक आदिम
मानले जातात आणि पुराव्यांमध्ये प्रत्येक संयोजकाचे सर्व प्रकार
तपासले जातात. पण त्यांपैकी काही प्रकार सरावासाठी सोडणे अधिक
चांगले असते. ओपन लॉजिक प्रकल्पाचे कामही अजून सुरू आहे.
सहयोग आणि सुधारणा वाढाव्यात म्हणून मसुदा स्वरूपातील किंवा
सरावप्रश्न नसलेली सामग्रीसुद्धा येथे समाविष्ट आहे. अभ्यासक्रमासाठी
तयार केलेल्या PDF मध्ये असे विभाग वगळले जातात.

अध्यापन आणि अभ्यासासाठी अधिक योग्य PDF पाहण्यासाठी
\href{http://builds.openlogicproject.org/}{OLP संकेतस्थळावरील नमुना अभ्यासक्रम}
पहा. स्वतःचा मजकूर तयार करण्यासाठी
\href{https://github.com/OpenLogicProject/OpenLogic/tree/master/courses/sample}{नमुना चालक फाइल}
वापरून सुरुवात करता येईल; अधिक सजवलेल्या आणि प्रगत उदाहरणांसाठी
या प्रकल्पातून तयार झालेल्या पाठ्यपुस्तकांचे स्रोत पाहता येतील.
\end{editorial}
\section{मूळ विभाग-चालकांच्या संपादकीय नोंदी}\label{source-driver:notes}
खालील नोंदी मूळ विभाग आणि प्रकरण-चालकांमधील आहेत. त्या मूळ ग्रंथाच्या संकलनमार्गांविषयी सांगतात; प्रस्तुत आवृत्तीविषयी नवे संपादकीय दावे नाहीत.
\subsection{संचांचे आकारमान — मूळ संपादकीय नोंद}\label{source-driver:OLP-0027}
\begin{editorial}
या प्रकरणात प्रगणने, गणनीयता आणि अगणनीयता यांची चर्चा आहे.
काही विभागांच्या दोन आवृत्त्या आहेत: अधिक प्राथमिक आवृत्तीत प्रगणने
म्हणजे याद्या किंवा $\PosInt$ पासूनची आच्छादने मानली आहेत;
अधिक अमूर्त आवृत्तीत प्रगणनांची व्याख्या $\Nat$ बरोबरच्या
एकास-एक आच्छादनांनी केली आहे.
\end{editorial}
\begin{editorial}
पुढील \ref{sfr:siz:enm-alt:sec},
\ref{sfr:siz:nen-alt:sec} आणि \ref{sfr:siz:red-alt:sec}
हे अनुक्रमे \ref{sfr:siz:enm:sec},
\ref{sfr:siz:nen:sec} आणि \ref{sfr:siz:red:sec}
यांचे पर्यायी विभाग आहेत. टिम बटन यांनी त्यांच्या Open Set Theory
या ग्रंथात वापरण्यासाठी ते तयार केले. ते किंचित अधिक प्रगत आहेत आणि
संच सिद्धांताच्या संदर्भात अधिक योग्य अशी गणनीयतेची भिन्न व्याख्या
वापरतात: यादी देता येणे किंवा $\PosInt$ पासूनच्या आच्छादक फलनाची
व्याप्ती असणे याऐवजी $\Nat$ किंवा त्याच्या आरंभीच्या खंडाशी एकास-एक
आच्छादन असणे.
\end{editorial}
\subsection{प्रथम-क्रम तर्कशास्त्र — मूळ संपादकीय नोंद}\label{source-driver:OLP-0138}
\begin{editorial}
  या भागात प्रथम-क्रम तर्कशास्त्राची अधिउपपत्ती संपूर्णतेपर्यंत मांडली
  आहे. सध्या ती विधानीय तर्कशास्त्राच्या स्वतंत्र विवेचनावर अवलंबून
  नाही; प्रत्येक गोष्ट सिद्ध केली आहे. ``FOL'' हा टॅग false असल्यास
  स्रोत फायली संख्यापकांवरील सामग्री वगळतात (आणि
  ``संरचना'' ऐवजी ``सत्य-मूल्यांकन'', $\Struct{M}$ ऐवजी
  $\pAssign{v}$, इत्यादी प्रतिस्थापित करतात). प्रत्यक्षात, विधानीय
  तर्कशास्त्राच्या भागातील बहुतेक सामग्री म्हणजे ``FOL'' टॅग बंद
  ठेवलेली प्रथम-क्रम सामग्रीच आहे.

  विधानीय तर्कशास्त्राचा भाग समाविष्ट केल्यास त्यामुळे बरीच
  पुनरावृत्ती होते. मात्र, या भागाला विधानीय तर्कशास्त्रातील सामग्री
  विचारात घेता यावी (आणि आधीच समाविष्ट केलेली सामग्री वगळता यावी,
  तसेच विधानीय भागातील संबंधित ठिकाणांचे संदर्भ देऊन सिद्धता
  संक्षिप्त करता याव्यात), अशी सोय करण्याचे नियोजन आहे.
\end{editorial}
\subsection{प्रतिमान उपपत्ती — मूळ संपादकीय नोंद}\label{source-driver:OLP-0182}
\begin{editorial}
  प्रतिमान उपपत्तीवरील साहित्य अपूर्ण आणि प्रयोगात्मक आहे. सध्या ते केवळ
  आल्डो अँटोनेली यांच्या प्रतिमान उपपत्तीवरील टिपणांचे रूपांतर आहे; मात्र
  प्रथम-क्रम तर्कशास्त्राच्या भागात समाविष्ट असलेले विषय (उपपत्त्या, संपूर्णता,
  संहतता) त्यातून वगळले आहेत. पाठ्यपुस्तकासाठी उपयुक्त ठरण्याकरिता त्यात बरीच
  अधिक प्रस्तावना, प्रेरणा व स्पष्टीकरण तसेच सरावप्रश्नांची गरज आहे. अँडी अराना
  खास या भागावर काम करण्याचे नियोजन करीत आहेत (\gitissue{65}).
\end{editorial}
\subsection{संगणनक्षमता — मूळ संपादकीय नोंद}\label{source-driver:OLP-0208}
\begin{editorial}
  हा भाग जेरेमी अविगॅड यांच्या संगणनक्षमता सिद्धांतावरील टिपांवर आधारित
  आहे. अजूनपर्यंत फक्त पुनरावर्ती फलनांच्या प्रकरणात सराव आहेत; तसेच
  प्रेरणा, उदाहरणे, तपशील आणि सराव घालून सर्वच सामग्रीचा विस्तार करणे
  उपयुक्त ठरेल.
\end{editorial}
\subsection{मोडल टॅब्लो — मूळ संपादकीय नोंद}\label{source-driver:OLP-0460}
\begin{editorial}
  मोडल तर्कशास्त्रातील पूर्वचिन्हयुक्त टॅब्लोंवरील
  हे प्रकरण अजून मसुदा आहे. अधिक उदाहरणे,
  संपूर्णतेच्या सिद्धता, आणि बंद टॅब्लो शोधण्याचे
  प्रयत्न निष्फळ ठरल्यास त्यातून प्रतिप्रतिमाने
  कशी मिळवता येतात याची चर्चा आवश्यक आहे.
\end{editorial}
\subsection{मोडल क्रमवर्ती कलन — मूळ संपादकीय नोंद}\label{source-driver:OLP-0470}
\begin{editorial}
  मोडल तर्कशास्त्रातील क्रमवर्ती कलनांवरील हे प्रकरण
  अजून मसुदा आहे. अधिक उदाहरणे, निर्दोषतेचे आणि
  संपूर्णतेचे पुरावे आवश्यक आहेत.
\end{editorial}
\subsection{कालिक तर्कशास्त्रे — मूळ संपादकीय नोंद}\label{source-driver:OLP-0476}
\begin{editorial}
  या प्रकरणात कालिक तर्कशास्त्रांचा आढावा घेतला आहे.
\end{editorial}
\subsection{संच, संबंध, फलने — मूळ संपादकीय नोंद}\label{source-driver:OLP-0720}
\begin{editorial}
  या फाइलमध्ये अनौपचारिक संचसिद्धान्तावरील या भागाची
  पहिली चार प्रकरणे मूळ, संथ गतीने मांडलेल्या आवृत्तीत
  समाविष्ट आहेत. संपूर्ण भाग
  \verb|sets-functions-relations-complete.tex| मध्ये आहे.
  या भागातील सामग्री मूलभूत अनौपचारिक संचसिद्धान्ताची
  बऱ्यापैकी पूर्ण ओळख करून देते. विद्यार्थ्यांना ही
  पार्श्वभूमी आधीच आहे असे गृहीत धरता येत नसेल, तर
  अभ्यासक्रमाची सुरुवात या सामग्रीच्या उजळणीने,
  निदान संच, फलने आणि संबंध या भागाच्या उजळणीने,
  करणे बहुधा योग्य ठरेल. त्यामुळे पुढील कोणत्याही
  तार्किक विभागांसाठी आवश्यक असलेल्या गणिती
  संकेतलेखनाचे मूलभूत कौशल्य सर्व विद्यार्थ्यांना
  प्राप्त होईल.
\end{editorial}
\subsection{संचांचे आकारमान — मूळ संपादकीय नोंद}\label{source-driver:OLP-0722}
\begin{editorial}
या प्रकरणात प्रगणने, गणनीयता आणि अगणनीयता यांची चर्चा आहे.
काही विभागांच्या दोन आवृत्त्या आहेत: संचसिद्धान्तावर अधिक
तपशीलवार चर्चा होत असल्यास वापरण्यासाठी एक अधिक प्राथमिक
आणि एक अधिक प्रगत आवृत्ती. या फाइलमध्ये फक्त प्राथमिक
आवृत्त्या समाविष्ट आहेत. अधिक प्रगत आवृत्त्या
\verb|size-of-sets-complete| मध्ये आहेत.
\end{editorial}
\chapter*{स्रोतावरील संपादकीय नोंदी}
\markboth{स्रोतावरील संपादकीय नोंदी}{स्रोतावरील संपादकीय नोंदी}
\addcontentsline{toc}{chapter}{स्रोतावरील संपादकीय नोंदी}
\section*{भाषांतरित स्रोत-एककांतील दुरुस्ती-नोंदी}
खालील नोंदी भाषांतरित स्रोत-एककांतील स्पष्ट दुरुस्त्या आणि निवडी सांगतात. विभागाजवळ आधीच दिसणाऱ्या नोंदी येथे पुन्हा दिलेल्या नाहीत.
\begin{quote}\small\textbf{OLP-0054 — स्रोतनोंद.} MRINF-003 पात्रता: गोठवलेल्या इंग्रजीतील अंतःस्थ cardeq रचना A\ensuremath{\approx}B\ensuremath{\approx}C अशी वैध रीतीने विस्तारते; ती स्रोतदोष नाही. मराठीतील स्पष्ट B\ensuremath{\approx}C निष्कर्ष त्याच तुल्यबलता-साखळीतील, पुढील उपयोगासाठी केंद्रित केलेली सममूल्य मांडणी आहे.\end{quote}
\begin{quote}\small\textbf{OLP-0058 — स्रोतनोंद.} MRPL-001: गोठवलेल्या इंग्रजीतील पहिल्या संक्षेपात जुळणारा आरंभीचा कंस नसतानाही शेवटचा कंस आहे. मराठीत मानक नकार-विकल्पयोग सूत्र दिले असून इंग्रजी बाइट्स बदललेले नाहीत.\end{quote}
\begin{quote}\small\textbf{OLP-0119 — स्रोतनोंद.} OLAXD-001: गोठवलेल्या इंग्रजी सूत्रात सर्वांत बाहेरचा उजवा कंस सुटला आहे. मराठी सूत्रात तो कंस जोडला असून गोठवलेले इंग्रजी बाइट्स बदललेले नाहीत.\end{quote}
\begin{quote}\small\textbf{OLP-0120 — स्रोतनोंद.} OLAXD-002: गोठवलेल्या इंग्रजी सूत्रात सर्वांत बाहेरचा एक उजवा कंस सुटला आहे. मराठी सूत्रात तो कंस जोडला असून गोठवलेले इंग्रजी बाइट्स बदललेले नाहीत.\end{quote}
\begin{quote}\small\textbf{OLP-0120 — स्रोतनोंद.} OLAXD-003: गोठवलेल्या इंग्रजी निष्कर्षात आवश्यक A-सशर्त भाग सुटला आहे. मराठी निष्कर्षात तो पुनर्स्थापित केला असून गोठवलेले इंग्रजी बाइट्स बदललेले नाहीत.\end{quote}
\begin{quote}\small\textbf{OLP-0122 — स्रोतनोंद.} OLAXD-004: गोठवलेल्या इंग्रजी मजकुरात दुसरा संदर्भ चुकून पुन्हा ax:land1 असा आहे. दुसऱ्या संयुक्ताचे निगमन सिद्ध करण्यासाठी मराठीत तो ax:land2 असा दुरुस्त केला असून गोठवलेले इंग्रजी बाइट्स बदललेले नाहीत.\end{quote}
\begin{quote}\small\textbf{OLP-0122 — स्रोतनोंद.} OLAXD-005: गोठवलेल्या इंग्रजी मजकुरात पहिला संदर्भ ax:lnot1 असा आहे; पण दाखवलेली दोन्ही सूत्रे थेट ax:lnot2 ची उदाहरणे आहेत. मराठीत संदर्भ दुरुस्त केला असून गोठवलेले इंग्रजी बाइट्स बदललेले नाहीत.\end{quote}
\begin{quote}\small\textbf{OLP-0123 — स्रोतनोंद.} OLAXD-006: गोठवलेल्या इंग्रजी मजकुरात शेवटचे पाऊल पुन्हा निगमन प्रमेयाने मिळते असे म्हटले आहे. प्रत्यक्षात आधीच्या सशर्त निष्कर्षावर सत्य स्थिरांकाचे स्वयंसिद्धक आणि विध्यात्मक अनुमान लावल्यावर अपेक्षित निष्कर्ष मिळतो; मराठीत हे कारण स्पष्ट केले असून गोठवलेले इंग्रजी बाइट्स बदललेले नाहीत.\end{quote}
\begin{quote}\small\textbf{OLP-0124 — स्रोतनोंद.} OLAXD-007: गोठवलेल्या इंग्रजी मजकुरातील QR निष्कर्षाच्या संख्यापक-व्याप्तीत B(x) आधी आवश्यक उद्गारचिन्ह सुटले आहे. मराठीत ते {!}B(x) असे पुनर्स्थापित केले असून गोठवलेले इंग्रजी बाइट्स बदललेले नाहीत.\end{quote}
\begin{quote}\small\textbf{OLP-0124 — स्रोतनोंद.} OLAXD-008: गोठवलेल्या इंग्रजी मजकुरातील एका पूर्तता-विधानात B(c) आधी आवश्यक उद्गारचिन्ह सुटले आहे. मराठीत ते {!}B(c) असे पुनर्स्थापित केले असून गोठवलेले इंग्रजी बाइट्स बदललेले नाहीत.\end{quote}
\begin{quote}\small\textbf{OLP-0125 — स्रोतनोंद.} OLAXD-009: गोठवलेल्या इंग्रजी स्रोताच्या अध्याय-शीर्ष टिप्पणीमध्ये जुने axiomatic-proofs नाव आहे; मार्ग आणि olfileid यांच्याशी सुसंगत असे axiomatic-deduction नाव येथे वापरले आहे.\end{quote}
\begin{quote}\small\textbf{OLP-0130 — स्रोतनोंद.} OLCOM-001: गोठवलेल्या इंग्रजी स्रोताच्या व्याख्याविषयक टिपेत सार्वत्रिक सूत्रातील A\_n नंतर (x\_n) सुटले आहे. मराठीत ते पुनर्स्थापित केले असून गोठवलेले इंग्रजी बाइट्स बदललेले नाहीत.\end{quote}
\begin{quote}\small\textbf{OLP-0132 — स्रोतनोंद.} OLCOM-002: गोठवलेल्या इंग्रजी स्रोताच्या सर्ववाची विगमन-प्रसंगातील अंतिम सूत्रात B ऐवजी A आले आहे. मराठीत प्रसंगाशी जुळणारे B पुनर्स्थापित केले असून गोठवलेले इंग्रजी बाइट्स बदललेले नाहीत.\end{quote}
\begin{quote}\small\textbf{OLP-0132 — स्रोतनोंद.} OLCOM-003: गोठवलेल्या इंग्रजी स्रोताच्या सत्यता पूर्वप्रमेयातील संख्यकीकृत प्रसंगांत चार वेळा पदे म्हटले आहे; संदर्भित निष्पत्ती आणि संतृप्त-दृष्टांत गुणधर्म बंद पदांसाठी आहेत. मराठीत बंद पदे असे स्पष्ट केले असून गोठवलेले इंग्रजी बाइट्स बदललेले नाहीत.\end{quote}
\begin{quote}\small\textbf{OLP-0133 — स्रोतनोंद.} OLCOM-004: गोठवलेल्या इंग्रजी स्रोताच्या फलन-पदातील t\_\{i+1\} नंतर सलग दोन स्वल्पविराम आहेत. मराठीत एकच स्वल्पविराम ठेवला असून गोठवलेले इंग्रजी बाइट्स बदललेले नाहीत.\end{quote}
\begin{quote}\small\textbf{OLP-0133 — स्रोतनोंद.} OLCOM-005: गोठवलेल्या इंग्रजी स्रोताच्या प्रतिनिधी-स्वातंत्र्य उदाहरणात दुसऱ्या प्रतिनिधी t' विषयी बोलताना न-पूर्ति सूत्रात पुन्हा t आले आहे. मराठीत t' पुनर्स्थापित केले असून गोठवलेले इंग्रजी बाइट्स बदललेले नाहीत.\end{quote}
\begin{quote}\small\textbf{OLP-0136 — स्रोतनोंद.} OLCOM-006: गोठवलेल्या इंग्रजी स्रोताच्या notFOL शाखेत सत्यता पूर्वप्रमेयासाठी prop:ccs च्या जागी prop:fsat-ccs वापरण्याचा गंतव्य-संदर्भ FOL टॅगच्या आत गेल्यामुळे दिसत नाही. मराठीत समान संदर्भ notFOL शाखेलाही स्पष्टपणे लागू केला असून सर्व गोठवलेले संदर्भ जतन केले आहेत.\end{quote}
\begin{quote}\small\textbf{OLP-0140 — स्रोतनोंद.} OLFOL-001: गोठवलेल्या इंग्रजी मजकुरातील या तार्किक निष्पन्नता-सूत्रात आणि पुढील दोन पुनरुक्तींमध्ये सार्वत्रिक सूत्राचा बंद करणारा चौकोनी कंस चुकीच्या जागी आहे किंवा गहाळ आहे, तर निष्कर्षानंतर एक जादा कंस आहे. मराठी मजकुरात तिन्ही ठिकाणी समान, सुव्यवस्थित आधारविधाने आणि निष्कर्ष दिले आहेत.\end{quote}
\begin{quote}\small\textbf{OLP-0143 — स्रोतनोंद.} OLFOL-002: गोठवलेल्या इंग्रजी मजकुरात अनेक स्थाने असू शकणारी चिन्हे constants म्हटली आहेत. स्थिरचिन्हे एकच वस्तू निर्देशित करतात; अनेक स्थाने असू शकणारी चिन्हे predicates असल्यामुळे मराठी मजकुरात ते दुरुस्त केले आहे.\end{quote}
\begin{quote}\small\textbf{OLP-0143 — स्रोतनोंद.} OLFOL-003: याच उदाहरणात प्रांत \{0, 1, 2\} असा ठरवला आहे; गोठवलेल्या इंग्रजी मजकुरातील \{1, 2, 3\} ही यादी 0 वगळते आणि प्रांताबाहेरील 3 समाविष्ट करते. मराठी मजकुरात प्रांताशी जुळणारी \{0, 1, 2\} ही यादी वापरली आहे.\end{quote}
\begin{quote}\small\textbf{OLP-0146 — स्रोतनोंद.} OLFOL-004: गोठवलेल्या इंग्रजीतील पहिल्या सार्वत्रिक सूत्रात Atom चा दुसरा फलसाधक चुकीने सार्वत्रिक सूत्राबाहेर ठेवला आहे. पुढील सर्वसाधारण रूपाशी जुळण्यासाठी मराठीत P(v\_0) हा पूर्ण आण्विक सूत्रभाग संख्यापकाच्या व्याप्तीत ठेवला आहे.\end{quote}
\begin{quote}\small\textbf{OLP-0152 — स्रोतनोंद.} OLFOL-005: गोठवलेल्या इंग्रजीतील पहिल्या संक्षेपात जुळणारा आरंभीचा कंस नसतानाही शेवटचा कंस आहे. मराठीत जुळत नसलेला कंस वगळला असून इंग्रजी बाइट्स बदललेले नाहीत.\end{quote}
\begin{quote}\small\textbf{OLP-0156 — स्रोतनोंद.} OLFOL-006: गोठवलेल्या इंग्रजी व्याख्येत k-स्थानी फलनासाठी m\_0 ते m\_k असे k+1 निर्देशांक आणि तितकेच युक्तिवाद दिले आहेत. येथे k युक्तिवादांसाठी m\_1 ते m\_k ही सुसंगत मांडणी वापरली आहे.\end{quote}
\begin{quote}\small\textbf{OLP-0156 — स्रोतनोंद.} OLFOL-007: हा पुरावा सामान्य भाषा L साठी आहे; गोठवलेल्या इंग्रजीतील या परिच्छेदाच्या दोन ठिकाणी चुकून विधानीय भाषा L\_0 आली आहे. येथे प्रमेयाशी सुसंगत L वापरले आहे.\end{quote}
\begin{quote}\small\textbf{OLP-0156 — स्रोतनोंद.} OLFOL-008: गोठवलेल्या इंग्रजीत या एकाच रचना-बाबीत विन्यासात्मक अनन्यता-चिन्हाऐवजी पर्यायी सममूल्यता-चिन्ह आले आहे. व्याख्या आणि इतर सर्व बाबींशी सुसंगत विन्यासात्मक अनन्यता येथे वापरली आहे.\end{quote}
\begin{quote}\small\textbf{OLP-0161 — स्रोतनोंद.} OLFOL-010: गोठवलेल्या इंग्रजी मजकुरातील ``single-two place relation'' हे जोडचिन्ह चुकीचे आहे; व्याख्येला आवश्यक असलेला एकच द्विस्थानी संबंध मराठीत स्पष्ट केला आहे.\end{quote}
\begin{quote}\small\textbf{OLP-0162 — स्रोतनोंद.} OLFOL-011: गोठवलेल्या इंग्रजी प्रदर्शनात पहिल्या ओळीच्या शेवटी आणि पुढील ओळीच्या सुरुवातीला सलग दोन समानता-चिन्हे येतात. मराठी प्रदर्शनात पहिल्या ओळीच्या शेवटचे अनावश्यक चिन्ह वगळले आहे.\end{quote}
\begin{quote}\small\textbf{OLP-0163 — स्रोतनोंद.} OLFOL-013: गोठवलेल्या इंग्रजी मजकुराच्या अस्तित्ववाची शाखेत पूर्तिची आवश्यक अट वगळली आहे. समांतर वैश्विक शाखा आणि औपचारिक व्याख्या यांच्याशी जुळणारी अट मराठीत पुनःस्थापित केली आहे.\end{quote}
\begin{quote}\small\textbf{OLP-0163 — स्रोतनोंद.} OLFOL-012: गोठवलेल्या इंग्रजी मजकुरात संबंधाचे अर्थनिर्धारण चर-मूल्यांकनावर अवलंबून असल्यासारखा अनावश्यक निर्देश जोडला आहे. रचनेतील संबंध स्थिर असल्यामुळे तो निर्देश मराठीत काढला आहे.\end{quote}
\begin{quote}\small\textbf{OLP-0163 — स्रोतनोंद.} OLFOL-014: गोठवलेल्या इंग्रजी उदाहरणात नकारानंतर दोन ठिकाणी एक जादा उघडा कंस आणि एका सूत्रात जादा स्वल्पविराम आहे. समांतर सूत्रे व व्याख्या यांनुसार मराठी आवृत्तीत ते दुरुस्त केले आहेत.\end{quote}
\begin{quote}\small\textbf{OLP-0163 — स्रोतनोंद.} OLFOL-015: गोठवलेल्या इंग्रजी उदाहरणात m = 2, 3, 4 साठी उत्तरवर्ती R(a,x) असत्य म्हटले आहे; m = 2 साठी ते सत्य आहे. सशर्त सूत्राचा पूर्ववर्ती R(x,a) या तिन्ही मूल्यांसाठी असत्य असल्यामुळे मराठीत तोच वापरला आहे.\end{quote}
\begin{quote}\small\textbf{OLP-0163 — स्रोतनोंद.} OLFOL-016: गोठवलेल्या इंग्रजी मजकुरात दुसऱ्या रंजक प्रकरणातील चलाचे m हे नाव गाळले आहे. संदर्भ आणि लगेच पुढील गणना यांनुसार ते मराठीत पुनःस्थापित केले आहे.\end{quote}
\begin{quote}\small\textbf{OLP-0163 — स्रोतनोंद.} OLFOL-017: गोठवलेल्या इंग्रजी सारांशात बाह्य वैश्विक चलाला चुकून n म्हटले आहे. सर्व चार बाह्य मूल्यांचा संदर्भ m नेच असल्यामुळे मराठीत m वापरले आहे.\end{quote}
\begin{quote}\small\textbf{OLP-0163 — स्रोतनोंद.} OLFOL-018: गोठवलेल्या इंग्रजी मजकुरात m = 1 आणि m = 2 दोन्हीसाठी n = 4 हा प्रतिदृष्टांत दिला आहे; पण (2,4) हा R मध्ये आहे. मराठीत अनुक्रमे n = 4 आणि n = 1 हे योग्य प्रतिदृष्टांत दिले आहेत.\end{quote}
\begin{quote}\small\textbf{OLP-0164 — स्रोतनोंद.} OLFOL-019: गोठवलेल्या इंग्रजी सिद्धतेतील दुसरी k-सदस्यीय क्रमित सूची t\_i पासून सुरू होते आणि त्यामुळे i मुक्त राहतो. पहिल्या सूचीशी व घटकनिहाय समानतेशी जुळण्यासाठी मराठीत ती t\_1 पासून सुरू केली आहे.\}\end{quote}
\begin{quote}\small\textbf{OLP-0164 — स्रोतनोंद.} OLFOL-020: गोठवलेल्या इंग्रजी वैश्विक बाबतीत s\_1' आणि s\_2' ही दोन्ही अपरिभाषित s पासून बनवली आहेत. पुढील पर्याय व एकमत दाव्यांना आवश्यक असल्याप्रमाणे मराठीत ती अनुक्रमे s\_1 आणि s\_2 पासून बनवली आहेत.\}\}\end{quote}
\begin{quote}\small\textbf{OLP-0164 — स्रोतनोंद.} OLFOL-021: गोठवलेल्या इंग्रजी व्याख्येच्या सुरुवातीच्या नामपदबंधात Gamma सलग दोनदा आले आहे. अर्थ न बदलणारी दुसरी पुनरावृत्ती मराठीत काढली आहे.\end{quote}
\begin{quote}\small\textbf{OLP-0165 — स्रोतनोंद.} OLFOL-022: गोठवलेल्या इंग्रजी गणनेच्या पहिल्या ओळीच्या शेवटी आणि पुढील ओळीच्या सुरुवातीला सलग दोन समानताचिन्हे दिसतात. मराठी आवृत्तीत पहिल्या ओळीचे अनावश्यक चिन्ह काढले आहे.\end{quote}
\begin{quote}\small\textbf{OLP-0165 — स्रोतनोंद.} OLFOL-023: या प्रकरणातील सूत्र-आदेशनाची व्याख्या t' हे A मध्ये x साठी आदेशनासाठी मुक्त असतानाच लागू होते. गोठवलेल्या इंग्रजी विधानात ही आवश्यक पूर्वअट नसल्यामुळे मराठीत ती स्पष्टपणे जोडली आहे.\end{quote}
\begin{quote}\small\textbf{OLP-0171 — स्रोतनोंद.} OLFOL-024: गोठवलेल्या इंग्रजीतील दुसऱ्या सूत्राच्या शेवटच्या v\_2 पुढे वस्तुभाषेचे निर्देशचिन्ह चुकले आहे. समान परिच्छेदातील सर्व इतर चर-आस्थिती आणि अभिप्रेत सूत्राशी जुळण्यासाठी मराठीत वस्तुभाषेतील v\_2 केले असून गोठवलेले इंग्रजी बाइट्स बदललेले नाहीत.\end{quote}
\begin{quote}\small\textbf{OLP-0178 — स्रोतनोंद.} OLFOL-025: बाब (6) मध्ये x चा प्रकार tau दिलेला असूनही गोठवलेल्या इंग्रजीतील पुढील स्पष्टीकरणात तो sigma म्हटला आहे. फलनाचा प्रांत आणि tau ते sigma हा प्रकार जुळण्यासाठी मराठीत tau केले असून गोठवलेले इंग्रजी बाइट्स बदललेले नाहीत.\end{quote}
\begin{quote}\small\textbf{OLP-0187 — स्रोतनोंद.} OLFOL-026 — गोठवलेल्या इंग्रजी सूत्रात M-प्राइममधील पदमूल्य काढताना f चे अर्थनिर्धारण चुकून M मधील दिले आहे. व्याख्या आणि पुढील गणना यांनुसार येथे M-प्राइम असणे आवश्यक आहे; मराठी सूत्रात ती खूण दुरुस्त केली आहे.\end{quote}
\begin{quote}\small\textbf{OLP-0189 — स्रोतनोंद.} OLFOL-027 — गोठवलेल्या इंग्रजी व्याख्येत क्रमिकेची लांबी आणि पुनरावर्तनाचा प्रतांक या दोन्हींसाठी n वापरल्यामुळे I\_0 येथे चलरहित आण्विक सूत्रांपुरते चुकून मर्यादित होते. समान क्रमिकेची लांबी k अशी स्वतंत्र खूण घेऊन मराठी व्याख्या दुरुस्त केली आहे.\end{quote}
\begin{quote}\small\textbf{OLP-0190 — स्रोतनोंद.} OLFOL-028 — गोठवलेल्या इंग्रजी पुराव्यात p च्या प्रांतात a आधीच असण्याचे आणि p रिक्त असण्याचे प्रसंग वगळले आहेत; परंतु I मध्ये रिक्त नकाशा स्पष्टपणे समाविष्ट आहे. मराठी पुराव्यात सूत्रे न बदलता हे दोन प्रसंग जोडले आहेत.\end{quote}
\begin{quote}\small\textbf{OLP-0192 — स्रोतनोंद.} OLFOL-029 — गोठवलेल्या इंग्रजीतील या अस्तित्ववाचक पुराव्याच्या विधेयात साक्षीदार x हा पहिला युक्तिवाद चुकून वगळला आहे. त्याच परिच्छेदातील सुसंगतता-विधान आणि पुढील प्रत्येक मानक संख्यांकाचे सूत्र या दोन्हींनुसार मराठी सूत्रात x पुनःस्थापित केला आहे.\end{quote}
\begin{quote}\small\textbf{OLP-0193 — स्रोतनोंद.} OLFOL-030 — गोठवलेल्या इंग्रजीतील शेवटच्या तर्कात x हा s च्या प्रांतात नसतो असे म्हटले आहे; परंतु s चा प्रांत M असल्यामुळे x त्यात असतोच. आच्छादकता अपयशी होण्यासाठी x हा s च्या व्याप्तीत नसणे आवश्यक आहे, आणि मराठी मजकुरात ते दुरुस्त केले आहे.\end{quote}
\begin{quote}\small\textbf{OLP-0194 — स्रोतनोंद.} OLFOL-031 — गोठवलेल्या इंग्रजीतील पहिल्या सूत्रात c चे अर्थनिर्धारण M मध्ये घेतले आहे; पण c हे विस्तारित भाषेचे चिन्ह असल्यामुळे L\_A-रचना M मध्ये ते परिभाषितच नाही. मराठी सूत्रात ते योग्य विस्तारित रचना M\textasciicircum{}c मध्ये घेतले आहे.\end{quote}
\begin{quote}\small\textbf{OLP-0194 — स्रोतनोंद.} OLFOL-033 — गोठवलेल्या इंग्रजी पुराव्यात मनमाना सांत उपसंच घेतल्यानंतर लगेच सर्वांत मोठा k निवडला आहे; पण त्या उपसंचात c-विषयक असमता एकही नसेल, तर असा k नसतो. मराठी पुराव्यात तो रिकामा प्रसंग स्वतंत्रपणे हाताळला आहे.\end{quote}
\begin{quote}\small\textbf{OLP-0194 — स्रोतनोंद.} OLFOL-032 — विधानात गणनीय अमानक प्रतिमानाचा दावा आहे; गोठवलेल्या इंग्रजी पुराव्यात संहततेने प्रतिमान मिळवल्यानंतर त्याची गणनीयता सिद्ध केलेली नाही. उपलब्ध पूर्वप्रमेयानुसार अधोगामी लोव्हेनहाइम--स्कोलेम प्रमेय लावून ती पायरी मराठीत स्पष्ट केली आहे.\end{quote}
\begin{quote}\small\textbf{OLP-0195 — स्रोतनोंद.} OLFOL-034 — गोठवलेल्या इंग्रजीतील प्रकरण-विभागणीत K च्या प्रांताबाहेरील b लिहिले आहे; पण लगेचची गणना y = a हेच प्रकरण हाताळते. मराठी स्पष्टीकरणात b ऐवजी a केले आहे.\end{quote}
\begin{quote}\small\textbf{OLP-0195 — स्रोतनोंद.} OLFOL-035 — गोठवलेल्या इंग्रजीतील तिसऱ्या गणनेच्या शेवटी मुक्त y आहे; त्या ओळीची डावी बाजू आणि L ची बेरीज-सारणी दोन्ही येथे a आवश्यक ठरवतात. मराठी सूत्रात y ऐवजी a केले आहे.\end{quote}
\begin{quote}\small\textbf{OLP-0196 — स्रोतनोंद.} OLFOL-036 — गोठवलेल्या इंग्रजी विधानात प्रत्येक x ला पूर्ववर्ती असल्याचा दावा आहे, पण त्याच विधानानुसार शून्य हा लघुतम घटक असून Q\_2 नुसार तो कोणाचाही उत्तरवर्ती नाही. मराठीत पूर्ववर्तीचा दावा अशून्य x पुरता मर्यादित केला आहे; पुढील पुराव्यातील पुनरुक्तीही त्यानुसार दुरुस्त केली आहे.\end{quote}
\begin{quote}\small\textbf{OLP-0196 — स्रोतनोंद.} OLFOL-037 — गोठवलेल्या इंग्रजी पुराव्यात या एकाच गणनेत आधी न परिभाषित केलेले opulus चिन्ह वापरले आहे; विभागभर बेरीजेसाठी nsplus निश्चित केलेले असल्यामुळे मराठी सूत्रांत तेच चिन्ह पुनःस्थापित केले आहे.\end{quote}
\begin{quote}\small\textbf{OLP-0196 — स्रोतनोंद.} OLFOL-038 — गोठवलेल्या इंग्रजी निष्कर्षात कोणतेही M गणनीय आहे असे गृहीत न धरता अमानक खंड गणनीय असल्याचे म्हटले आहे; अमानक PA-प्रतिमाने अगणनीयही असू शकतात. पुढील वाक्याच्या गणनीय-प्रतिमान व्याप्तीशी जुळवून मराठी दावा गणनीय प्रतिमानापुरता स्पष्ट केला आहे.\end{quote}
\begin{quote}\small\textbf{OLP-0197 — स्रोतनोंद.} OLFOL-040 — गोठवलेल्या इंग्रजी संच-वर्णनात क्रमित जोडीतील x ऐवजी असंबद्ध n वर अट आहे. K च्या घोषित संबंधाशी आणि आधीच्या उदाहरणाशी जुळवून मराठी सूत्रात x बांधला आहे.\end{quote}
\begin{quote}\small\textbf{OLP-0197 — स्रोतनोंद.} OLFOL-039 — गोठवलेल्या इंग्रजी सूत्रात n>0 साठी g(n)=n+1 दिल्यामुळे 0 आणि 1 प्रतिमेत येत नाहीत आणि g द्विएकी राहत नाही. पुढील सर्व स्थानांतरित क्रियांशी सुसंगत असा अपेक्षित नकाशा g(n)=n-1 आहे.\end{quote}
\begin{quote}\small\textbf{OLP-0200 — स्रोतनोंद.} OLFOL-041 — गोठवलेल्या इंग्रजीतील या निष्पन्नतेच्या उजव्या बाजूस आधी परिभाषित केलेल्या H ऐवजी असंबद्ध ग्रीक डेल्टा छापले आहे. पुढील वाक्यातील प्रतिपरिवर्तनासाठी आणि संपूर्ण युक्तिवादासाठी येथे H चाच निषेध आवश्यक आहे; मराठी सूत्रात H पुनःस्थापित केला आहे.\end{quote}
\begin{quote}\small\textbf{OLP-0201 — स्रोतनोंद.} OLFOL-042 — गोठवलेल्या इंग्रजीतील या स्थानांतरण-सूत्रात h ला M'\_2 मधील P चे अर्थनिर्धारण दिले आहे; पण h चा प्रांत M\_1 आहे आणि पुढील घटकनिहाय अटही M'\_1 मधील P चे अर्थनिर्धारणच स्थानांतरित करते. मराठी सूत्रात M'\_1 पुनःस्थापित केले आहे.\end{quote}
\begin{quote}\small\textbf{OLP-0202 — स्रोतनोंद.} OLFOL-043 — गोठविलेल्या इंग्रजी विधानातील “if and only if” या पदबंधात “if” गाळले आहे; मराठी विधानात आशयाला आवश्यक असलेला द्विशर्त संबंध स्पष्ट केला आहे.\end{quote}
\begin{quote}\small\textbf{OLP-0202 — स्रोतनोंद.} OLFOL-044 — गोठविलेल्या इंग्रजी सूत्राच्या शेवटच्या अणुसूत्रात P'c\_1…c\_n अशी अपूर्ण कच्ची मांडणी आहे; येथे आधीपासून वापरलेले सुबद्ध अणुसूत्र पुनर्स्थापित केले आहे.\end{quote}
\begin{quote}\small\textbf{OLP-0205 — स्रोतनोंद.} OLFOL-045 — गोठवलेल्या इंग्रजीत M' ही L शी अनुरूप असल्याचे म्हटले आहे; पण पुनर्नामकरणात भाषा L शी नव्हे, तर रचना M शी अनुरूप रचना अपेक्षित आहे. मराठी वाक्यात M पुनःस्थापित केले आहे.\end{quote}
\begin{quote}\small\textbf{OLP-0206 — स्रोतनोंद.} OLFOL-046 — गोठवलेल्या इंग्रजीत या मोठ्या रचनेसाठी आधीच्या डाव्या उपरचनेचेच M चिन्ह पुन्हा वापरले आहे आणि N-तारांकित प्रांतातील तारका कंसाबाहेर चुकीच्या पातळीवर ठेवली आहे. मराठीत मोठ्या रचनेला A म्हटले असून N च्या प्रांताची तारका M च्या समांतर केली आहे.\end{quote}
\begin{quote}\small\textbf{OLP-0206 — स्रोतनोंद.} OLFOL-047 — सापेक्षीकरणाने मिळणारे D\_1 हे सर्वसाधारणपणे अमूर्त तर्कशास्त्रातील वाक्य असते, प्रथम-क्रम वाक्य नव्हे; D\_2 मात्र प्रथम-क्रम आहे. गोठवलेल्या इंग्रजीतील दोन्हींना first-order म्हणणारा दावा मराठीत दुरुस्त केला आहे.\end{quote}
\begin{quote}\small\textbf{OLP-0206 — स्रोतनोंद.} OLFOL-048 — गोठवलेल्या इंग्रजीत गणनीय मोठ्या प्रतिमानासाठी आणि त्यातील डाव्या उपरचनेसाठी M\_0 हेच चिन्ह दोनदा वापरले आहे. मराठीत मोठे प्रतिमान A\_0 आणि उपरचना M\_0, N\_0 अशी भिन्न चिन्हे वापरली आहेत.\end{quote}
\begin{quote}\small\textbf{OLP-0207 — स्रोतनोंद.} OLFOL-049 — गोठवलेल्या इंग्रजीतील प्रमेय-विधानात normal हे आवश्यक गृहीतक राहून गेले आहे; सिद्धता समरूपता, विस्तार आणि सापेक्षीकरण यांसह सामान्यतेचे गुणधर्म स्पष्टपणे वापरते. मराठी विधानात सामान्य हे गृहीतक पुनःस्थापित केले आहे.\end{quote}
\begin{quote}\small\textbf{OLP-0207 — स्रोतनोंद.} OLFOL-050 — गोठवलेल्या इंग्रजीत या पायरीतील उपअनुक्रमाला M\_n ऐवजी M चा उपअनुक्रम म्हटले आहे. मराठीत योग्य n-अधोलिखित पुनःस्थापित केले आहे.\end{quote}
\begin{quote}\small\textbf{OLP-0207 — स्रोतनोंद.} OLNOTE-001 — सर्व M\_n चा संयोग स्पष्टपणे घडवण्यासाठी स्थिरांच्या समान अर्थनिर्धारणांबाहेर परस्पर विभक्त प्रांत असलेल्या समरूपी प्रती निवडल्या आहेत. ही निवड पुरेशी आणि सोयीची आहे; संयोगासाठी ती अनिवार्य नाही.\end{quote}
\begin{quote}\small\textbf{OLP-0207 — स्रोतनोंद.} SOL6-A019 — स्थिरांचे समानतेचे आकृतिबंध भिन्न असतील, तर समरूपतेने सर्व रचनांतील स्थिरांना त्याच वस्तू नेमता येत नाहीत. प्रथम समान आण्विक सिद्धांताचा अनंत उपअनुक्रम निवडला आहे, त्यानंतर स्थिरांच्या अर्थनिर्धारणांचे एकीकरण केले आहे. वाढत्या मूळ निर्देशांकांचा नव्या निर्देशांकांपेक्षा लहान नसण्याचा गुणधर्म आवश्यक सममूल्यतेची खोली राखतो.\end{quote}
\begin{quote}\small\textbf{OLP-0207 — स्रोतनोंद.} OLFOL-051 — गोठवलेल्या इंग्रजीत या मोठ्या संकेतबद्ध रचनेसाठी पुन्हा M हेच चिन्ह वापरले आहे, जरी M आधीच सर्व M\_n च्या संयोगासाठी वापरलेले आहे. मराठीत मोठ्या रचनेसाठी A वापरून दोन वेगळ्या वस्तू स्पष्ट केल्या आहेत.\end{quote}
\begin{quote}\small\textbf{OLP-0207 — स्रोतनोंद.} OLFOL-052 — E हे अमूर्त तर्कशास्त्रातील वाक्य आहे; म्हणून M\_n आणि N\_n वरील दोन पूर्ति-चिन्हांना L अधोलिखित आवश्यक आहे. गोठवलेल्या इंग्रजीतील साधे models चिन्ह मराठीत models\_L असे दुरुस्त केले आहे.\end{quote}
\begin{quote}\small\textbf{OLP-0207 — स्रोतनोंद.} OLFOL-053 — गोठवलेल्या इंग्रजीतील संहततेचा नेहमीचा युक्तिवाद संक्षिप्त आहे. मराठी सिद्धतेत नवीन स्थिर प्रत्येक मानक घटकाच्या वर ठेवणारी वाक्ये स्पष्ट करून अमानक साक्षीचा निष्कर्ष उलगडला आहे.\end{quote}
\begin{quote}\small\textbf{OLP-0207 — स्रोतनोंद.} OLFOL-054 — गोठवलेल्या इंग्रजीत अंतिम मोठ्या प्रतिमानासाठी M\textasciicircum{}* या दोन विसंगत TeX रूपांचा वापर आहे आणि ते आधीच्या M\textasciicircum{}* वर्धित रचनेशीही धडकते. मराठीत संकेतबद्ध मोठ्या प्रतिमानाचे A\textasciicircum{}* हेच चिन्ह सातत्याने वापरले आहे.\end{quote}
\begin{quote}\small\textbf{OLP-0220 — स्रोतनोंद.} SOL6-A029 — रिकाम्या क्रमिकेसाठी 0 आणि 1 दोन्ही संकेतांक मान्य आहेत. पहिल्या concat पुनरावर्तनात दुसरी क्रमिका रिकामी असेल तर पहिला संकेतांक जतन होतो; परिबद्ध लघुतम शोध मात्र रिकाम्या परिणामासाठी 0 निवडतो. दोन पर्यायी संकेतांक-फलनांतील हा मर्यादित भेद येथे स्पष्ट केला आहे.\end{quote}
\begin{quote}\small\textbf{OLP-0226 — स्रोतनोंद.} OLCMP-012 — गोठवलेल्या इंग्रजीत काही e आंशिक पुनरावर्ती फलनाचा निर्देशांक नसू शकतो असे म्हटले आहे. पण आधीच दिलेल्या प्रमाण रूपातील T आणि U मुळे प्रत्येक नैसर्गिक e हा एका आंशिक पुनरावर्ती फलनाचा निर्देशांक आहे. मराठीत ही अपवादस्थिती काढून सिद्धतेतील तदनुरूप प्रकरणे दुरुस्त केली आहेत.\end{quote}
\begin{quote}\small\textbf{OLP-0232 — स्रोतनोंद.} OLCMP-013 — गोठवलेल्या इंग्रजी स्पष्टीकरणात कार्यक्रमाचा निर्देशांक दोन ठिकाणी x असा चुकून आला आहे; प्रमेयात आणि त्याच स्पष्टीकरणात तो e आहे. मराठीत दोन्ही ठिकाणी e ठेवला आहे. इंग्रजीतील ``It you think'' ही अक्षरदोषयुक्त सुरुवातही ``जर'' अशी वाचली आहे.\end{quote}
\begin{quote}\small\textbf{OLP-0233 — स्रोतनोंद.} OLCMP-014 — गोठवलेल्या इंग्रजीत शेवटचे g हे केवळ ``recursive'' म्हटले आहे. f आंशिक पुनरावर्ती असल्याने g सुद्धा आंशिक पुनरावर्तीच हमखास असते; ते सर्वत्र परिभाषित असणे आवश्यक नाही. मराठीत हा गुणधर्म स्पष्ट केला आहे.\end{quote}
\begin{quote}\small\textbf{OLP-0234 — स्रोतनोंद.} OLCMP-015 — गोठवलेल्या इंग्रजीतील पहिल्या वाक्यात ``total for the partial computable functions'' असे आले आहे. मागील प्रमेय आणि पुढील विरोध पाहता येथे अभिप्रेत शब्द ``universal'' आहे; मराठीत तोच अर्थ घेतला आहे.\end{quote}
\begin{quote}\small\textbf{OLP-0236 — स्रोतनोंद.} OLCMP-016 — गोठवलेल्या इंग्रजीत S \ensuremath{\in} S नंतर X \ensuremath{\notin} S असे आले आहे, पण X परिभाषित नाही. S च्या व्याख्येवरून आवश्यक विधान S \ensuremath{\notin} S आहे; मराठीत तोच निर्देशांक ठेवला आहे.\end{quote}
\begin{quote}\small\textbf{OLP-0239 — स्रोतनोंद.} OLCMP-017 — गोठवलेल्या इंग्रजीतील व्याप्तीच्या उलट दिशेच्या युक्तिवादात x अनिर्दिष्ट आहे. येथे साक्षीदार z चा पहिला घटक (z)\_0 हा योग्य आदान आहे; मराठीत तो स्पष्ट केला आहे.\end{quote}
\begin{quote}\small\textbf{OLP-0239 — स्रोतनोंद.} OLCMP-018 — गोठवलेल्या इंग्रजीत या दिशेची सिद्धता थेट S ही संगणनक्षम फलनाची व्याप्ती आहे असे धरते; पण व्याख्येनुसार रिकामा S देखील c.e. आहे आणि तो सर्वत्र परिभाषित फलनाची व्याप्ती असू शकत नाही. रिकाम्या संचासाठी कुठेही परिभाषित नसलेल्या फलनाचा स्वतंत्र आधार मराठीत दिला आहे.\end{quote}
\begin{quote}\small\textbf{OLP-0240 — स्रोतनोंद.} SOL6-A032 — गोठवलेल्या इंग्रजीतील len(z)=2 ही कसोटी अवैध prime-power संख्यांनाही लागू होऊ शकते; 42 मध्ये 2 आणि 3 आहेत, 5 नाही पण 7 आहे. परिभाषित length helper चा सर्वांत लहान साक्षीदार index 1 आहे, म्हणून length 2 येते. जोडीची आवश्यक-संपूर्ण कसोटी z=tuple((z)\_0,(z)\_1) ठेवून K\_0 चा अचूक domain मिळतो.\end{quote}
\begin{quote}\small\textbf{OLP-0241 — स्रोतनोंद.} OLCMP-019 — गोठवलेल्या इंग्रजीतील पहिल्या दोन सिद्धता दोन्ही c.e. संचांना सर्वत्र परिभाषित प्रगणक आहेत असे धरतात. c.e. च्या व्याख्येत रिकामा संचही आहे, पण तो अशा फलनाची व्याप्ती नसतो. रिकाम्या संचाची स्वतंत्र सोपी स्थिती येथे स्पष्ट केली आहे.\end{quote}
\begin{quote}\small\textbf{OLP-0242 — स्रोतनोंद.} OLCMP-020 — गोठवलेल्या इंग्रजी सिद्धतेत A चे परिभाषाक्षेत्र cfind\_d असे ठरवूनही अंतिम कसोटीत T(e,x,h(x)) आणि cfind\_e लिहिले आहे; ती कसोटी प्रत्यक्षात A च्या पूरकाची आहे. येथे d वापरले आहे. पुढील अनौपचारिक स्पष्टीकरणातील e,f ही जोडीही d,e अशी दुरुस्त केली आहे.\end{quote}
\begin{quote}\small\textbf{OLP-0243 — स्रोतनोंद.} OLCMP-021 — गोठवलेल्या इंग्रजीत K\_0 ची पहिली व्याख्या जोडी (e,x) वापरते, पण लगेचची समतुल्य व्याख्या (x,e) अशी उलट क्रमाने लिहिते. W\_e मध्ये x असल्याचा अर्थ पहिल्या व्याख्येतील (e,x) जोडीशी जुळतो; म्हणून येथे तीच क्रमरचना ठेवली आहे.\end{quote}
\begin{quote}\small\textbf{OLP-0244 — स्रोतनोंद.} OLCMP-022 — गोठवलेल्या इंग्रजी प्रश्नात अनेक-एक न्यूनीकरण f चे प्रकारलेखन f:A\ensuremath{\to}B असे आहे; मात्र आधीची व्याख्या f:N\ensuremath{\to}N अशी आहे आणि जातिबोधक फलांचे समीकरण संपूर्ण N वर तपासायचे आहे. म्हणून येथे आधीच्या व्याख्येशी सुसंगत प्रकारलेखन वापरले आहे.\end{quote}
\begin{quote}\small\textbf{OLP-0245 — स्रोतनोंद.} OLCMP-023 — गोठवलेल्या इंग्रजीतील शेवटचे सिद्धतावाक्य K चे K\_0 कडे न्यूनीकरण सांगते. त्यामुळे K\_0 च्या संपूर्णतेवरून K ची संपूर्णता सिद्ध होत नाही; आवश्यक दिशा K\_0 चे K कडे अशी आहे. पुढील स्वतंत्र प्रश्न मात्र मूळ K चे K\_0 कडे न्यूनीकरणच विचारतो, आणि तो तसाच ठेवला आहे.\end{quote}
\begin{quote}\small\textbf{OLP-0246 — स्रोतनोंद.} SOL6-A031 — अनेक-एक न्यूनीकरण संपूर्ण N वर अट पूर्ण करणारे total computable फलन असते. गोठवलेल्या इंग्रजीतील g(w) जोडीचे घटक वेगळे काढून थेट वापरते; मात्र प्रत्येक नैसर्गिक संख्या या prime-power coding मधील वैध जोडीचा संकेतांक नाही. अचूक पुनःसंकेतबद्धीकरणाची कसोटी आणि अवैध संकेतांकासाठी nowhere-defined फलनाचा निर्देशांक येथे जोडला आहे.\end{quote}
\begin{quote}\small\textbf{OLP-0250 — स्रोतनोंद.} OLCMP-024 — गोठवलेल्या इंग्रजी प्रमेयात आंशिक संगणनक्षम f नसल्याचे म्हटले आहे, पण सिद्धतेची पहिली ओळ फक्त सर्वत्र परिभाषित संगणनक्षम f घेते. पुढील g ची रचना f(e) अपरिभाषित असतानाही आंशिक संगणनक्षम राहते; म्हणून येथे आवश्यक सर्वसाधारण आंशिक f घेतले आहे.\end{quote}
\begin{quote}\small\textbf{OLP-0254 — स्रोतनोंद.} SOL6-A040 — गोठवलेल्या इंग्रजी परिचयात शीर्ष सर्वांत डावीकडील, म्हणजे end-marker च्या घरावर सुरू होते. पुढील औपचारिक initial configuration आणि सम यंत्राची सर्व उदाहरणे आदानाच्या पहिल्या घरावर सुरू होतात. मराठी परिचय त्या operative व्याख्येशी सुसंगत केला; गोठवलेले source bytes जतन आहेत.\end{quote}
\begin{quote}\small\textbf{OLP-0254 — स्रोतनोंद.} SOL6-A037 — गोठवलेल्या इंग्रजी इतिहासातील कोलॉससचे ट्यूरिंग यांना दिलेले बांधणी-श्रेय, SSEM ला सर्वसाधारण वापराचे म्हटलेला दावा आणि सर्व आधीची यंत्रे हाताने नव्याने जोडावी लागत हा दावा दुरुस्त केला आहे. Z3 साठी पालकांचा दिवाणखाना हे अनिश्चित ठिकाण वगळले; Zuse archive मध्ये तो तपशील Z1 शी जोडला आहे. प्रत्यक्ष पुन्हा पाहिलेल्या संस्थात्मक स्रोतांचे search extractions: \href{https://www.tnmoc.org/colossus}{राष्ट्रीय संगणन संग्रहालय}, \href{https://www.manchester.ac.uk/about/news/university-celebrates-babys-65th-birthday/}{मँचेस्टर विद्यापीठ} आणि \href{https://zuse.zib.de/z1}{झुसे संग्रह}. गणितीय व्याख्यांमध्ये बदल नाही.\end{quote}
\begin{quote}\small\textbf{OLP-0255 — स्रोतनोंद.} SOL6-A038 — गोठवलेल्या इंग्रजीत कोणत्याही यंत्राच्या कोणत्याही आदानासाठी स्थितिवर्णने पाहून प्रदान निश्चित करता येते असे म्हटले आहे. हा दावा थांबणाऱ्या संगणनांपुरता स्पष्ट केला; याच विभागातील विषम आदानावरचे सम यंत्र कायम चालते आणि त्याला अंतिम प्रदान नसते.\end{quote}
\begin{quote}\small\textbf{OLP-0255 — स्रोतनोंद.} OLTUR-001 — गोठवलेल्या इंग्रजीत आरंभीची अवस्था one म्हटली आहे; लगेच आधीचे चित्र आणि संक्रमण फलन शून्य अवस्था दाखवतात. मराठीत त्यांच्याशी सुसंगत शून्य अवस्था वापरली आहे.\end{quote}
\begin{quote}\small\textbf{OLP-0257 — स्रोतनोंद.} OLTUR-002 — या औपचारिक प्रारंभिक स्थितिवर्णनात शीर्ष आदानाच्या पहिल्या घरावर आहे; गोठवलेल्या इंग्रजी परिचयात ते फितीच्या डाव्या टोकाच्या घरावर असल्याचे म्हटले आहे. दोन्ही मूळ कथने गोठवलेल्या स्रोतांत जतन आहेत; मराठी परिचय आणि येथे दिलेली व्याख्या एकाच प्रारंभिक स्थानाशी सुसंगत केली आहेत.\end{quote}
\begin{quote}\small\textbf{OLP-0257 — स्रोतनोंद.} OLTUR-004 — गोठवलेली व्याख्या रिकाम्या आदानालाही परवानगी देते; पण दाखवलेल्या त्रयीत शीर्षाचा क्रमांक एक आणि फितीच्या क्रमिकेची लांबीही एक असल्याने आधीची काटेकोर लहानपणाची अट मोडते. मराठीत रिकामे घर समाविष्ट करण्याची मर्यादित स्पष्टता दिली आहे; सूत्र जतन आहे.\end{quote}
\begin{quote}\small\textbf{OLP-0257 — स्रोतनोंद.} OLTUR-003 — गोठवलेल्या इंग्रजीत या स्पष्टीकरणात आदान डाव्या टोकाच्या चिन्हाच्या डावीकडे सुरू होते असे शब्दशः येते; लगेच आधीचे सूत्र आणि परिच्छेद उजवी बाजू निश्चित करतात. मराठीत उजवीकडे असे स्पष्ट केले आहे.\end{quote}
\begin{quote}\small\textbf{OLP-0257 — स्रोतनोंद.} SOL6-A039 — सांत धावेतील अंतिम स्थितिवर्णनाला क्रमिकेत पुढचे स्थितिवर्णन नसते; अनंत धावेतील प्रत्येकाला असते. गोठवलेल्या व्याख्येतील प्रत्येक या शब्दाचा scope स्पष्ट केला. डावीकडील हालचालीच्या clause मधील अन्यथा फक्त डावीकडील हालचालीला लागू होते; उजवीकडील सूचनेला डाव्या टोकावर थांबवणारा पर्यायी successor मिळत नाही. रिकाम्या आदानासाठीचे आधीचे blank-padding स्पष्टीकरण कायम ठेवून मुख्य displayed initial formula फक्त nonempty input साठी म्हटले आहे.\end{quote}
\begin{quote}\small\textbf{OLP-0257 — स्रोतनोंद.} OLTUR-005 — आधीच्या विभागात डाव्या टोकाचे चिन्ह बदलून लिहिण्यास परवानगी आहे; गोठवलेल्या प्रदान-सूत्रात अंतिम फीत मात्र त्या चिन्हानेच सुरू होत असल्याचे गृहीत आहे. मराठीत या मर्यादेची स्पष्ट नोंद केली आहे.\end{quote}
\begin{quote}\small\textbf{OLP-0258 — स्रोतनोंद.} SOL6-A041 — गोठवलेल्या addition diagram मध्ये प्रारंभिक अवस्थेत रेघ वाचणारा loop दुसऱ्या अवस्थेकडे जातो. त्यामुळे मधला blank भरला जात नाही; उदाहरणार्थ तीन आणि दोन या आदानांचे output सलग पाच रेघा नसते. हा loop प्रारंभिक अवस्थेकडे परत नेला आहे; separator भरल्यावरच पुढच्या अवस्थेत जाते.\end{quote}
\begin{quote}\small\textbf{OLP-0258 — स्रोतनोंद.} SOL6-A041 — या mover exercise मधील गोठवलेला संदर्भ disciplined doubler कडे होता. आधीच्या परिच्छेदातील रचना मूळ सहा-अवस्थांच्या doubler ला mover जोडते; संदर्भ त्याच यंत्राकडे दुरुस्त केला.\end{quote}
\begin{quote}\small\textbf{OLP-0260 — स्रोतनोंद.} OLTUR-006 — गोठवलेल्या मजकुरात शीर्ष डाव्या टोकाचे चिन्ह सापडेपर्यंत डावीकडे नेण्याची पद्धत दिली आहे. पण आधीची यंत्र-व्याख्या ते चिन्ह पुसण्यास किंवा इतर घरांत लिहिण्यास मनाई करत नाही. त्यामुळे हा केवळ मर्यादित उदाहरणार्थ उपाय आहे; सर्वसाधारण शिस्तबद्ध रूपांतरासाठी अधिक अचूक चिन्ह-व्यवस्था आवश्यक आहे.\end{quote}
\begin{quote}\small\textbf{OLP-0260 — स्रोतनोंद.} SOL6-A042 — गोठवलेल्या disciplined addition diagram मध्येही प्रारंभिक stroke-loop चुकीने दुसऱ्या अवस्थेत जात होता. पहिला input block पार करताना प्रारंभिक अवस्थेतच राहणारा loop दिला; separator भरल्यावर पुढची अवस्था घेतली जाते. अंतिम head घर एकवर आणि designated h अवस्थेत पोहोचण्याची रचना कायम.\end{quote}
\begin{quote}\small\textbf{OLP-0261 — स्रोतनोंद.} OLTUR-007 — गोठवलेल्या इंग्रजीतील संक्रमण फलनाच्या पहिल्या शाखेत फक्त अवस्था Q मध्ये असणे ही अट आहे; त्यामुळे तिसऱ्या, अपरिभाषित-संक्रमण शाखेशी ती छेदते. येथे पहिली शाखा \ensuremath{\delta} परिभाषित असतानाच लागू होईल असे स्पष्ट केले आहे. सर्व चिन्हे आणि इतर शाखा जतन आहेत.\end{quote}
\begin{quote}\small\textbf{OLP-0261 — स्रोतनोंद.} SOL6-A043 — या विभागातील तिन्ही addition diagrams मध्ये पहिला stroke-loop चुकीने दुसऱ्या अवस्थेकडे होता; पहिला input block पार करताना प्रारंभिक अवस्थेत राहणारा loop केला. अंतिम combined machine रेघांचा एक दुप्पट सलग खंड ठेवते, पण त्याच्या डावीकडे blanks उरतात. म्हणून ते रेघांचा खंड दुप्पट करते एवढेच उदाहरण सांगते; unary numerical-output व्याख्येनुसार फलन संगणित करण्यासाठी आधीच्या mover प्रमाणे output फितीच्या आरंभी आणावे लागेल.\end{quote}
\begin{quote}\small\textbf{OLP-0262 — स्रोतनोंद.} SOL6-A044 — गोठवलेल्या boundary-simulation footnote मध्ये दुसऱ्या घरात end-marker लिहिण्याचे प्रकरण आहे; घर शून्यवरील marker पुसण्याला मात्र मूळ machine definition परवानगी देते. त्या स्थानाची ओळख जतन करणारी व्यवस्था लागते हे स्पष्ट केले. अनेक tapes ची संख्या प्रत्येक machine साठी सांत आहे; अनंत grid चे वेगळे उदाहरण कायम.\end{quote}
\begin{quote}\small\textbf{OLP-0265 — स्रोतनोंद.} SOL6-A045 — सांत संख्येच्या आदानांवरच परिभाषित असलेल्या फलनाचा दावा आंशिक संगणनक्षमतेचा आहे. त्या सांत सारणीतील मूल्ये दिली असता table lookup आणि बाहेरच्या inputs वर न थांबणे याने ते फलन संगणित करता येते; प्रत्येक नैसर्गिक आदानावर परिभाषित असल्याचा दावा नाही.\end{quote}
\begin{quote}\small\textbf{OLP-0267 — स्रोतनोंद.} SOL6-A048 — गोठवलेली यादी पहिल्या क्रमांकापासून लिहिलेली आहे, पण universal function चे निर्देशांक-परिभाषाक्षेत्र सर्व नैसर्गिक संख्या आहे आणि त्यात शून्य येते. येथे एकाच प्रभावी प्रगणनाला शून्यापासून निर्देशांक दिले; प्रत्येक नैसर्गिक निर्देशांकाला यंत्र मिळते. एकापासून निर्देशांक देणे स्वतः चुकीचे नाही, पण त्या पद्धतीत शून्याचा स्वतंत्र विस्तार सांगावा लागला असता.\end{quote}
\begin{quote}\small\textbf{OLP-0267 — स्रोतनोंद.} SOL6-A046 — गोठवलेल्या universal simulation sketch मध्ये zero-based tape-block numbering, रिकामे input padding, शून्य head वर left-stay आणि उजवीकडील tape extension स्पष्ट नव्हते. ते operative configuration नियमांशी जुळवले. Output cleanup मध्ये marker आणि trailing blanks वगळून शुद्ध unary output तपासणे आवश्यक आहे; प्रत्येक non-stroke block वर लगेच rejection केल्यास सर्व वैध initial markers आणि अनेक halting tapes चुकीने नाकारले जातात. अमान्य output असल्यास simulator थांबतो पण nonnumeric output ठेवतो, म्हणून halting equivalence आणि partial-function undefinedness दोन्ही जपली जातात. प्रत्यक्ष पूर्ण universal instruction table बांधल्याचा दावा नाही.\end{quote}
\begin{quote}\small\textbf{OLP-0268 — स्रोतनोंद.} SOL6-A049 — universal machine विभागात निश्चित केलेले तेच शून्यापासूनचे प्रभावी प्रगणन येथे वापरले आहे. गोठवलेल्या एकापासूनच्या यादीत शून्य निर्देशांकाचे यंत्र सांगितलेले नव्हते; या मांडणीत प्रत्येक नैसर्गिक निर्देशांकाला यंत्र मिळते.\end{quote}
\begin{quote}\small\textbf{OLP-0268 — स्रोतनोंद.} OLTUR-008 — गोठवलेल्या इंग्रजी पुराव्यात पहिले घर पाहत थांबणे म्हणजेच शिस्तबद्ध असणे असे कंसात म्हटले आहे. आधीच्या व्याख्येनुसार शिस्तबद्ध यंत्राला याशिवाय खास थांबण्याची अवस्था, टोकाचे चिन्ह राखणे आणि घर ० वरून डावीकडे न जाणे या अटीही आहेत. येथे आधीच्या रूपांतरानुसार पूर्ण शिस्तबद्ध यंत्र गृहीत धरले आहे; त्यातील पहिल्या घरावर थांबण्याचा गुणधर्म या जोडणीसाठी वापरला जातो.\end{quote}
\begin{quote}\small\textbf{OLP-0270 — स्रोतनोंद.} OLTUR-010 — गोठवलेल्या इंग्रजी मजकुरात A(x,y) बनवणाऱ्या अटींना sentences म्हटले आहे. प्रदर्शित अटींमध्ये x आणि y मुक्त आहेत; म्हणून त्या स्वतंत्र वाक्ये नसून उघडी सूत्रे आहेत. पुढील संक्रमण-स्वयंसिद्धकांत त्या सार्वत्रिक संख्यापकांच्या व्याप्तीत वापरल्या जातात.\end{quote}
\begin{quote}\small\textbf{OLP-0270 — स्रोतनोंद.} OLTUR-009 — गोठवलेल्या इंग्रजी डावीकडील संक्रमणाच्या पहिल्या फलितात A(x,y) आहे. पण त्या संक्रमणात चिन्ह लिहिले जाते x' या घरात; A च्या आधीच्या व्याख्येनुसार त्याने लिहिले जाणारे घर वगळले पाहिजे. येथे A(x',y) केले आहे. दुसऱ्या, घर शून्यवरील शाखेतील A(0,y) बदललेले नाही.\end{quote}
\begin{quote}\small\textbf{OLP-0271 — स्रोतनोंद.} OLTUR-013 — गोठवलेल्या इंग्रजी आढाव्यात एका निष्पन्नतेत T(M,w) ऐवजी आवश्यक {!}T(M,w) मधील {!} गाळलेले आहे; पुढील उदाहरणात संपूर्ण बांधणी M साठी असताना अचानक T यंत्र म्हटले आहे. मराठीत दोन्ही ठिकाणी सुसंगत {!}T(M,w) आणि M वापरले आहेत.\end{quote}
\begin{quote}\small\textbf{OLP-0271 — स्रोतनोंद.} SOL6-A051 — गोठवलेल्या इंग्रजी व्याख्येत आरंभीचे सर्वांत उजवे रिकामे नसलेले घर घेतले आहे. आदानाच्या शेवटी रिकामी चिन्हे असतील, तर ती मर्यादा आदानाच्या पूर्ण लांबीपेक्षा लहान होऊ शकते. दिलेली क्रम-स्वयंसिद्धके सर्व रचनांत लागोपाठच्या संख्यांकांमध्ये इतर घटक नसल्याची खात्री देत नाहीत; त्यामुळे लहान मर्यादेपलीकडील सार्वत्रिक रिकामेपण आरंभीच्या अटींवरून मिळत नाही. संपूर्ण आदानाची लांबी घेऊन मूळ विगमनाचा तर्क राखला आहे.\end{quote}
\begin{quote}\small\textbf{OLP-0271 — स्रोतनोंद.} OLTUR-012 — गोठवलेल्या इंग्रजी पुराव्यात आधी थांबणारी जोडी (q,\ensuremath{\sigma}) निवडल्यानंतर या निष्पन्नतेत अचानक (q',\ensuremath{\sigma}') येते, आणि एका S-विधेयापूर्वी Obj चिन्हही गाळले आहे. येथे आधी निवडलेलीच (q,\ensuremath{\sigma}) जोडी व भाषेतील तेच S-विधेय वापरले आहे; अस्तित्वविधानाचा निष्कर्ष बदलत नाही.\end{quote}
\begin{quote}\small\textbf{OLP-0271 — स्रोतनोंद.} OLTUR-011 — गोठवलेल्या इंग्रजी लेम्मात यंत्र n पायऱ्यांनंतर थांबलेले नसावे अशी अट आहे, पण पुढील पुरावा n व्या पायरीवर मिळणाऱ्या थांबणाऱ्या स्थितिवर्णनासाठीही तेच लेम्मा वापरतो. येथे चालक्रम n-पायऱ्यांच्या स्थितिवर्णनापर्यंत पोहोचणे ही अचूक अट घेतली आहे; ते स्थितिवर्णन थांबणारे असू शकते.\end{quote}
\begin{quote}\small\textbf{OLP-0271 — स्रोतनोंद.} OLTUR-014 — गोठवलेल्या इंग्रजी पडताळणीत डावीकडील संक्रमणाच्या पहिल्या शाखेत A(x,y) आणि त्याच्या उदाहरणात A(l,n) पुन्हा आले आहे. OLTUR-009 नुसार लिहिलेले घर x'=l'=m वगळले पाहिजे. म्हणून येथे A(x',y) आणि A(l',n) वापरले आहेत; घर शून्यवरील शाखा बदललेली नाही.\end{quote}
\begin{quote}\small\textbf{OLP-0271 — स्रोतनोंद.} OLTUR-015 — गोठवलेल्या इंग्रजी उजवीकडील व डावीकडील पुढील-स्थिती प्रदर्शनांत प्रत्येक घरासाठी जुने \ensuremath{\sigma}\_i चिन्हच पुन्हा दाखवले आहे, जरी आधीचे वाक्य लिहिलेल्या घरात \ensuremath{\sigma}' येते असे म्हणते आणि पुढील स्पष्टीकरण जुने त्या घराचे विधान वगळते. येथे वरची + खूण पुढील-स्थिती चिन्हांसाठी आहे: लिहिलेल्या घरात नवे चिन्ह आणि इतरत्र जुने चिन्ह. त्यामुळे प्रदर्शनाचे वाचन स्थितिवर्णनाच्या व्याख्येशी जुळते.\end{quote}
\begin{quote}\small\textbf{OLP-0272 — स्रोतनोंद.} SOL6-A052 — गोठवलेल्या इंग्रजी बांधणीत E सर्वांत डाव्या फीतघरावर थांबते म्हटले आहे, पण ते घर अंतखूणेचे आहे; पुढे चालवायच्या D चे प्रारंभिक शीर्षस्थान पहिले आदानघर आहे. येथे E ने तयार केलेले वाक्य त्या घरापासून ठेवून त्याच घरावर थांबण्याची अट घेतली आहे.\end{quote}
\begin{quote}\small\textbf{OLP-0272 — स्रोतनोंद.} OLTUR-016 — गोठवलेल्या इंग्रजी पुराव्यात आदान-वाक्य प्रथम B असे लिहिले आहे, पण त्याच परिच्छेदात पुढे {!}B सरकवले जाते आणि संपूर्ण प्रमेयात {!}B हेच वाक्य आहे. येथे प्रारंभीही {!}B लिहिले आहे; नकार आणि पूर्ततायोग्यतेची युक्ती अपरिवर्तित आहे.\end{quote}
\begin{quote}\small\textbf{OLP-0273 — स्रोतनोंद.} OLTUR-017 — गोठवलेल्या इंग्रजी सांत प्रतिरूपात वरची मर्यादा थांबण्याच्या पायऱ्या k आणि आदानाची लांबी यांतील मोठी संख्या घेतली आहे. यंत्र घर एकपासून सुरू होऊन k उजव्या हालचालींनंतर k+1 या घरात पोहोचू शकते. तसेच सर्वोच्च घटकावर n<n ठेवलेल्या सांत क्रमामुळे n हे शेवटचे आदान-चिन्ह असलेले घर नसावे; अन्यथा रिकाम्या-शेपटीची अट त्या घरालाच रिकामे म्हणेल. म्हणून मराठीत दोन्ही संख्यांपेक्षा एकाने मोठी मर्यादा घेतली आहे.\end{quote}
\begin{quote}\small\textbf{OLP-0273 — स्रोतनोंद.} OLTUR-018 — गोठवलेल्या इंग्रजी सरावातील एकमेव सूचना पुढील अवस्थेला अव्याख्यायित q म्हणते, जरी एकच अवस्था q0 दिली आहे; तसेच M'' च्या उत्तरवर्ती फलनात एका ठिकाणी M' लिहिले आहे. मराठीत दोन्ही ठिकाणी दिलेलीच एक-अवस्थेची यंत्रणा आणि M'' रचना वापरली आहे.\end{quote}
\begin{quote}\small\textbf{OLP-0273 — स्रोतनोंद.} OLTUR-019 — गोठवलेल्या इंग्रजी सांत-प्रतिरूप बांधणीत डावीकडे जाण्याच्या सकारात्मक-स्थान शाखेत स्थिरता-अट लिहिलेल्या घराऐवजी नव्या शीर्षस्थानाला वगळते; त्याच शाखेत नव्या वेळेसाठी B अटही गाळली आहे, जरी लगतचा मजकूर प्रत्येक संक्रमणात ती जोडल्याचे सांगतो. येथे आधीच्या OLTUR-009 नुसार लिहिलेले घर वगळले आहे आणि या शाखेतही नव्या वेळेची अट घातली आहे. शून्यावरील शाखा बदललेली नाही.\end{quote}
\begin{quote}\small\textbf{OLP-0273 — स्रोतनोंद.} OLTUR-020 — गोठवलेल्या इंग्रजी लेम्माच्या पुराव्यात आणि लगेचच्या सरावात थांबण्याचे वाक्य E असे लिहिले आहे; सरावात T' ऐवजी जुने T आले आहे. येथे दोन्ही ठिकाणी लेम्मातील संपूर्ण T' आणि {!}E संयुक्त विधानच तपासायचे आहे.\end{quote}
\begin{quote}\small\textbf{OLP-0273 — स्रोतनोंद.} OLTUR-021 — गोठवलेल्या इंग्रजी प्रतिपरिवर्तनात गृहीत प्रतिरूप अचानक जुन्या T आणि {!}E संयुक्त विधानाचे म्हटले आहे, जरी लेम्मा आणि पुढील निष्पन्नता T' वर अवलंबून आहेत. येथे लेम्मातीलच T' आणि {!}E गृहीत घेतले आहेत.\end{quote}
\begin{quote}\small\textbf{OLP-0273 — स्रोतनोंद.} OLTUR-022 — गोठवलेल्या इंग्रजी मजकुरात आरंभीच्या वेळेसाठीही B(0) निष्पन्न असल्याचा दावा आहे. पण B(y') ही अट संक्रमणांच्या निष्कर्षातच आहे; आरंभीच्या वेळेसाठी ती दिलेली नाही. येथे B(n) फक्त किमान एक पायरी झाल्यावर वापरले आहे. परस्परभिन्न संख्यांक मिळवण्यासाठी हेच पुरेसे आहे.\end{quote}
\begin{quote}\small\textbf{OLP-0273 — स्रोतनोंद.} SOL6-A053 — गोठवलेल्या इंग्रजी उपप्रमेयातील निष्पत्तीपद्धतीसाठी पुढील पुरावा निष्पत्तींचे प्रभावी प्रगणन वापरतो. ती आवश्यक अट येथे स्पष्ट केली आहे; केवळ सांत-वैध वाक्यांचा अप्रभावी स्वयंसिद्धकसंच घेतल्याने विरोधाभास मिळत नाही. सर्व सांत रचनांतील सत्यता हा सममूल्यतेचा पूर्ण उजवा भाग आहे.\end{quote}
\begin{quote}\small\textbf{OLP-0276 — स्रोतनोंद.} SOL6-A054 — गोठवलेल्या इंग्रजी परिच्छेदात फ्रेगेचे logicist उद्दिष्ट सर्व गणितासाठी सांगितले आहे. त्याच्या 1884 च्या Grundlagen मध्ये अंकगणित विश्लेषक, तर भूमिती अनुभवपूर्व संश्लेषक मानली आहे; येथे हा फरक जपला.\end{quote}
\begin{quote}\small\textbf{OLP-0276 — स्रोतनोंद.} SOL6-A054 — 16 जून 1902 चे रसेलचे फ्रेगेला पत्र McMaster Russell Archives मध्ये नोंदलेले आहे. येथे 1902 हा विरोधाभास कळवण्याचा वर्षसंदर्भ आहे; शोध लागल्याची तारीख म्हणून तो दिलेला नाही.\end{quote}
\begin{quote}\small\textbf{OLP-0276 — स्रोतनोंद.} SOL6-A054 — मूळ ऐतिहासिक सारांशातील सर्व प्रणालींबद्दलचा दावा येथे आवश्यक प्रभावीपणा व अंकगणितीय सामर्थ्य यांच्या अटींसह दिला आहे. केवळ सुसंगततेवरून स्वतःच्या सुसंगतता-विधानाच्या खंडनाची अशक्यता मिळत नाही.\end{quote}
\begin{quote}\small\textbf{OLP-0277 — स्रोतनोंद.} SOL6-A055 — मूळ परिच्छेदातील प्रतिनिधित्वाची अट पहिल्या अपूर्णता प्रमेयासाठी पुरेशी सामर्थ्याची अट आहे. दुसऱ्या प्रमेयासाठी अतिरिक्त अटी पुढील विभागात दिल्या आहेत. सुसंगतता व प्रभावी स्वयंसिद्धकीय मांडणी या इतर आवश्यक अटीही निष्कर्षात स्पष्ट केल्या.\end{quote}
\begin{quote}\small\textbf{OLP-0278 — स्रोतनोंद.} SOL6-A056 — गोठवलेल्या इंग्रजी आढाव्यात मूळ Gödel sentence आणि केवळ consistency गृहीत धरणाऱ्या व्यापक प्रमेयाची गल्लत होते. मूळ वाक्याच्या स्वतंत्रतेसाठी अतिरिक्त गृहीतक आणि Rosser च्या बदललेल्या रचनेचा फरक येथे स्पष्ट केला.\end{quote}
\begin{quote}\small\textbf{OLP-0283 — स्रोतनोंद.} SOL6-A063 — मर्यादा सिद्ध करताना गरज नसलेली पदे वगळून केवळ अंतिम पदाच्या उपपदांची रचनाक्रमिका निवडली आहे. मनमानी रचनाक्रमिकेत इतर मोठी पदे असू शकतात; त्यांच्यासाठी ही लांबीची मर्यादा दावा केलेली नाही.\end{quote}
\begin{quote}\small\textbf{OLP-0284 — स्रोतनोंद.} SOL6-A064 — मूळ आण्विक-सूत्र चाचणीतील argument-code tuple साठी z<x ही मर्यादा चुकीची आहे. एका वैध द्विस्थानी विधेयाला दोन स्थिरपदे दिल्यास हा tuple code अंतिम formula code पेक्षा मोठा असू शकतो. योग्य आदिम पुनरावर्ती मर्यादा आणि tuple length ही predicate arity इतकीच असण्याची अट येथे स्पष्ट केली.\end{quote}
\begin{quote}\small\textbf{OLP-0284 — स्रोतनोंद.} SOL6-A064 — मूळ sentence चाचणीत सूत्र असण्याची Frm(x) अट सुटली आहे. तिच्याशिवाय रिक्त अनुक्रमाच्या code वर दोन्ही bounded universal अटी रिक्तपणे खऱ्या ठरतात. येथे आवश्यक सूत्र-अट जोडली.\end{quote}
\begin{quote}\small\textbf{OLP-0286 — स्रोतनोंद.} SOL6-A065 — मूळ proof-code चाचण्या decoded fields वरच विसंबतात. येथे sequence codes, नेमका दोन-बाजूंचा sequent code, नेमकी initial/unary/binary proof-node रूपे आणि दाखवलेल्या दोन्ही नियमांचे syntax domains तपासले. लांबी दोन असली तरी extra prime factor असलेल्या 42 चा sequent code म्हणून स्वीकार किंवा initial node ला तिसरा अनावश्यक घटक जोडण्याचा स्वीकार यामुळे थांबतो.\end{quote}
\begin{quote}\small\textbf{OLP-0287 — स्रोतनोंद.} SOL6-A066 — मूळ OpenAssum च्या मार्गसंकेतांकासाठी पूर्ण subtree क्रमिकेचा code हा strict upper bound घेतला आहे. फक्त एक गृहीतक असलेल्या निष्पत्तीसाठी आवश्यक singleton मार्गाचा code त्या क्रमिकेच्या code इतकाच असतो; म्हणून मूळ कसोटी त्या गृहीतकाला उघडे मानत नाही. येथे वैध nonempty मार्गाला सर्वसाधारण sequence bound दिला. Exact node guard ने extra fields वगळले, आणि Assum ने प्रत्यक्ष subtree क्रमिकेच्या पूर्ण लांबीपर्यंत शोध घेतला.\end{quote}
\begin{quote}\small\textbf{OLP-0288 — स्रोतनोंद.} SOL6-A067 — decoded ओळी व लांबी तपासणे म्हणजे canonical क्रमिकासंकेतांक तपासणे नव्हे. एका axiom च्या singleton proof-code ला त्याच्या support बाहेरील अविभाज्य घटकाने गुणले तरी जुन्या length आणि line projections बदलत नाहीत; जुनी Deriv कसोटी ते अवैध संख्या स्वीकारते. येथे Deriv(d) आणि antecedent-list s साठी मागील Seq अट जोडली.\end{quote}
\begin{quote}\small\textbf{OLP-0292 — स्रोतनोंद.} SOL6-A070 — समशेषतेचा समान शेष हा अर्थ धन भाजकासाठी आहे. शून्य भाजकाचा totalized शेष शून्य घेतला तरी तो divisibility द्वारे दिलेल्या समशेषतेशी समतुल्य नसतो. म्हणून या व्याख्येत आणि सुनझीच्या प्रमेयात धन भाजक स्पष्ट केला. बीटा फलाच्या बांधणीत सर्व भाजक आधीपासून धन आहेत.\end{quote}
\begin{quote}\small\textbf{OLP-0326 — स्रोतनोंद.} SOL6-A091: पर्याय मूल्यांकन निश्चित x, X किंवा u वर बदलते; त्याचे कोणत्याही y वरील मूल्य डाव्या बाजूला दाखवले आहे. पूरक संचासाठी disjointness बरोबर पूर्ण आधारक्षेत्र झाकणे आवश्यक आहे. Conjunction encoding मधील X मुक्त नसण्याची अट capture टाळते.\end{quote}
\begin{quote}\small\textbf{OLP-0329 — स्रोतनोंद.} SOL6-A092: Inf डेडेकिंड अनंतता व्यक्त करते. येथे सर्वसाधारण अनंततेशी आणि Fin च्या सांततेशी सममूल्यत्वासाठी निवडीचे स्वयंसिद्धक मानणारा सामान्य संचसैद्धान्तिक संदर्भ घेतला आहे. Orbit ची पुनरावृत्ती काढल्यावर repetition-free प्रगणना मिळते.\end{quote}
\begin{quote}\small\textbf{OLP-0331 — स्रोतनोंद.} SOL6-A093: संगणनक्षम प्रगणनाचे विधान प्रभावी भाषांसाठी आणि प्रभावी सिद्धता प्रणालींसाठी आहे. द्वितीय क्रमाच्या मानक चिन्हार्थमीमांसेत निर्दोष व संपूर्ण प्रभावी प्रणाली नसते. प्रगणनीय प्रतिमानाचा लोव्हेनहाइम स्कोलेम दावा प्रगणनीय भाषेसाठी आहे.\end{quote}
\begin{quote}\small\textbf{OLP-0332 — स्रोतनोंद.} SOL6-A094: विगमन रूपबंधातील इतर मुक्त parameters सार्विकरीत्या बद्ध घेतले आहेत. Order proof मध्ये tail चा अंतर्भाव ही त्या सूत्राच्या पूर्ततेसाठी आवश्यक अट आहे; ती पुरेशी होण्यासाठी successor closure हवी. Tail स्वतः वरील सूत्र पूर्ण करतो, म्हणून ते पूर्ण करणाऱ्या सर्व संचांचा छेद योग्य order relation देतो.\end{quote}
\begin{quote}\small\textbf{OLP-0333 — स्रोतनोंद.} SOL6-A095: गैरस्वयंसिद्धकीकरणाचा निष्कर्ष प्रभावी सिद्धता प्रणालींविषयी आहे. Reduction मध्ये प्रथम क्रमाची अंकगणितीय वाक्ये वापरली असून नव्या चलांची निवड capture टाळणारी आहे.\end{quote}
\begin{quote}\small\textbf{OLP-0338 — स्रोतनोंद.} SOL6-A096: Injectivity फक्त तुलना केल्या जाणाऱ्या X वरील प्रतिबंधासाठी आहे. फलन पूर्ण प्रांतावर परिभाषित असले तरी X बाहेरील मूल्ये मनमानी असू शकतात; त्यांवर injectivity मागितलेली नाही.\end{quote}
\begin{quote}\small\textbf{OLP-0339 — स्रोतनोंद.} SOL6-A097: मूळ Inf, Count आणि Aleph-one सूत्रांचे आधी निदर्शित दोष येथे प्रत्यक्ष दुरुस्त केले आहेत. X-to-X closure, restricted injectivity, empty-set countability आणि strictly smaller cardinality च्या उपसंचांची अट आवश्यक आहेत.\end{quote}
\begin{quote}\small\textbf{OLP-0340 — स्रोतनोंद.} SOL6-A098: दुरुस्त cardinality definitions नंतर powerset coding आणि unique-code Cont वापरले आहेत. पूर्ण प्रांत continuum इतकाच असण्याची अट आता full-domain Y आणि Cont(Y) मागते; जुन्या अपुऱ्या identity-function range अटीऐवजी ही अट आहे. CH आणि NCH चे dependent explanations त्यानुसार दुरुस्त केले आहेत.\end{quote}
\begin{quote}\small\textbf{OLP-0344 — स्रोतनोंद.} SOL6-A102: फक्त कोणत्याही लॅम्डाच्या व्याप्तीत येणे हा बद्ध वापराचा निकष नाही. त्याच नावाचे संबंधित अमूर्तीकरण त्या वापराला बांधते. नामांतर करताना मुक्त वापर आणि इतर बंधन-संबंध जपले आहेत.\end{quote}
\begin{quote}\small\textbf{OLP-0345 — स्रोतनोंद.} SOL6-A103: आदेशनात चराचे फक्त मुक्त वापर बदलतात. आदेशित पदातील मुक्त चर पकडले जाऊ नयेत म्हणून आवश्यक बद्ध चरांचे आधी नामांतर केले जाते.\end{quote}
\begin{quote}\small\textbf{OLP-0347 — स्रोतनोंद.} SOL6-A105: या करीकरणाच्या न्यूनीकरणांत अमूर्त केलेली चरे परस्पर भिन्न आणि सर्व आदेशित पदांतील मुक्त चरांपासून वेगळी निवडली आहेत. आवश्यक नामांतर आधी केले जाते; अन्यथा displayed intermediate expression मुक्त चर पकडू शकते.\end{quote}
\begin{quote}\small\textbf{OLP-0350 — स्रोतनोंद.} SOL6-A111: क्लिनी प्रमाण रूपाने total primitive-recursive functions, total predicate चे लघुतमीकरण आणि search ला force करणारे शेवटचे उपयोजन पुरते आहे. या सिद्धतेत naive partial composition किंवा partial primitive-recursion construction गृहीत धरलेली नाही.\end{quote}
\begin{quote}\small\textbf{OLP-0351 — स्रोतनोंद.} SOL6-A108: आधीच्या recursive-function व्याख्येतील शून्य फलन एकस्थानी आहे. चर्चचा शून्य अंक आणि तो अंक नेहमी परत करणारे एकस्थानी फलन यांत अतिरिक्त अमूर्तीकरणाचा फरक आहे.\end{quote}
\begin{quote}\small\textbf{OLP-0352 — स्रोतनोंद.} SOL6-A109: सरळ nested application ची construction total numeric functions साठी आहे. आंशिक strict composition चे सामान्य विधान पुढील स्वतंत्र computability equivalence पूर्ण झाल्यावर मिळते; main proof येथे total closure पुरती वापरते.\end{quote}
\begin{quote}\small\textbf{OLP-0353 — स्रोतनोंद.} SOL6-A110: प्राचलांसहित प्रतिनिधी आणि प्राचले सामावलेली फलनमूल्ये यांची नावे वेगळी घेतली आहेत. मागील recursion value ला फक्त उरलेली प्राचले उपयोजली आहेत; recursion index पुन्हा उपयोजलेला नाही. Numeric total प्रकरण आणि arbitrary beta-recursion equations यांच्या सिद्धतांचे कार्य वेगळे स्पष्ट केले आहे.\end{quote}
\begin{quote}\small\textbf{OLP-0354 — स्रोतनोंद.} SOL6-A113: दिलेल्या पदातील मुक्त चर पकडले जाऊ नयेत म्हणून diagonal construction मधील अमूर्त चर आधी ताजे निवडले आहे. सर्व fixed-point expressions जपली आहेत.\end{quote}
\begin{quote}\small\textbf{OLP-0355 — स्रोतनोंद.} SOL6-A112: Displayed selector-search construction प्रथम total numeric test साठी सिद्ध केली आहे. Search थांबत नसल्यास त्याचे head computationही थांबत नाही. क्लिनी प्रमाण रूपातील अंतिम extractor आधी search force केल्याने general computability theorem पूर्ण होते; arbitrary partial test चे closure त्यानंतर घेतले आहे.\end{quote}
\begin{quote}\small\textbf{OLP-0362 — स्रोतनोंद.} SOL6-A120: एका inner substitution ची definedness मिळाल्यावर दुसऱ्याची definedness आपोआप मिळत नाही. सर्व आवश्यक नावे एका सांत forbidden संचाबाहेर निवडून, सर्वांत आतल्या अमूर्तीकरणांपासून नामांतर केले आहे. Binding-tree argument ने arbitrary alpha-equivalent प्रतिनिधींसाठी परिणामी alpha-class ची uniqueness सिद्ध केली आहे.\end{quote}
\begin{quote}\small\textbf{OLP-0363 — स्रोतनोंद.} SOL6-A121: नावांच्या यादीत एकाच नावाचे अनेक वापर असतील, तर पहिली जागा घ्यायची आहे. त्यामुळे सर्वांत जवळचे अमूर्तीकरण चराला बांधते. पुढील उलट रूपांतर फक्त प्रत्येक निर्देशांकाला उपलब्ध जागा असलेल्या पदांसाठी परिभाषित आहे.\end{quote}
\begin{quote}\small\textbf{OLP-0369 — स्रोतनोंद.} SOL6-A125: दोन आदेशनांच्या क्रमाची अदलाबदल करताना आवश्यक ताजेपणाच्या अटी आणि आदेशनाचे संयोजनसूत्र स्पष्ट केले आहे. आधी दुरुस्त केलेला अमूर्तीकरणाच्या प्रकरणातील न्यूनीत आदेशित प्रतिनिधी जपला आहे.\end{quote}
\begin{quote}\small\textbf{OLP-0371 — स्रोतनोंद.} SOL6-A126: इटा संकुचनाने फलनस्थानील अमूर्तीकरण नाहीसे होऊ शकते. त्यामुळे बीटा पुराव्यातील पहिली चार प्रकरणे तशीच लागू होत नाहीत. मुख्य triangle दावा आणि अमूर्तीकरणाच्या उपयोजनाचा सहाय्यक दावा यांचे एकत्र रचनात्मक विगमन देऊन ही उणीव पूर्ण केली आहे.\end{quote}
\begin{quote}\small\textbf{OLP-0375 — स्रोतनोंद.} SOL6-A128: मूळ e b हे पद शून्य घातांकासाठी identity मध्ये न्यूनीत होते; ते बीटा-सामान्य चर्च अंक एक नाही. बाहेरून चर्च अंकाचे दोन प्राचल लावल्याने सर्व घातांकांसाठी आवश्यक बीटा-सामान्य अंक मिळतो. शून्याच्या शून्य घाताचा येथे एक असा पुनरावर्ती संकेत आहे.\end{quote}
\begin{quote}\small\textbf{OLP-0394 — स्रोतनोंद.} SOL6-A140: दिलेल्या मोडल कोष्टकांत मधल्या U मूल्यावरील विरोधाभासाची संभाव्यता T होते; तिचे निषेधन F होते, मूळात दिलेले U नाही. मॅट्रिक्सचा चार-संयोजकीय भाषाखंड स्पष्ट केला. मूळ 1920 चा व्याख्यानवृत्तांत प्रत्यक्ष पाहून ऐतिहासिक वर्ष दुरुस्त केले; उद्धरणातील noon म्हणजे दुपारी बारा वाजता.\end{quote}
\begin{quote}\small\textbf{OLP-0395 — स्रोतनोंद.} SOL6-A141: दोन्ही क्लिनी मॅट्रिक्स येथे मानक L0 च्या असत्यता-स्थिर नसलेल्या चार-संयोजकीय खंडावर आहेत. सर्व चलांना U देणाऱ्या मूल्यनिर्णयनावरील सर्वतः सत्य सूत्र नसण्याचा पुरावा या व्याप्तीत लागू होतो.\end{quote}
\begin{quote}\small\textbf{OLP-0397 — स्रोतनोंद.} SOL6-A142: फक्त निर्दिष्ट मूल्ये बदलून लूकासेविच मॅट्रिक्सला R-Mingle म्हणणे अयोग्य होते. R-Mingle चे वेगळे सशर्त कोष्टक Kooi आणि Tamminga यांच्या primary paper मधील विभाग 4, पृ. 10 शी जुळवले; जोडलेल्या असत्यता-स्थिराचे मूल्य F स्पष्ट केले. LP आणि Hallden यांच्या classical tautologies ची समानता केवळ समान चार-संयोजकीय खंडात आहे.\end{quote}
\begin{quote}\small\textbf{OLP-0400 — स्रोतनोंद.} SOL6-A143: अनंतमूल्य आणि सांतमूल्य लूकासेविच मॅट्रिक्स येथे मानक L0 च्या चार-संयोजकीय खंडावर आहेत. याने मागील त्रिमूल्य मॅट्रिक्सशी भाषेची समानता जपली. सांत आधारसंचांसाठीचे प्रतिलोम हे आधीचे योग्य restriction जपले.\end{quote}
\begin{quote}\small\textbf{OLP-0404 — स्रोतनोंद.} SOL6-A144: प्रमेयाच्या क्रमवर्तीत फक्त निर्दिष्ट सत्यतामूल्यांच्या जागी सूत्र असते; इतर जागा रिकाम्या असतात. हा आवश्यक scope स्पष्ट केला. व्युत्पादन सांत वृक्ष आहे आणि त्यात तार्किक व संरचनात्मक नियम दोन्ही येतात. दाखवलेले Gamma उदाहरण सांत क्रमिकेसाठी; अनंत आधारसंचासाठी आधीच्या सांत-उपसंच व्याख्येचा वापर होतो.\end{quote}
\begin{quote}\small\textbf{OLP-0411 — स्रोतनोंद.} SOL6-A146: क्रमाने केलेल्या आदेशनाच्या दोन्ही मध्यवर्ती displays मध्ये संपूर्ण सूत्र दुसऱ्या आदेशनाच्या आत घेतले. पूर्वपक्षातील चलही बदलले पाहिजेत; केवळ उत्तरपक्ष बदलणे अपुरे आहे. अंतिम विस्तारित सूत्रे मूळातच योग्य होती.\end{quote}
\begin{quote}\small\textbf{OLP-0417 — स्रोतनोंद.} SOL6-A149: रूपबंधाच्या notation मध्ये केवळ विधानचल बदलतात; असत्यासारखे तार्किक स्थिरांक बदलत नाहीत. आदेशन-दृष्टांताच्या वैधतेच्या सिद्धतेतील नवीन valuation सर्व विधानचलांवर पूर्ण केली आणि तिचे truth sets प्रतिमानाच्या W मधील जगांपुरते स्पष्ट केले.\end{quote}
\begin{quote}\small\textbf{OLP-0421 — स्रोतनोंद.} SOL6-A150: एका असत्य दृष्टांताचे प्रदर्शन म्हणजे नेमका एकच दृष्टांत असत्य आहे असा दावा नाही. त्याच एक-जगी प्रतिमानात असत्य स्थिरांक घातल्यावर आणखी असत्य दृष्टांत मिळतो.\end{quote}
\begin{quote}\small\textbf{OLP-0423 — स्रोतनोंद.} SOL6-A151: D च्या proof मध्ये V हा model tuple मध्ये आहे; त्याची विशेष निवड करण्याची गरज नाही, असा independence चा अचूक अर्थ घेतला. B आणि 5 च्या proofs मध्ये वापरलेली पद्धत प्रतिपरिवर्तन आहे.\end{quote}
\begin{quote}\small\textbf{OLP-0426 — स्रोतनोंद.} SOL6-A153: मानक भाषांतराच्या प्रत्येक मोडल पायरीत नवीन बद्ध चल निवडण्याची अट स्पष्ट केली. Monadic हे द्वितीय-क्रम संख्यापकांपुरते बंधन आहे; व्यक्तिचलांवरील प्रथम-क्रम संख्यापक अनुमत आहेत. परावर्तितेच्या उदाहरणात व्यक्तिचलाचे मूल्य, प्रत्यक्ष जग आणि successor set वेगळे करून नेमणूक स्पष्ट केली.\end{quote}
\begin{quote}\small\textbf{OLP-0428 — स्रोतनोंद.} SOL6-A155: सामान्य मोडल तर्कशास्त्र मिळण्याचा दावा K आणि आवश्यक असल्यास Dual हे स्वयंसिद्धके घेतल्याच्या संदर्भात आहे. निर्दोषतेसाठी फक्त characteristic सूत्रे नव्हे, तर सर्व स्वयंसिद्धकांचे सर्व आदेशन-दृष्टांत प्रतिमानाच्या सर्व जगांत सत्य असणे आवश्यक आहे.\end{quote}
\begin{quote}\small\textbf{OLP-0431 — स्रोतनोंद.} SOL6-A156: दहा ओळींच्या K-सिद्धतेत ओळ 6 मधील implication चे बाह्य कोष्ठक उघडले आणि ओळी 8 व 9 मधील प्रत्येकी एक जादा बंद कोष्ठक काढले. ओळी 3 व 5 चे K-दृष्टांत आणि सर्व inference references आधीच योग्य होते; ते जपले.\end{quote}
\begin{quote}\small\textbf{OLP-0432 — स्रोतनोंद.} SOL6-A157: replacement-rule च्या संक्षिप्त नोंदीत “नवीन सूत्र for जुने सूत्र” असा क्रम वापरला आहे. म्हणून C(A) वरून C(B) मिळवताना label B for A आहे; हा क्रम खालील double-negation examples मधील p for not-not-p शी जुळतो.\end{quote}
\begin{quote}\small\textbf{OLP-0436 — स्रोतनोंद.} SOL6-A159: निर्दोषता विभागाच्या प्रस्तावनेतील axioms सत्य असण्याची अट पुढील प्रमेयाप्रमाणे सर्व आदेशन-दृष्टांत सर्व जगांत सत्य असण्याची आहे. मूळ शब्दांकनातील characteristic सूत्र आणि पूर्ण schema यांतील संदिग्धता दूर केली; प्रमेय व त्याची सिद्धता बदललेली नाही.\end{quote}
\begin{quote}\small\textbf{OLP-0438 — स्रोतनोंद.} SOL6-A160: प्रणालीतील derivability च्या व्याख्येचा संदर्भ त्याच्या योग्य विभागाकडे वळवला. गृहीतकांची संख्या शून्य असू शकते हे स्पष्ट केले; रिकाम्या premise set साठी कोणतीही गृहीतके न वापरणारी सिद्धता मूळ प्रणालीतील सिद्धताच असते.\end{quote}
\begin{quote}\small\textbf{OLP-0440 — स्रोतनोंद.} SOL6-A161: सापेक्ष विसंगतता आणि योग्य प्रतिमानांचा वर्ग. एखाद्या मजबूत प्रणालीच्या सापेक्ष विसंगत Gamma मनमानी Kripke प्रतिमानात satisfiable असू शकतो. उदाहरणार्थ KT च्या सापेक्ष Boxp आणि not-p विसंगत आहेत, पण nonreflexive model मध्ये एकत्र सत्य होतात. No-model दावा संबंधित system च्या सर्व axiom instances सर्व worlds मध्ये सत्य असलेल्या models पुरता केला. पुढील consistency definition आणि तीन facts बदलले नाहीत.\end{quote}
\begin{quote}\small\textbf{OLP-0442 — स्रोतनोंद.} SOL6-A162: सर्वसाधारण canonical construction मध्ये अनेक जगांची गरज. मूळ प्रेरक वाक्य कोणतेही A सत्य करण्यासाठी multiple worlds आवश्यक आहेत असे म्हणत होते. p सारखी सूत्रे single-world model मध्ये सत्य होतात; मात्र सर्व consistent formulas साठी single-world construction पुरेशी नाही. आवश्यकतेचा scope सर्वसाधारण रचनेवर केला. प्रत्यक्ष canonical world set सर्व complete consistent sets घेतो हा पुढील युक्तिवाद जपला.\end{quote}
\begin{quote}\small\textbf{OLP-0443 — स्रोतनोंद.} SOL6-A163: Complete consistent sets मधील disjunction आणि joint scope. Disjunction membership च्या iff proof मध्ये मूळाने फक्त forward direction दिली होती. दोन्ही classical introduction tautologies व deductive closure ने reverse direction जोडली. Biconditional च्या reverse proof मधील दोन्ही नसतील या Marathi wording मुळे not-both आणि both-not यांची गल्लत होत होती; दोन्ही statements एकत्र true असण्याच्या प्रत्येक शक्यतेचा निषेध स्पष्ट केला. आधीचे negation आणि biconditional-antecedent repairs जपले.\end{quote}
\begin{quote}\small\textbf{OLP-0444 — स्रोतनोंद.} SOL6-A164: Lindenbaum proof मधील रिकाम्या finite premise list चा निर्देशांक. Finite witness list शून्य-लांबीची असू शकते. मूळ proof मध्ये त्या list मधील indices चा largest घेतला होता, पण empty list ला largest नसतो. Max of zero plus the listed indices ही नेमकी convention दिली; k=0 साठी n=0 मिळतो. Increasing chain मधील finite witness containment आणि सर्व Delta\_n ची consistency proof जपली. पूर्वीची exhaustive enumeration साठी length-at-most-n दुरुस्ती योग्य.\end{quote}
\begin{quote}\small\textbf{OLP-0445 — स्रोतनोंद.} SOL6-A165: Canonical accessibility चे iff आणि शून्य गृहीतकांचा boxing lemma. Diamond branch मधील explanatory clause ने R खरे होण्यासाठी sufficient अट दिली होती; अपेक्षित implication साठी R ते condition ही दिशा आवश्यक आहे. औपचारिक पुढील definition प्रमाणे iff स्पष्ट केला. Boxing lemma च्या k=0 प्रकरणात RK नव्हे तर NEC थेट लागू केला; k>=1 साठी मूळ RK proof जपला. Modal equivalence च्या वाक्यात missing शी जोडले. पूर्वीचे k/n आणि system-subscript repairs योग्य.\end{quote}
\begin{quote}\small\textbf{OLP-0446 — स्रोतनोंद.} SOL6-A166: Canonical model चा nonempty domain आणि valuation ची world व्याप्ती. Normal logic च्या आधीच्या definition मध्ये inconsistent full logic वगळलेली नाही. अशा Sigma साठी complete consistent sets नसतात आणि W रिकामा होतो, जे basic model definition शी विसंगत आहे. Canonical model construction consistent normal systems साठी केली आणि Lindenbaum ने W nonempty दाखवले. Valuation truth set स्पष्टपणे Delta in W पुरता केला. Inconsistent systems साठी completeness चा explosion case स्वतंत्र हाताळता येतो; रिकामा set एक वैध Kripke model असल्याचा दावा नाही.\end{quote}
\begin{quote}\small\textbf{OLP-0447 — स्रोतनोंद.} SOL6-A167: Truth lemma मधील accessibility bridge आणि विगमन गृहीतक. Diamond च्या forward case मध्ये Box-inverse containment वरून Diamond-image containment मिळवण्यासाठी box-iff-diamond lemma आवश्यक आहे; जुन्या prop-diamond reference ने ती containment थेट मिळत नाही. योग्य lemma दाखवला. Reverse case मध्ये B चे membership आणि Diamond satisfaction यांच्या मध्ये induction hypothesis चे B satisfaction पाऊल जोडले. Theorem मध्ये consistent normal Sigma व Delta in W scope स्पष्ट केला; सर्व Boolean व Box cases आणि prior probAnd repair जपले.\end{quote}
\begin{quote}\small\textbf{OLP-0448 — स्रोतनोंद.} SOL6-A168: Determination theorem च्या hypotheses आणि canonical proof पद्धतीचा scope. Canonical model अस्तित्वासाठी consistent normal Sigma ही hypothesis determination theorem मध्ये स्पष्ट केली. K ची consistency पूर्वीच्या soundness आणि falsity च्या nonvalidity वरून मिळते, म्हणून K-completeness proof ला ही अट लागू होते. Arbitrary class completeness साठी canonical model membership नेहमी आवश्यक असे नाही; ही अट या canonical-model proof पद्धतीसाठी आवश्यक असल्याचे शब्दांकन केले. Determination iff आणि K corollary ची दोन्ही directions जपल्या.\end{quote}
\begin{quote}\small\textbf{OLP-0449 — स्रोतनोंद.} SOL6-A169: Canonical frame theorem ची consistency आणि weak-density witness. Canonical frame theorem आणि supplementary proposition सुसंगत normal systems पुरते स्पष्ट केले. Functionality proof मध्ये biconditional schema वरून D मिळतो हे नोंदले; partial functionality plus seriality देऊन exact one successor मिळतो. Weak-density proof मधील finite Diamond-premise list रिकामी मानल्यास त्याचे displayed nonempty conjunction steps अपूर्ण होतात. Box-inverse Delta1 हा consistent Delta2 चा subset असल्याने inconsistent witness ला किमान एक Diamond premise लागतो; m>=1 व n>=0 हे नेमके bounds दिले. पाच canonical relation proofs व all remaining displayed steps जपले.\end{quote}
\begin{quote}\small\textbf{OLP-0451 — स्रोतनोंद.} SOL6-A170: finite-model search मध्ये निश्चित frame व सूत्रातील वेगवेगळ्या चलांच्या valuations एवढाच count आहे. Frame-class test साठी दिलेली सांत schema-list आणि stopping bound साठी effective computation स्पष्ट केली. Universal-frame प्रेरक quotient केवळ लक्ष्य सूत्रातील सांत variables वर घेतला; सर्व भाषेतील variables वर agreement मागितल्यास finite quotient मिळेलच असे नाही. Formal filtration मध्ये सर्व चलांची valuation class-image set ने होते, पण selected formulas पुरते truth preservation असते. Box induction ची दुसरी दिशा पूर्ण केली.\end{quote}
\begin{quote}\small\textbf{OLP-0456 — स्रोतनोंद.} SOL6-A171: प्रणालीतील प्रत्येक सूत्र प्रतिमानाच्या प्रत्येक जगावर सत्य असणे आणि चौकट-गुणधर्म असणे हे स्वतंत्र आहेत. S5 च्या कोणत्याही प्रतिमानात लक्ष्य सूत्र सत्य असेल तर त्याचा निषेध सिद्ध होत नाही; संपूर्णतेवरून परावर्ती युक्लिडी प्रतिमान मिळते. त्यानंतर वैश्विक प्रतिमान आणि सांत निस्यंदनाचे मूळ पाऊल लागू होते.\end{quote}
\begin{quote}\small\textbf{OLP-0457 — स्रोतनोंद.} SOL6-A172: निर्णेयतेच्या प्रस्तावनेत दिलेली सांत आणि प्रभावी रूपबंध-यादी स्पष्ट केली. S5 च्या प्रतिमान-शोधात प्रत्येक आकारासाठी निश्चित जगांची नावे, वैश्विक संबंध आणि लक्ष्य सूत्रातील चलांचीच सांत मूल्यांकन-यादी वापरली जाते; उर्वरित चलांचे मूल्यांकन रिक्त ठेवता येते.\end{quote}
\begin{quote}\small\textbf{OLP-0458 — स्रोतनोंद.} SOL6-A173: अतिरिक्त अटीमुळे प्राप्यतेच्या जोड्या वाढत नाहीत; त्या कमी होतीलच असे नाही. दोन्ही निस्यंदने समान होऊ शकतात, म्हणून अधिक सूक्ष्मपणाच्या तुलनेत समतेची शक्यता स्पष्ट केली.\end{quote}
\begin{quote}\small\textbf{OLP-0459 — स्रोतनोंद.} SOL6-A174: रिक्त संचही मोडलदृष्ट्या बंद असतो; अनंततेचा दावा रिक्त नसलेल्या संचांसाठी आहे. संक्रमकतेच्या सिद्धतेत C1 ज्या modal सूत्रांवर परिमाणित आहे ती Gamma मध्ये असल्याची अट स्पष्ट घेतली; त्यातून उपसूत्र व मोडल बंदतेचा उपयोग वैध ठरतो.\end{quote}
\begin{quote}\small\textbf{OLP-0462 — स्रोतनोंद.} SOL6-A175: नव्या पूर्वचिन्हाच्या अटीचे प्रतिउदाहरण दाखवणाऱ्या केवळ Box-पर्यायात असत्य Box सूत्रावर लागू होणाऱ्या नियमाचे लेबल असत्य Box असे दुरुस्त केले. उदाहरणातील जाणीवपूर्वक शिथिल केलेली नवेपणाची अट बदललेली नाही.\end{quote}
\begin{quote}\small\textbf{OLP-0464 — स्रोतनोंद.} SOL6-A176: निर्दोषतेच्या सिद्धतेत मोडल प्रतिमान आणि मोडल पूर्तीची संज्ञा वापरली. नियमानुसार तयार झालेल्या शाखेची पूर्वचिन्हे आरंभीच्या रिक्त नसलेल्या खंडांखाली बंद असतात. म्हणून नवे पूर्वचिन्ह वाढवताना त्याचे आधीपासूनचे वंशज नसतात आणि फक्त पालकापासून नव्या पूर्वचिन्हापर्यंतची प्राप्यता तपासणे पुरेसे असते.\end{quote}
\begin{quote}\small\textbf{OLP-0465 — स्रोतनोंद.} SOL6-A179: अतिरिक्त नियमांच्या उदाहरणात 5 स्वयंसिद्धकाचे योग्य सूत्र आणि त्याचा बंद टॅब्लो दिला. मूळातील Box A वरून Box Diamond A हे वेगळे सूत्र होते. योग्य Diamond A गृहीतकासाठी सत्य Diamond नियमाने दुसरे नवे पूर्वचिन्ह आणले; शेवटचा असत्य Diamond नियम त्याच वापरलेल्या पूर्वचिन्हावर लागू होतो.\end{quote}
\begin{quote}\small\textbf{OLP-0468 — स्रोतनोंद.} SOL6-A180: संपूर्णतेच्या रचनेत बंद आणि संपूर्ण शाखा वेगळ्या केल्या. गृहीतकांचे समान पूर्वचिन्ह 1 स्पष्ट केले; रिक्त संचाचे प्रकरण वेगळे हाताळले. सांत रचनेचे थांबणे मोडल खोली आणि सांत संततीवरून स्पष्ट केले. असांत संचासाठी प्रत्येक सांत उपसंचाची पूर्ततायोग्यता, सुसंगततेची सांत सिद्धता, लिंडेनबाउमचे पूर्वप्रमेय आणि कॅनॉनिकल सत्यता पूर्वप्रमेय वापरून सर्वसाधारण निष्कर्ष जपला.\end{quote}
\begin{quote}\small\textbf{OLP-0469 — स्रोतनोंद.} SOL6-A181: निर्णय प्रक्रियेतील नवे पूर्वचिन्ह देणारा नियम प्रत्येक पूर्वचिन्हयुक्त चिन्हांकित सूत्रासाठी एकदा लावायचा असतो; वापरलेली पूर्वचिन्हे घेणाऱ्या नियमाची व्याप्ती त्याच्या तात्काळ विस्तारांपुरती आहे. थांबण्याचा सांत इनपुटवरील पुरावा स्पष्ट केला. मनमानी गृहीतकसंचासाठी संपूर्णतेचा निष्कर्ष जपून, प्रत्यक्ष सांत निर्णय प्रक्रियेचे विधान सांत गृहीतकसंचापुरते केले. उदाहरणांतील वापरलेली पूर्वचिन्हे 1.n या आकाराची आहेत हे स्पष्ट केले; सर्व सहा टॅब्लो आणि दोन्ही प्रतिमाने जपली.\end{quote}
\begin{quote}\small\textbf{OLP-0471 — स्रोतनोंद.} SOL6-A183: S5 साठी कट-मुक्त कलन ज्ञात नसल्याचा मूळ दावा ऐतिहासिक मसुद्यापुरता केला. प्रत्यक्ष वाचलेल्या Aghaei आणि Mohammadi यांच्या 2018 लेखाचा स्रोत जोडला; त्यातील G3S5 आणि पुढे दिलेला विशिष्ट LK-विस्तार यांची व्याप्ती वेगळी ठेवली.\end{quote}
\begin{quote}\small\textbf{OLP-0473 — स्रोतनोंद.} SOL6-A184: Dual च्या सिद्धतेतील पहिल्या नकार-पायरीत उजव्या बाजूचे सूत्र नकारासह डावीकडे जाते. त्यामुळे चुकीचे उजवा-नकार नियम-नाव डावा-नकार असे केले. अंतिम संयोग-नियमाच्या पायरीत द्विशर्ती दोन सशर्त सूत्रांच्या संयोगाचे संक्षिप्त रूप म्हणून घेतली आहे हे स्पष्ट केले; मूळ सिद्धतेतील सूत्रे जपली.\end{quote}
\begin{quote}\small\textbf{OLP-0474 — स्रोतनोंद.} SOL6-A185: मोडल क्रमवर्ती प्रणाल्यांच्या तक्त्यात मूळ नियमांची नावे primitive Box/Diamond निवडीनुसार केली; दोन्ही primitive असताना दोन्ही नियम दाखवले. Diamond-only कट-सिद्धतेतील डाव्या सूत्रांच्या अदलाबदलीचे चुकीचे उजवे नियम-नाव डावे केले. कटविना असंपूर्णतेचा दावा या विभागातील विशिष्ट कलनापुरता स्पष्ट केला.\end{quote}
\begin{quote}\small\textbf{OLP-0481 — स्रोतनोंद.} SOL6-A187: इतिहासात कालिक स्थान स्पष्ट होण्यासाठी प्रत्येक अवस्था-घटना स्वतंत्र मानली; प्रणालीची स्थिती पुन्हा आली तरी तिच्या वेगळ्या वेळच्या घटना वेगळ्या आहेत. Boolean उपसूत्रे त्याच बिंदू व इतिहासावर तपासायची हे स्पष्ट केले. शेवटच्या उदाहरणातील कधीच पुढे सत्य होत नाही हा सार्विक नकार दुरुस्त केला.\end{quote}
\begin{quote}\small\textbf{OLP-0484 — स्रोतनोंद.} SOL6-A188: मोडल-विरहिततेची अट सर्व कर्त्यांसाठी केली. समूहज्ञानाचा संयोग फक्त सीमित समूहासाठी साधा formula abbreviation आहे हे स्पष्ट केले; अनंत समूहाचे स्वतंत्र अर्थनिर्धारण व सामायिक ज्ञानाचा भाषा-विस्तार वेगळा दिला. वैयक्तिक आणि सामायिक ज्ञान नेहमी काटेकोरपणे वेगळे असतात हा अर्थ टाळला.\end{quote}
\begin{quote}\small\textbf{OLP-0485 — स्रोतनोंद.} SOL6-A192: प्रतिमानातील R हा कर्त्यानुसार निर्देशित संबंधांचा परिवार आहे हे स्पष्ट केले. दोन कर्त्यांचे संबंध सारखे असले तरी त्यांची नावे या परिवारात राखली जातात. V(p) हा W चा उपसंच आणि कर्त्याचे a हे गणिती चिन्ह स्पष्ट केले.\end{quote}
\begin{quote}\small\textbf{OLP-0486 — स्रोतनोंद.} SOL6-A189: संक्रमण-संवरणाच्या powers मध्ये मूळ zero-based indexing ऐवजी पूर्वीच्या व्याख्येशी जुळणारे R\textasciicircum{}1=R आणि n>0 union वापरले. Positive-path आणि reflexive-path conventions वेगळ्या ठेवल्या; मूळ common-knowledge semantics बदलली नाही.\end{quote}
\begin{quote}\small\textbf{OLP-0487 — स्रोतनोंद.} SOL6-A191: सकारात्मक व नकारात्मक आत्मनिरीक्षणाच्या प्रतिउदाहरणांना वेगळे पूर्वांग दिले. एखादी गोष्ट माहीत असणे हे नकारात्मक आत्मनिरीक्षणाचे पूर्वांग नाही; त्या तत्त्वासाठी माहीत नसणे आवश्यक आहे. तक्त्यातील मूळ सूत्रे जपली.\end{quote}
\begin{quote}\small\textbf{OLP-0488 — स्रोतनोंद.} SOL6-A190: द्विअनुकरणासाठी graphs समरूप असण्याची गरज नाही; forth/back उत्तरवर्ती अनुरूपता आवश्यक आहे हे स्पष्ट केले. चित्रात न दिलेल्या valuations आणि इतर कर्त्यांचे संबंध नेमके निवडल्यावरच दाखवलेला तीन-जोड्यांचा संबंध द्विअनुकरण आहे असे मांडले; मूळ graph जपला.\end{quote}
\begin{quote}\small\textbf{OLP-0490 — स्रोतनोंद.} SOL6-A193: p ची घोषणा आणि p असूनही b ला p माहीत नाही याची घोषणा समान प्रतिमान देतात हा मूळ दावा चुकीचा होता. पहिल्यात दोन जगे, दुसऱ्यात फक्त w1 उरते; actual arrows व valuations वरून हे स्पष्ट केले. असत्य घोषणा-पूर्वांगात उरलेल्या प्रतिमानात मूल्यांकन करायचे नाही हेही स्पष्ट केले; म्हणून रिकामा restriction undefined point निर्माण करत नाही.\end{quote}
\begin{quote}\small\textbf{OLP-0493 — स्रोतनोंद.} SOL6-A345: SOL6-A195 मधील गणिताने वैध पण ऐच्छिक अपरिमेयता-सिद्धता भाषांतरातून मागे घेतली. मूळाची संक्षिप्त मांडणी, निवडलेली मूल्ये आणि घाताची गणना जपली; विस्तृत सिद्धता अंमलात न आणलेली सुधारणा-सूचना म्हणून नोंदवली आहे.\end{quote}
\begin{quote}\small\textbf{OLP-0499 — स्रोतनोंद.} SOL6-A196: अंतःप्रज्ञावादी रचनाशील पडताळणी आणि साधारण epistemic knowledge operator यांतील फरक स्पष्ट केला. मूळ disjunction चे स्पष्टीकरण सामान्य knowledge operator साठी चुकीचे होते; मागील प्रत्यक्ष प्रतिमानातील counterexample आणि पुढील विभागातील योग्य local forcing clauses दिल्या.\end{quote}
\begin{quote}\small\textbf{OLP-0523 — स्रोतनोंद.} SOL6-A213: सर्वांत लहान पूर्वांग-गोलकाचा उल्लेख आकृतीपुरता स्पष्ट केला. दोन शेजारच्या दर्शवलेल्या थरांची जवळीक आणि मनमानी गोलकांतील जवळीक वेगळी आहे; सामान्य गोलक-अट सर्वात जवळचे जग नसण्यालाही लागू होते.\end{quote}
\begin{quote}\small\textbf{OLP-0524 — स्रोतनोंद.} SOL6-A214: रिकाम्या गोलकाची परवानगी व अनंत उपकुलांवरील closure स्पष्ट केली. सर्वांत लहान पूर्वांग-गोलक नसण्याचा पूर्ण प्रतिमान दिला; फक्त उतरत्या क्रमिकेवरून तो निष्कर्ष काढता येत नाही. पूर्वीचा लहान पूर्वांग-गोलकांसाठीचा source repair जपला.\end{quote}
\begin{quote}\small\textbf{OLP-0525 — स्रोतनोंद.} SOL6-A217: या आकृत्यांत सर्वाधिक जवळची पूर्वांग-जगे उपलब्ध आहेत, ही अट स्पष्ट केली. सर्वांत जवळचे जग नसलेल्या प्रतिमानांत सामान्य गोलक-अटच वापरायची; कठोर अभिव्यंजनाशी अवश्यतेची तुलना S5 च्या संदर्भात आहे.\end{quote}
\begin{quote}\small\textbf{OLP-0533 — स्रोतनोंद.} SOL6-A223: ही स्तराधारित संचरचनेची अनौपचारिक रूपरेषा आहे. रसेलचा तोच युक्तिवाद रोखणे आणि संपूर्ण औपचारिक उपपत्तीची सुसंगतता सिद्ध करणे वेगळे दावे आहेत; येथे पहिल्या दाव्याचे कारण दिले आहे.\end{quote}
\begin{quote}\small\textbf{OLP-0534 — स्रोतनोंद.} SOL6-A224: आरंभीच्या टप्प्याला आधीचे टप्पे नसतात. संचयी मांडणीमुळे रसेलचा अमर्याद संच तयार होत नाही हा युक्तिवाद आहे; कोणत्याही औपचारिक संच उपपत्तीच्या पूर्ण सुसंगततेची सिद्धता नाही.\end{quote}
\begin{quote}\small\textbf{OLP-0536 — स्रोतनोंद.} SOL6-A226: शेवटच्या चरणात सदस्यत्वाचा द्वितीय-क्रम संख्यापक असलेले संकल्पनारचनेचे उदाहरण वापरले आहे; पूर्ण रूपबंधापेक्षा या विशिष्ट उदाहरणाची उपलब्धता आवश्यक आहे. मर्यादित रूपबंधांची व्याप्ती वेगळी असू शकते; Sean Walsh, Fragments of Frege’s Grundgesetze and Gödel’s Constructible Universe, विभाग 2, व्याख्या 2.1--2.2, \url{https://arxiv.org/abs/1407.3861}, येथे ही मर्यादा स्पष्ट केली आहे.\end{quote}
\begin{quote}\small\textbf{OLP-0539 — स्रोतनोंद.} SOL6-A227: स्वसदस्यत्व-विरोधी युक्तिवाद प्रथम निर्माण-टप्पा आणि त्याच्या आधी तयार झालेले सदस्य वापरतो. सदस्यांचा शेवटचा टप्पा असावा किंवा पुढचा टप्पा लगतचा असावा अशी अतिरिक्त अट नाही.\end{quote}
\begin{quote}\small\textbf{OLP-0543 — स्रोतनोंद.} SOL6-A229: मूळ मनमाना अनंतता-साक्षी संचावर एकास-एकपणा थेट मानण्याऐवजी लघुतम संवरणावर तो सिद्ध केला. नियमिततेचे स्वयंसिद्धक वापरलेले नाही. इटाय बेन-याकोव यांच्या MATH 770: Foundations of Mathematics, विभाग 5.2.1, पृ.\textasciitilde{}86, मधील लघुतम संच व त्याच्या सदस्यांचे गुणधर्म या मार्गाशी जुळतात: \url{https://people.math.wisc.edu/~awmille1/old/m571-08/ben-yaacov.pdf}.\end{quote}
\begin{quote}\small\textbf{OLP-0550 — स्रोतनोंद.} SOL6-A230: सुक्रमणाचा संबंध दिलेल्या संचावर आहे. उपसंचात पूर्ववर्ती नसणे आणि त्यातील इतर सर्व सदस्यांपेक्षा लहान असणे यांतील भाषांतरातील भेद स्पष्ट केला. मूळाची संक्षिप्त सिद्धता जपली आहे; स्वतंत्र विस्तार केलेला नाही.\end{quote}
\begin{quote}\small\textbf{OLP-0551 — स्रोतनोंद.} SOL6-A231: द्विदिश अटी आणि विभक्तीकरणाचा मराठी शब्द दुरुस्त केला. मूळातील वैध प्रतिलोम-तर्क, प्रतिमेचा समतुल्य संकेत आणि संक्षिप्त तुलना-सिद्धता जपली आहेत; त्यांचे स्वतंत्र स्पष्टीकरण भाषांतरात घातलेले नाही.\end{quote}
\begin{quote}\small\textbf{OLP-0554 — स्रोतनोंद.} SOL6-A234: संच-फलनाच्या प्रतिमेत पदमूल्य वापरताना घटक फलनाच्या प्रांतात असावा ही अट स्पष्ट केली. मनमाना संच घेतल्यास प्रतिमा फलनाच्या आलेखावरून परिभाषित करून विभक्तीकरणाने मिळते; यासाठी प्रतिस्थापन लागत नाही.\end{quote}
\begin{quote}\small\textbf{OLP-0569 — स्रोतनोंद.} SOL6-A245: मूळाच्या प्रतिमा-तुलनेत निवडीची पुरेशी अट स्पष्ट केली. आच्छादक फलनापासून उलट दिशेचे एकास-एक फलन निवडीच्या स्वयंसिद्धकाखाली मिळते; प्रतिस्थापनाला स्वतः निवड आवश्यक आहे किंवा ही तुलना पूर्ण निवडीशी समतुल्य आहे असा दावा नाही.\end{quote}
\begin{quote}\small\textbf{OLP-0573 — स्रोतनोंद.} SOL6-A246: सर्व संक्रमक संचांत विभक्तीकरण सत्य असते हा मूळ दावा चुकीचा आहे. आवश्यक उपसंच-समावेशकत्वाची अट सिद्धतेच्या लघुतम प्रतिमानात आणि सरावाच्या विधानात जोडली. सरावाची सिद्धता दिलेली नाही. सांत स्वयंसिद्धकीकरणीयतेच्या पहिल्या निष्कर्षात सुसंगततेची अट आणि प्रतिमानांच्या अरिक्ततेचा आधार स्पष्ट केला.\end{quote}
\begin{quote}\small\textbf{OLP-0579 — स्रोतनोंद.} SOL6-A249: घातांकनाच्या येथे दिलेल्या दोन्ही व्याख्यांत आधार शून्येतर असण्याची अट स्पष्ट केली. शून्य आधारासाठी सीमा-सूत्रातील संयोगात शून्य घातांकाचा समावेश केल्यास चुकीचा परिणाम मिळतो; आधीच्या सरावप्रश्नात ही अट होती, पण व्याख्येत नव्हती. मूळ घातांकन सूत्रे व न सोडवलेला सराव जपले.\end{quote}
\begin{quote}\small\textbf{OLP-0597 — स्रोतनोंद.} SOL6-A263: निवड नसलेल्या उदाहरणात आधीच्या सुक्रमित संचांकांवर आधारित बेथ संकेताऐवजी वास्तव संख्यांचा नेमका दावा दिला. फेफरमन–लेव्ही यांच्या मूळ सारांशात, ZF सुसंगत असेल तर वास्तव संख्यांचा संच गणनीय संचांचे गणनीय युतीकरण असणेही सुसंगत आहे, असे म्हटले आहे: Notices of the American Mathematical Society 10 (1963), पृ. 593, सारांश 63T-389, https://www.ams.org/journals/notices/196310/196310FullIssue.pdf. येथे पूर्ण forcing सिद्धतेचे स्वतंत्र प्रमाणीकरण केलेले नाही.\end{quote}
\begin{quote}\small\textbf{OLP-0622 — स्रोतनोंद.} SOL6-A272: कँटर यांच्या उपचाराच्या आरंभाचे स्रोतामधील 1894 हे वर्ष 1884 केले आहे. ऑलिव्हर डाइझर यांच्या Einführung in die Mengenlehre मधील चरित्रकालक्रम, पृ. 218, मे 1884 मध्ये हाल येथील आरोग्यधामातील वास्तव्य नोंदवतो; लिओपोल्डिनाचे सदस्यचरित्रही 1884 पासून आजार आणि पुन्हा पुन्हा उपचार नोंदवते: https://www.leopoldina.org/en/members/member-list/detail/georg-cantor. येथे संपूर्ण वैद्यकीय नोंदींचे स्वतंत्र पुनरावलोकन केलेले नाही.\end{quote}
\begin{quote}\small\textbf{OLP-0624 — स्रोतनोंद.} SOL6-A273: गेंटझेन यांच्या तार्किक निगमनावरील I–II लेखांसाठी Springerची मूळ प्रकाशकनोंद 21 जुलै 1933 हा प्राप्तिदिनांक आणि 1935 हे प्रकाशनवर्ष देते. म्हणून 1934 शी जोडलेले निर्मितीवर्ष सांगण्याऐवजी निश्चित प्रकाशनवर्ष दिले आहे. मूळ लेख: https://doi.org/10.1007/BF01201353 आणि https://doi.org/10.1007/BF01201363. येथे संपूर्ण चरित्राचे स्वतंत्र ऐतिहासिक प्रमाणीकरण केलेले नाही.\end{quote}
\begin{quote}\small\textbf{OLP-0625 — स्रोतनोंद.} SOL6-A274: गोडेल यांचा प्रबंध 1929 मध्ये पूर्ण झाला आणि डॉक्टरेटची पदवी 1930 मध्ये मिळाली, हा भेद स्पष्ट केला आहे. प्रबंधानंतर वर्षभरात मिळालेल्या अपूर्णता निष्कर्षांचा निर्देश प्रबंधाशीच जोडला आहे; 1931 प्रकाशनवर्ष जपले आहे. व्हिएन्ना विद्यापीठाचे इतिहासवर्णन आणि IAS पदवीनोंद: https://geschichte.univie.ac.at/de/artikel/von-der-berechenbarkeit-zum-modernen-computer आणि https://www.ias.edu/scholars/godel.\end{quote}
\begin{quote}\small\textbf{OLP-0629 — स्रोतनोंद.} SOL6-A275: मूळातील सहा महिन्यांचा उल्लेख शिक्षेचा आहे; प्रत्यक्ष सुटका १४ सप्टेंबर १९१८ रोजी झाली. तसेच आयुष्यभर बिनशर्त शांततावाद हा दावा मागे घेतला: मॅकमास्टरच्या 30909 या पत्रनोंदीतील रसेल यांच्या स्वतःच्या विधानात काही परिस्थितींमध्ये बलप्रयोग समर्थनीय मानला आहे. स्रोत: https://russell-letters.mcmaster.ca/chronology आणि https://bracers.mcmaster.ca/30909.\end{quote}
\begin{quote}\small\textbf{OLP-0631 — स्रोतनोंद.} SOL6-A276: GCHQच्या नोंदीनुसार जोन क्लार्क आणि ट्यूरिंग यांचा साखरपुडा १९४१ मध्ये झाला. बॉम्ब यंत्र विद्युत् आणि यांत्रिक तंत्रावर चालत होते, असे The National Museum of Computing स्पष्ट करते; ते इलेक्ट्रॉनिक कोलॉससपेक्षा वेगळे होते. स्रोत: https://www.gchq.gov.uk/information/joan-clarke आणि https://www.tnmoc.org/bombe.\end{quote}
\begin{quote}\small\textbf{OLP-0632 — स्रोतनोंद.} SOL6-A277: झुरिकमधील १९१० चे पद वेतन मिळणारे होते; त्यामुळे त्यापूर्वीची गॉटिंगेनमधील १९०५ ची प्राध्यापक नेमणूक नाकारली जात नाही. फ्रायबुर्गची मानद नेमणूक १९२६ मध्ये झाली. हिटलरला अभिवादन करण्यास नकार दिल्यामुळे तक्रार झाल्यानंतर १९३५ मध्ये राजीनामा द्यावा लागला, अशी विद्यापीठाची नोंद आहे. स्रोत: https://uni-freiburg.de/math/geschichte-i/ आणि https://home.mathematik.uni-freiburg.de/mildenberger/freiburg2014/zermelo.html.\end{quote}
\begin{quote}\small\textbf{OLP-0635 — स्रोतनोंद.} OLINC-314: स्रोताच्या सीमा-व्याख्येत x=c वगळलेले नाही; मराठी सूत्रात 0<|x-c| ही आवश्यक अट घातली आहे.\end{quote}
\begin{quote}\small\textbf{OLP-0638 — स्रोतनोंद.} SOL6-A281: एकच द्विमान निरूपण-पद्धत वापरून एकास-एक फलनाच्या सिद्धतेतील अंक-एकमेवपणाची अट पूर्ण केली आहे. मूळ प्रतिमेबाहेरील संख्येचे उदाहरण [0,1] वर प्रतिमेतच मिळते; खरे प्रतिउदाहरण 5/6 दिले आहे. हा आधुनिक द्विमान युक्तिवाद मूळ दशमान पत्रातील शब्दशः सिद्धता म्हणून मांडलेला नाही.\end{quote}
\begin{quote}\small\textbf{OLP-0639 — स्रोतनोंद.} SOL6-A282: मूळ चित्रे जपली आहेत. समान प्राचल-अंतरालांवरील जाळी-घरांची सुसंगत विभागणी स्पष्ट केल्याने सीमा समसमान असल्याचे आणि प्रतिमा पूर्ण चौरस असल्याचे सिद्ध होते; सातत्याची अट प्रत्येक आदान बिंदूवर पूर्ण केली आहे.\end{quote}
\begin{quote}\small\textbf{OLP-0639 — स्रोतनोंद.} OLINC-317: स्रोताने x चा प्रांत चौरस दिला; येथे एकक रेषा हा योग्य प्रांत जपला आहे.\end{quote}
\begin{quote}\small\textbf{OLP-0639 — स्रोतनोंद.} OLINC-318: मूळ बिंदुनिहाय-सीमा युक्तिवादातील उणीव SOL6-A282 मध्ये दुरुस्त केली: समान प्राचल-अंतरालांवरील घर-अंतर्भाव, समसमान सीमा आणि अंतर्भूत बंद अंतरालांनी प्रतिमेतील प्रत्येक बिंदूची साक्ष मिळते.\end{quote}
\begin{quote}\small\textbf{OLP-0639 — स्रोतनोंद.} OLINC-319: प्रत्येक p साठी कुठलातरी खुला अंतराल देणे ही सातत्याची अट नव्हती. SOL6-A282 मध्ये प्रत्येक दिलेल्या x0 वरील epsilon-delta अट समसमान अभिसरणाने सिद्ध केली आहे.\end{quote}
\begin{quote}\small\textbf{OLP-0639 — स्रोतनोंद.} OLINC-320: मूळ शेवटच्या परिच्छेदातील असंबद्ध (a,b) उल्लेख आणि 'show to show' पुनरुक्ती नवीन सिद्धतेत आलेली नाहीत.\end{quote}
\begin{quote}\small\textbf{OLP-0655 — स्रोतनोंद.} SOL6-A294: उलटवण्यात प्रत्येक कोटीच्या कटांची संख्या वाढत नाही, ही सिद्ध सीमा वापरली आहे. अणु-कटातील विगमन कमी उंचीच्या तत्काळ उपसिद्धतेवर लावले आहे. कमी कोटीचे कट असतानाच्या दुर्बलीकरण व आकुंचनाचा आवश्यक आधार स्पष्ट केला आहे; सरावाचे प्रसंग सोडवलेले नाहीत.\end{quote}
\begin{quote}\small\textbf{OLP-0655 — स्रोतनोंद.} OLINC-339: स्रोताचे Delta=Delta,A हे चक्रीय समीकरण; आवश्यक विभागणी Delta=Delta',A अशी दुरुस्त केली आहे.\end{quote}
\begin{quote}\small\textbf{OLP-0655 — स्रोतनोंद.} OLINC-340: स्रोत येथे सर्वोच्च कटासाठीचे वेगळे उपप्रमेय दाखवतो; या परिच्छेदातील कोटी-कमीकरण मागील उपप्रमेयावर आधारित आहे.\end{quote}
\begin{quote}\small\textbf{OLP-0656 — स्रोतनोंद.} SOL6-A295: CutCS फक्त कट-सूत्र काढतो आणि उरलेला सामायिक संदर्भ जपतो. अस्तित्व-सूत्राचा संदर्भ पहिल्या वृक्षात जपला आहे; अभिव्यंजनाच्या वृक्षात आवश्यक दुर्बलीकरणाने दोन्ही आधारांचा संदर्भ जुळवला आहे आणि खालच्या कटचा योग्य निष्कर्ष दिला आहे.\end{quote}
\begin{quote}\small\textbf{OLP-0656 — स्रोतनोंद.} OLINC-341: स्रोत दुसऱ्या प्रसंगात pi\_2 लिहितो; उजवा आधार स्वयंसिद्धक असल्यास आवश्यक सिद्धता डाव्या आधाराची pi\_1 असते.\end{quote}
\begin{quote}\small\textbf{OLP-0656 — स्रोतनोंद.} OLINC-342: तिसऱ्या वृक्षात स्रोताचे Delta हे Delta' हवे; अन्यथा संदर्भातील B-and-C ची अतिरिक्त प्रत येते.\end{quote}
\begin{quote}\small\textbf{OLP-0656 — स्रोतनोंद.} OLINC-343: अंतिम वृक्षात Delta' आणि B-and-C च्या दोन प्रती हव्यात; स्रोताच्या Delta मुळे तीन प्रती येतात.\end{quote}
\begin{quote}\small\textbf{OLP-0658 — स्रोतनोंद.} SOL6-A296: चर-साक्षीचा प्रसंग भाषा-सीमेवरून पूर्ण केला आहे. सामान्य अंतर्वेशक सूत्ररूपात असतो; बंद संदर्भांतील उपयोगांसाठी मुक्त चरे सार्वत्रिकपणे बांधून वाक्यरूप घेतले आहे. OLINC-344 मधील ऐतिहासिक missing-case नोंद मूळ स्रोतासाठी कायम आहे, सध्याचा प्रसंग दुरुस्त आहे.\end{quote}
\begin{quote}\small\textbf{OLP-0658 — स्रोतनोंद.} SOL6-A303: बेथच्या परिभाष्यता परिच्छेदात उरलेल्या इंग्रजी predicate ऐवजी मराठी विधेय-संज्ञा वापरली आहे.\end{quote}
\begin{quote}\small\textbf{OLP-0659 — स्रोतनोंद.} SOL6-A297: सामायिक संदर्भ असलेल्या कटच्या वृक्षात CutCS नियमलेबल वापरले आहे.\end{quote}
\begin{quote}\small\textbf{OLP-0661 — स्रोतनोंद.} SOL6-A298: अग्र-संख्यापक सूत्राची matrix संख्यापकविरहित असल्याचे स्पष्ट केले आहे. आकुंचन-उलटवण्याने क्रममान वाढत नाही; अनुमाने काढल्यास ते कमी होऊ शकते.\end{quote}
\begin{quote}\small\textbf{OLP-0662 — स्रोतनोंद.} SOL6-A300: कलम-संयोजनापूर्वी बाहेरच्या सिद्धतेच्या मुक्ती-खुणा आत बसवलेल्या सिद्धतेच्या उघड्या गृहीतकांच्या खुणांपासून वेगळ्या ठेवायच्या; अन्यथा समान सूत्र असले तरी गृहीतक अनवधानाने मुक्त होऊ शकते.\end{quote}
\begin{quote}\small\textbf{OLP-0665 — स्रोतनोंद.} SOL6-A301: या प्रकरणाच्या बंद पदे वापरण्याच्या संकेताशी आदेशनाची अट जुळवली आहे. एका अस्वच्छ आयगेन-चर अनुमानाची दुरुस्ती केल्यावर अस्वच्छ अनुमानांची संख्या किमान एकाने घटते; नेमकी एकानेच घटते असा दावा नाही.\end{quote}
\begin{quote}\small\textbf{OLP-0671 — स्रोतनोंद.} SOL6-A302: अभिजात असत्यता नियमातून अंतिम सूत्र असत्य असल्यास आणखी कट वापरून रिकामा उत्तरपक्ष मिळवला आहे; कोणतेही गृहीतक मुक्त न होणाऱ्या प्रसंगाचाही अपेक्षित निष्कर्ष स्पष्ट केला आहे.\end{quote}
\begin{quote}\small\textbf{OLP-0672 — स्रोतनोंद.} SOL6-A307: उपसूत्राच्या गुणधर्माची परिचयातील सूचना पुढील अचूक निष्कर्षाशी जुळवली: उघडी गृहीतके, संख्यापकांची पद-उदाहरणे आणि आवश्यक असत्य समाविष्ट आहेत.\end{quote}
\begin{quote}\small\textbf{OLP-0675 — स्रोतनोंद.} SOL6-A304: क्रमपरिवर्तनापूर्वी बद्ध गृहीतक-खुणा सुरक्षित नवी नावे द्यायची; उपसिद्धतेच्या प्रत्येक प्रतीतील आयगेन स्थिरांक स्वतंत्र ठेवून नियमितता जपायची.\end{quote}
\begin{quote}\small\textbf{OLP-0676 — स्रोतनोंद.} SOL6-A305: न्यूनीकरणात सुरक्षित खुणा व आयगेन स्थिरांक वापरायचे. नव्या उघड्या गृहीतकांचा संच मूळच्या संचात समाविष्ट असतो; तो नेहमी तंतोतंत समान असेल किंवा खूण x पूर्णपणे नाहीशी होईल असा दावा नाही.\end{quote}
\begin{quote}\small\textbf{OLP-0676 — स्रोतनोंद.} SOL6-A309: नव्या परिच्छेदातील सिद्धता-नावे योग्य ग्रीक मुद्रण-संकेताने पुनर्स्थापित केली आहेत.\end{quote}
\begin{quote}\small\textbf{OLP-0678 — स्रोतनोंद.} SOL6-A306: मुख्य शाखेतील नवे गृहीतक नवी खूण वापरते. असत्य निष्कर्षासाठी रिकामा उत्तरपक्ष जपला आहे. उपसूत्राच्या गुणधर्मात उघडी गृहीतके, संख्यापकांची उदाहरणे आणि असत्य यांचा नेमका अपवाद स्पष्ट केला आहे.\end{quote}
\begin{quote}\small\textbf{OLP-0678 — स्रोतनोंद.} SOL6-A308: संकुचनाचा नियम योग्य मुद्रण-संकेताने दिला आहे. असत्य आधारासाठी मूळ रिकामा उत्तरपक्ष आणि दुर्बलीकरणाने मिळणारा आधार यांना स्वतंत्र सिद्धता-नावे दिली आहेत.\end{quote}
\begin{quote}\small\textbf{OLP-0678 — स्रोतनोंद.} SOL6-A310: नव्या अटींतील सूत्रे, सिद्धता-नावे आणि नियमांची नावे योग्य TeX संकेतांनी पुनर्स्थापित केली आहेत.\end{quote}
\begin{quote}\small\textbf{OLP-0679 — स्रोतनोंद.} SOL6-A311: सिद्धता न सापडता होणारे अपयशी थांबणे वेगळे केले आहे. न्याय्यतेत अणुरूप नसलेल्या आस्थितींचे निर्देशांक वापरले आहेत; वाक्यांच्या क्रमवर्तींसाठीची बंद-पद व्याप्ती स्पष्ट आहे.\end{quote}
\begin{quote}\small\textbf{OLP-0680 — स्रोतनोंद.} SOL6-A314: पुढील पद-प्रतिमानावर आधारित संपूर्णतेची चर्चा वाक्यांच्या क्रमवर्तींसाठी आहे, ही व्याप्ती परिचयात स्पष्ट केली आहे.\end{quote}
\begin{quote}\small\textbf{OLP-0683 — स्रोतनोंद.} SOL6-A312: सर्व योग्य पदे वापरून झाल्यावर संबंधित अरिक्त संचातील स्थिरांक पुन्हा वापरायचा; संचाबाहेरचा स्थिरांक आणून पुढील आयगेन-चर अट मोडायची नाही.\end{quote}
\begin{quote}\small\textbf{OLP-0684 — स्रोतनोंद.} SOL6-A316: सत्य-चिन्हांकित असत्यानेही शाखा बंद होते. पुन्हा वापरायच्या संख्यापकांसाठी अरिक्त स्थिरांक-संच, नवे उदाहरण नसतानाची कृती आणि संतृप्त उघड्या शाखेची अट स्पष्ट केली आहे.\end{quote}
\begin{quote}\small\textbf{OLP-0687 — स्रोतनोंद.} SOL6-A318: कच्च्या पद-विन्यासाऐवजी योग्य प्रकाराच्या पदांची व्याप्ती स्पष्ट केला. अपूर्ण रेडेक्स-मोजणी युक्तिवादाऐवजी आधीच्या नैसर्गिक निगमन सामान्यरूपीकरण प्रमेयाचा आणि सर्व अनुरूप रूपांतरणांचा वापर केला; बीटा-प्रबळ सामान्यीकरण हा स्वतंत्र अधिक बलवान परिणाम आहे.\end{quote}
\begin{quote}\small\textbf{OLP-0687 — स्रोतनोंद.} SOL6-A326: बेरीज-प्रकाराच्या उदाहरणांत तक्त्यातील injection चिन्ह वापरले आणि त्याची वरची खूण दुसरा प्रकार सांगते हे स्पष्ट केले.\end{quote}
\begin{quote}\small\textbf{OLP-0688 — स्रोतनोंद.} SOL6-A323: खुणित संदर्भ सुसंगत संच आहे आणि पदांची रचना विगमनाने आहे, हे आधीच्या नियमांशी जुळवले.\end{quote}
\begin{quote}\small\textbf{OLP-0689 — स्रोतनोंद.} SOL6-A324: येथे दिलेल्या पद-नियमांसाठी सिद्धतेची विधानीय व्याप्ती स्पष्ट केला; संख्यापकांचे पद-नियम दिले असल्याचा दावा नाही.\end{quote}
\begin{quote}\small\textbf{OLP-0691 — स्रोतनोंद.} SOL6-A321: मुख्य बीटा-नियम आणि पूर्ण सामान्यरूपीकरणातील क्रमपरिवर्तन व असत्य-सरलीकरण वेगळे केले; आत हलवलेल्या पदांचे मुक्त चल सुरक्षित ठेवले.\end{quote}
\begin{quote}\small\textbf{OLP-0693 — स्रोतनोंद.} SOL6-A320: संदर्भात खुणित प्रकार असल्यावर क्रमवर्ती स्वयंसिद्धक मिळते, अशी नियमाची दिशा दुरुस्त केली.\end{quote}
\begin{quote}\small\textbf{OLP-0696 — स्रोतनोंद.} SOL6-A322: प्रकार-संरक्षण संदर्भ वाढवता येणाऱ्या प्रकार-पद्धतीत आहे; अचूक वापरलेली गृहीतके घटू शकतात. क्रमपरिवर्तनाला बीटा-पायरी मानलेले नाही.\end{quote}
\begin{quote}\small\textbf{OLP-0696 — स्रोतनोंद.} SOL6-A327: पद-रूपांतरणे या शब्दाशी अनुरूप मराठी लिंग-सुसंगती दुरुस्त केली.\end{quote}
\begin{quote}\small\textbf{OLP-0697 — स्रोतनोंद.} SOL6-A319: प्रत्येक पदाला प्रकार असतो हा असत्य दावा काढला: योग्य प्रकार मिळाल्यास तो एकमेव असतो. सर्व संदर्भ-खुणांची सुसंगती आणि बद्ध चलांची सुरक्षित नावे स्पष्ट केली.\end{quote}
\begin{quote}\small\textbf{OLP-0698 — स्रोतनोंद.} SOL6-A328: दुर्बलीकरणापूर्वी आयगेन स्थिरांकांची सुरक्षित नावे घेतली; उंची-शून्य प्रसंगात डावीकडील असत्य-स्वयंसिद्धकही समाविष्ट केले.\end{quote}
\begin{quote}\small\textbf{OLP-0699 — स्रोतनोंद.} SOL6-A329: दोन-नियमांच्या विधानात असत्य स्थिरांकाचा अपवाद स्पष्ट केला; दुर्बलीकरणाची आवश्यकता दाखवणाऱ्या उदाहरणाला योग्य अणुरूप व नियम-संच scope दिला.\end{quote}
\begin{quote}\small\textbf{OLP-0700 — स्रोतनोंद.} SOL6-A330: दोन्ही क्रमवर्ती तक्त्यांतील असत्याचे स्वतंत्र स्वयंसिद्धक प्रस्तावनेतही स्पष्ट केले; बहुसंचाचे सर्व अर्थ जपले.\end{quote}
\begin{quote}\small\textbf{OLP-0701 — स्रोतनोंद.} SOL6-A331: व्युत्क्रमण व आकुंचनाच्या आधारप्रसंगांत असत्याचे स्वयंसिद्धक समाविष्ट केले; प्रतिस्थापनापूर्वी अंतर्गत आयगेन स्थिरांक सुरक्षित ठेवले. उंची आणि आधारविधानांची संख्या वेगळ्या अक्षरांनी लिहिली.\end{quote}
\begin{quote}\small\textbf{OLP-0702 — स्रोतनोंद.} SOL6-A332: संयोग-परिचयाच्या निष्कर्षात दोन पूर्वांग प्रती ठेवून पुढील आकुंचन खरे केले; खोलीच्या आणि स्वयंसिद्धकाच्या दाव्यांचा योग्य scope व निष्फळ संख्यापकाचा अपवाद स्पष्ट केला.\end{quote}
\begin{quote}\small\textbf{OLP-0703 — स्रोतनोंद.} SOL6-A313: क्रमवर्ती संख्यापक-नियमांच्या बंद पदे वापरण्याच्या संकेताशी आदेशनाची अट जुळवली आहे; आधीची नियमितीकरण आणि विगमनातील घट जपली आहे.\end{quote}
\begin{quote}\small\textbf{OLP-0705 — स्रोतनोंद.} SOL6-A333: लघुतम पद्धतीच्या व्याख्येत तक्त्यातील नेमके असत्याचे स्वयंसिद्धक आणि उजवे दुर्बलीकरण वगळले आहे.\end{quote}
\begin{quote}\small\textbf{OLP-0708 — स्रोतनोंद.} SOL6-A334: G1m ची व्याख्या G1i च्या तक्त्याशी जुळवली; उजवे दुर्बलीकरण आणि असत्याचे स्वयंसिद्धक दोन्ही वगळावे लागतात.\end{quote}
\begin{quote}\small\textbf{OLP-0710 — स्रोतनोंद.} SOL6-A335: तक्त्यात परिभाषित न केलेल्या mG1i आणि mG1m बद्दलचे मूळ विधान वगळले; दिलेले mG3i नियम जपले.\end{quote}
\begin{quote}\small\textbf{OLP-0711 — स्रोतनोंद.} SOL6-A336: सिद्धतेच्या व्याख्येतील विगमनाची संज्ञा सुसंगत केली आणि उघडे अवतरणचिन्ह बंद केले.\end{quote}
\begin{quote}\small\textbf{OLP-0713 — स्रोतनोंद.} SOL6-A337: G1c ते G3c भाषांतरातील उलट दिशेचा संदर्भ बदलला; दोन्ही दिशांत मुख्य संख्यापकीय सूत्र जपणाऱ्या नियमांचे आवश्यक दुर्बलीकरण किंवा आकुंचन स्पष्ट केले.\end{quote}


या नोंदी अनुवादकाने वेगळ्या जोडल्या आहेत; त्या मूळ मजकुराचा भाग नाहीत.
मूळ इंग्रजी स्रोताचे बाइट्स बदललेले नाहीत. मराठी TeX फाइलांतील नोंदवलेल्या स्रोतदुरुस्त्या या संपादकीय नोंदी आणि निर्णयनोंदवहीत स्पष्ट केल्या आहेत.
\begin{enumerate}
\item विभाग \ref{sfr:rel:set:sec} मधील $K=L\cup I$ आणि $H=G\cup I$ येथे
$I$ चे स्वतंत्र नामकरण राहिले आहे. आधीच्या कर्णाच्या चर्चेनुसार त्याचा अर्थ
$\Id{\Nat}$ हा एकरूपता संबंध असा घ्यावा.
\item विभाग \ref{sfr:rel:tre:sec} मधील शाखेच्या व्याख्येत $z\in X\setminus B$
असे छापले आहे. वृक्षाचा आधारसंच $A$ आहे; येथे $X$ ऐवजी $A$ अभिप्रेत दिसतो.
\item त्याच विभागातील आरंभीच्या खंडांनुसार बंद असणाऱ्या उपवृक्षाच्या उदाहरणात
$A$ अरिक्त असणे आवश्यक आहे: आधीच्या व्याख्येनुसार वृक्षाला मूळ असते,
म्हणून रिक्त संच त्या व्याख्येत बसत नाही.
\item फलन प्रकरणातील OLFUN-001 ते OLFUN-005 या पाच
गोठवलेल्या-स्रोत निष्कर्षांची दुरुस्ती संबंधित परिच्छेदांलगत स्वतंत्र
``स्रोतदुरुस्ती'' नोंदींमध्ये दिली आहे. मूळ इंग्रजी बाइट्स बदललेले नाहीत.
\item संचांच्या आकारमानाच्या प्रकरणातील OLSIZ-001 ते OLSIZ-010
या दहा सामायिक गोठवलेल्या-स्रोत निष्कर्षांची आणि MRSIZ-001 ते MRSIZ-004
या चार स्थानिक निरीक्षणांची नोंद संबंधित परिच्छेदांलगत दिली आहे. मूळ
इंग्रजी बाइट्स बदललेले नाहीत.
\item सामायिक OLSIZ-011 सूचना अचूक बाइट-पुनर्तपासणीनंतर चुकीची ठरून
मागे घेण्यात आली; तिच्यावर आधारित कोणतीही मराठी दुरुस्ती केलेली नाही.
\item अंकगणितीकरण प्रकरणातील MRARITH-001 ते MRARITH-012 या बारा
गोठवलेल्या-स्रोत निरीक्षणांतील दुरुस्त्या किंवा स्पष्ट खुलासे संबंधित
परिच्छेदांलगत स्वतंत्र ``स्रोतदुरुस्ती'' नोंदींमध्ये दिले आहेत. मूळ
इंग्रजी बाइट्स बदललेले नाहीत.
\item अनंत संच प्रकरणातील MRINF-001 ते MRINF-003 या तीन स्रोत-निरीक्षणांची
नोंद ठेवली आहे. MRINF-001 मधील व्याकरणदोषाचा अभिप्रेत अर्थ थेट दिला आहे;
MRINF-002 मधील अप्रकट आधारसंचाचे गृहीतक तज्ज्ञ-पुनरावलोकनासाठी राखले आहे;
MRINF-003 मधील स्रोतदोषाचा दावा पुनर्तपासणीनंतर चुकीचा ठरून मागे घेतला आहे.
अंतःस्थ \texttt{cardeq} रचना वैध आहे; मराठी निष्कर्ष सममूल्य केंद्रित
मांडणी म्हणून जतन केला आहे. मूळ इंग्रजी बाइट्स बदललेले नाहीत.
\item विधानीय तर्कशास्त्र प्रकरणातील MRPL-001 आणि MRPL-002 या दोन
गोठवलेल्या-स्रोत चिन्हदोषांच्या दुरुस्त्या संबंधित परिच्छेदांलगत स्वतंत्र
``स्रोतदुरुस्ती'' नोंदींमध्ये दिल्या आहेत. मूळ इंग्रजी बाइट्स बदललेले नाहीत.
\item \textbf{MRPRF-001}: गोठवलेल्या इंग्रजी टॅब्लो उदाहरणातील
शेवटच्या दोन पायऱ्या $\sFmla{\True}{\varphi \land \psi}$ चे घटक उलगडतात,
पण त्यांची नियम-खूण चुकून $\TRule{\True}{\lif}$ अशी आहे. या वाचकात ती
अर्थानुसार $\TRule{\True}{\land}$ केली आहे; संरेखित इंग्रजी आणि मराठी
स्रोत-बाइट्स बदललेले नाहीत.

\item \textbf{OLSEQ-001}: OLP-0075 मधील चार पायऱ्यांत पूर्वपक्षातील
सूत्रांची अदलाबदल असूनही नियम-खूण उजवी अदलाबदल अशी आहे. या वाचकात त्या
चार खुणा डावी अदलाबदल अशा केल्या आहेत; संरेखित स्रोतांत मूळ रूप जतन आहे.
\item \textbf{OLSEQ-002}: त्याच उदाहरणातील दोन कथनसूत्रांत दुसऱ्या
सूत्रापूर्वीचे नकारचिन्ह गाळले आहे. लगतचे ध्येय, सममितीचे स्पष्टीकरण आणि
सिद्धता-वृक्ष यांनुसार या वाचकात ती दोन नकारचिन्हे पुनःस्थापित केली आहेत.
\item \textbf{OLSEQ-003}: OLP-0072 च्या शेवटच्या परिच्छेदातील ``पदावर
निर्बंध नाही'' याचा अर्थ पद बंद असण्याव्यतिरिक्त आणखी निर्बंध नाही असा आहे;
मराठी मजकूर ही आधी स्पष्ट केलेली अट सांगतो.
\item \textbf{OLND-001}: OLP-0087 मध्येही संख्यापक-नियमांसाठी पद बंद
असण्याची उभी अट शेवटच्या स्पष्टीकरणात स्पष्ट केली आहे.
\item \textbf{OLND-002}: OLP-0087 मधील आयगेन चलाचा एकत्रित सारांश सर्व
आधारविधानांतून चल वगळल्यासारखा वाचतो; मराठी सारांश प्रत्येक नियमाची योग्य
आधारविधान, निष्कर्ष आणि अमुक्त गृहीतकांची अट वेगवेगळी राखतो.
\item \textbf{OLND-003}: OLP-0089 मध्ये आकृतीतील उजवी शाखा चुकून डावी
म्हटली आहे; वाचकात आकृती आणि पुढील परिच्छेदाशी सुसंगत ``उजवीकडील'' आहे.
\item \textbf{OLND-004}: OLP-0089 मधील एकच अनुमान कधी नकार-विलोपन तर
कधी असत्यता-प्रवेशन असे खूणांकित आहे. अनुमान अर्थतः समान असल्याने मूळची
पर्यायी खूण जतन केली आहे.
\item \textbf{OLND-005}: OLP-0090 मधील आयगेन-चल तपासणीत मुख्य
आधारविधानातील नकार गाळला आहे आणि नियमाची व्याप्ती अपुरी सांगितली आहे;
मराठी मजकूर योग्य नकार व तात्पुरत्या गृहीतकाचा अपवाद स्पष्ट करतो.
\item \textbf{OLND-006}: OLP-0090 मधील सार्वत्रिक विलोपनासाठी ``कोणतेही
पद'' म्हणजे उभ्या नियमानुसार कोणतेही बंद पद; मराठी स्पष्टीकरण ती अट राखते.
\item \textbf{OLND-007}: OLP-0095 च्या एका सरावातील नियम-खूण
\verb|\Elim{\forall}| ही प्रकरणात इतरत्र असलेल्या
\verb|\Elim{\lforall}| शी दृश्यतः समतुल्य आहे. संरेखित रूप जतन केले आहे.

\item \textbf{OLTAB-001--OLTAB-004}: टॅब्लो प्रकरणातील आधीच्या
उदाहरणांमध्ये विषय, बंद पदांची अट आणि एका ओळ-संदर्भाशी संबंधित स्रोतदोष
दुरुस्त किंवा स्पष्ट केले आहेत; संरेखित इंग्रजी आणि मराठी स्रोत-बाइट्स जतन आहेत.
\item \textbf{OLTAB-005--OLTAB-009}: टॅब्लो सिद्धतायोग्यता आणि सुसंगतता
विभागातील संच-लेखन, नियम-पूर्वपद, ओळ-संदर्भ आणि शाखा-वर्णनातील निश्चित
स्रोतदोषांचे मराठी रूपात मर्यादित निराकरण केले आहे; मूळ इंग्रजी बाइट्स जतन आहेत.
\item \textbf{OLTAB-010}: विधानीय परिणामांच्या आठ टॅब्लो नोड्समधील
चिन्हांकित-सूत्र मॅक्रोची गहाळ युक्ती-सिमा पुनर्स्थापित केली आहे; संरेखित
स्रोतांत गोठवलेले रूप जतन आहे.
\item \textbf{OLTAB-011}: टॅब्लो निर्दोषता सिद्धतेतील सार्वत्रिक नियमांच्या
पुनरावृत्त B/A रूपबंध-विसंगतीत निष्कर्ष आणि अर्थविषयक पायऱ्यांशी सुसंगत
A-रूप वापरले आहे; गोठवलेले रूप संरेखित स्रोतांत जतन आहे.
\item \textbf{OLTAB-012}: एकरूपता-टॅब्लोच्या सममिती आणि संक्रमकता
स्पष्टीकरणांत लगतच्या वृक्षांशी विसंगत असलेली दोन कथनसूत्रे दुरुस्त केली आहेत;
औपचारिक वृक्ष अपरिवर्तित आहेत.

\item \textbf{OLAXD-001--OLAXD-009}: स्वयंसिद्धकीय निष्पत्ती प्रकरणातील
कंस, अंतिम निष्कर्ष, स्वयंसिद्धक-संदर्भ, संख्यापक-नियम आणि जुनी
फाइल-माहिती यांतील नऊ निश्चित स्रोतदोष मराठी संरेखित स्रोतांत मर्यादितपणे
दुरुस्त केले आहेत; गोठवलेले इंग्रजी बाइट्स जतन आहेत.
\item \textbf{OLCOM-001--OLCOM-006}: संपूर्णता प्रकरणातील सूत्र-युक्तिवाद,
चल-सुसंगती, बंद-पद अट, विरामचिन्हे, प्रतिनिधी आणि टॅग-संदर्भ यांतील सहा
निश्चित स्रोतदोषांची दुरुस्ती किंवा स्पष्टता मराठी स्रोतांत नोंदवली आहे.
 \item \textbf{OLFOL-001--OLFOL-040}: प्रथम-क्रम तर्कशास्त्र आणि प्रतिमान
उपपत्तीच्या पुढील प्रकरणांतील एकोणचाळीस सूत्रात्मक, निर्देशांकविषयक,
अर्थविषयक आणि सिद्धताविषयक स्रोतदोष किंवा अपूर्ण पायऱ्या मर्यादितपणे
दुरुस्त किंवा स्पष्ट केल्या आहेत. OLFOL-009 मधील अतिरिक्त आकुंचित कंसाची
सूचना पुनर्तपासणीनंतर चुकीची ठरून मागे घेतली आहे; गोठवलेली रचना आधीपासून
संतुलित आहे. प्रत्येक नोंदीचा अचूक स्रोत-स्थान, पर्याय आणि
पुनरावलोकन-प्रश्न सोबतच्या तज्ज्ञ-पुनरावलोकन नोंदींत आहे.

\item \textbf{OLCMP-001--OLCMP-024}:\par संगणनक्षमता प्रकरणांतील औपचारिक
मांडणी, चिन्हांकने आणि सिद्धतापायऱ्यांबद्दलच्या स्रोत-निरीक्षणांची संरेखित
स्रोतांत स्वतंत्र नोंद आहे. मूळ इंग्रजी बाइट्स जतन आहेत; मराठी वाचक-मजकूर
नोंदवलेल्या मर्यादित स्पष्टता आणि दुरुस्त्या वापरतो.

\item \textbf{OLTUR-001--OLTUR-007}: ट्यूरिंग यंत्रांच्या उदाहरणांतील
आरंभीची अवस्था, फितीचे टोक, रिकामे आदान, प्रदान आणि यंत्र-संयोजन
यांबद्दलची सात स्रोत-निरीक्षणे संरेखित स्रोतांत नोंदवली आहेत. मूळ
इंग्रजी बाइट्स जतन आहेत; मराठी वाचक-मजकूर मर्यादित स्पष्टता वापरतो.

\item \textbf{OLTUR-008--OLTUR-022}: अनिर्णेयता आणि सांत प्रतिरूपांच्या
प्रकरणातील स्रोत-निरीक्षणे संबंधित मराठी मजकुरालगत दिली आहेत. विशेषतः
यंत्र-चालक्रमाचे निरूपण, शेवटच्या पायरीचे स्थितिवर्णन आणि सांत प्रतिरूपाची
मर्यादा यांतील दुरुस्त्या मूळ इंग्रजी मजकूर न बदलता नोंदवल्या आहेत.

\item \textbf{OLINC-001--OLINC-005}: अपूर्णतेच्या परिचयातील उद्दिष्टांची
क्रमवारी, प्रमेयांच्या गृहीतकांची व्याप्ती आणि संबंध-निरूपणातील उपपत्तीचा
निर्देश यांसंबंधीची निरीक्षणे नोंदवली आहेत. विकर्ण सिद्धतेतील दोन
वगळलेले निर्देशांक आणि विस्तारित उपपत्तींच्या सुसंगततेची अट मराठीत
स्पष्ट केली आहे; मूळ इंग्रजी मजकूर बदललेला नाही.

\item \textbf{OLINC-006--OLINC-043}: विन्यासमीमांसेच्या
अंकगणितीकरणात आणि $Q$ मधील निरूपणीयतेत मूळ इंग्रजीतील
सूत्रे, संख्यांकांचे निर्देश, सिद्धतापायऱ्या व गृहीतके
यांबद्दलची स्रोत-निरीक्षणे स्वतंत्र नोंदवली आहेत.
संबंधित मराठी विभागांत त्यांची मर्यादित दुरुस्ती
किंवा स्पष्टता दिली आहे; मूळ इंग्रजी बाइट्स बदललेले नाहीत.

\item \textbf{OLINC-044--OLINC-054}: उपपत्ती आणि संगणनक्षमता या
प्रकरणातील मूळ इंग्रजीतील सूत्रे, गृहीतकांची व्याप्ती आणि
प्रतिमानातील सत्य यांबद्दलची स्रोत-निरीक्षणे स्वतंत्र नोंदवली
आहेत. संबंधित मराठी विभागांत त्यांची मर्यादित दुरुस्ती
किंवा स्पष्टता दिली आहे; मूळ इंग्रजी बाइट्स बदललेले नाहीत.

\item \textbf{OLINC-055--OLINC-059}: अपूर्णता आणि सिद्धतायोग्यता
या प्रकरणातील मूळ इंग्रजीत वस्तुभाषेतील सिद्धतायोग्यता
विधेय, गोडेल-वाक्याचा संकेतनांक आणि सुसंगततेचे सूत्र
यांमध्ये आढळलेल्या विसंगती स्वतंत्र नोंदवल्या आहेत.
मराठी आवृत्तीत संबंधित सूत्रांची मर्यादित दुरुस्ती
केली आहे; मूळ इंग्रजी बाइट्स बदललेले नाहीत.

\item \textbf{OLINC-060--OLINC-061}: द्वितीय-क्रम चिन्हार्थमीमांसेच्या
मूळ इंग्रजीत संक्रमणाविषयी वापरलेले $R^*$ हे चिन्ह आधीच्या प्रकरणातील
व्याख्येशी विसंगत आहे; मराठी आवृत्तीत सकारात्मक संक्रमणासाठी $R^+$
ठेवले आहे. मर्यादित संचाच्या गणनेविषयीच्या सिद्धतेत शेवटच्या घटकाचे
फलनमूल्य स्पष्ट केले आहे. दोन्ही मुद्दे स्वतंत्र स्रोत-नोंदींत नोंदले आहेत.

\item \textbf{OLINC-062--OLINC-065}: द्वितीय-क्रम अंकगणितात
बेरीज परिभाषित करणाऱ्या सूत्रातील बांधलेला चल सुसंगत केला आहे.
स्वयंसिद्धकीकरणाच्या सिद्धतेतील तुटलेली गणित-सीमा दुरुस्त केली आहे.
संहततेच्या प्रमेयाला स्वतंत्र खूण दिली आहे आणि सांत उपसंचावरील
कमाल बंधाचा निर्देश त्या सांत उपसंचाकडेच केला आहे.
मूळ इंग्रजी बाइट्स बदललेले नाहीत.

\item \textbf{OLINC-066--OLINC-072}: द्वितीय-क्रम
संचसिद्धान्तावरील मूळ प्रकरण स्वतःच निष्कर्षांत दोष
असू शकतात अशी सूचना देते. उपसंचांच्या अनंतता,
गणनीयता आणि अलेफ-एक यांसाठीची काही मूळ सूत्रे
अपेक्षित गुणधर्म व्यक्त करत नव्हती. मराठी सूत्रांत
\textbf{SOL6-A096--SOL6-A098} या दुरुस्त्या केल्या आहेत:
एकास-एक अट संबंधित उपसंचापुरती घेतली आहे; अनंतता,
गणनीयता आणि अलेफ-एक यांच्या अटी दुरुस्त केल्या आहेत;
सांतत्याचे सांकेतीकरण आणि पूर्ण प्रांताचे आकारमान
त्यांनुसार दुरुस्त केले आहे. निवडीसह संचसिद्धान्ताचा
आधी स्पष्ट केलेला संदर्भ लागू आहे. मूळ इंग्रजी बाइट्स
आणि ऐतिहासिक दोषनोंदी जतन आहेत.

\item \textbf{OLINC-073--OLINC-081}: लॅम्डा कलनाच्या
परिचयात एका विभागाचा जुना प्रकरण-संकेत, करीकरणाच्या
अंतिम आदेशनातील चुकलेले पद, फलनाच्या स्थानसंख्येतील
विसंगती, आणि आदिम पुनरावर्तन व लघुतमीकरणाच्या
पुराव्यांतील स्थानिक त्रुटी नोंदवल्या आहेत. मराठी
मजकुरातील नेमक्या दुरुस्त्या व स्पष्टीकरणे स्रोत-नोंदीत
दर्शवली आहेत; मूळ इंग्रजी बाइट्स बदललेले नाहीत.

\item \textbf{OLINC-082--OLINC-102}: लॅम्डा पदांच्या
विन्यासमीमांसेतील एकमेव वाचनीयता, मुक्त चर, आदेशन,
अल्फा-परिवर्तन, सममूल्यतावर्ग आणि ईटा नियम यांमधील
गोठवलेल्या इंग्रजी स्रोताच्या स्थानिक त्रुटी स्वतंत्र
स्रोत-नोंदीत नोंदवल्या आहेत. गणिती अर्थ स्पष्ट करणाऱ्या
मराठी दुरुस्त्या तिथे नेमक्या ओळखता येतात; मूळ इंग्रजी
बाइट्स बदललेले नाहीत. अल्फा-परिवर्तनाचा सरावातील
दोन सारख्या प्रश्नजोड्या स्रोताप्रमाणे राखल्या आहेत.

\item \textbf{OLINC-103--OLINC-113}: चर्च–रॉसर प्रकरणातील
गोठवलेल्या इंग्रजी स्रोताच्या स्थानिक सूत्रचिन्ह, नियम,
चल आणि सिद्धता-अंतराच्या नोंदी स्वतंत्र स्रोत-नोंदीत
आहेत. दुरुस्त्या मराठी आवृत्तीत स्पष्ट केल्या आहेत;
बीटा–ईटा पूर्ण विकासातील आच्छादन-प्रकरणांची उणीव
\textbf{SOL6-A126} मध्ये दुरुस्त केली आहे: मुख्य
दावा आणि साहाय्यक दावा एकत्रित रचनात्मक विगमनाने
सिद्ध केले आहेत. मूळ इंग्रजी बाइट्स
बदललेले नाहीत.

\item \textbf{OLINC-114--OLINC-128}: लॅम्डा-परिभाष्यता
प्रकरणातील सूत्रे, आदानांची संख्या, न्यूनीकरणाच्या पायऱ्या,
पुनरावर्तन आणि स्थिरबिंदू संयोजकांच्या नोंदी स्वतंत्र
स्रोत-नोंदीत आहेत. जिथे गणिती अर्थ स्पष्ट आहे तिथे
मराठीतील स्थानिक दुरुस्त्या नोंदवल्या आहेत.
क्रमगुणिताच्या स्व-आदेशानंतरच्या सूत्रातील बाह्य
लॅम्डाच्या व्याप्तीचा दोष \textbf{SOL6-A132} मध्ये
दुरुस्त केला आहे: दोन्ही शाखा एकाच बाह्य शरीरात
आहेत; उरलेला Fac स्वसंदर्भ जपला आहे.
मूळ इंग्रजी बाइट्स बदललेले नाहीत.

\item \textbf{OLINC-129--OLINC-131}: बहुमूल्य तर्कशास्त्रातील
अभिजात तुलनेच्या गोठवलेल्या इंग्रजी प्रमेयाची व्याप्ती
अतिविस्तृत आहे आणि अंतिम सिद्धतेत पूर्तीऐवजी निष्पन्नतेची
चिन्हे चुकून वापरली आहेत. मराठी आवृत्तीत विधान सामाईक
चार-संयोजक खंडापुरते स्पष्ट केले आहे, सर्व विधानीय चलांसाठी
अभिजात मूल्यांची अट नमूद केली आहे आणि प्रत्युदाहरणातील
पूर्तीची चिन्हे दुरुस्त केली आहेत. मूळ इंग्रजी बाइट्स बदललेले नाहीत.

\item \textbf{OLINC-132--OLINC-137}: त्रिमूल्य तर्कशास्त्रांच्या
गोठवलेल्या इंग्रजी मजकुरात संधीच्या समानतेतील पुनरावृत्ती,
एका सरावातील अतिरिक्त कंस, विगमनाच्या आधारातील अटीविना
दिलेली समानता आणि संधीच्या पुराव्यात दुसऱ्या घटकाऐवजी
पहिलाच घटक पुन्हा लिहिलेला आहे. मराठी आवृत्तीत स्पष्ट
गणिती दुरुस्त्या नोंदवून केल्या आहेत. R-Mingle च्या
मॅट्रिक्समधील चुकीचा अभिव्यंजन-तक्ता आणि अतिरिक्त
असत्य-स्थिरांकाची अस्पष्ट व्याप्ती \textbf{SOL6-A142}
मध्ये दुरुस्त केली आहे. योग्य RM3 तक्ता आणि स्पष्ट
असत्य-स्थिरांक विस्तार दिला आहे. मूळ इंग्रजी
बाइट्स बदललेले नाहीत.

\item \textbf{OLINC-138--OLINC-142}: अनंतमूल्य तर्कशास्त्रांच्या
गोठवलेल्या इंग्रजी मजकुरात $V_\infty$ च्या संचवर्णनातील
शून्य हर, $V_m$ च्या संचवर्णनातील गणनाभेद आणि G\"odel
सत्यता-फलनातील चुकीची गणिती सीमाचिन्हे दुरुस्त केली आहेत.
अनंत गृहीतकांसाठी ससीम-मूल्य निष्पन्नतेवरून अनंतमूल्य
निष्पन्नता सर्वसाधारणपणे मिळत नाही; म्हणून केवळ ही
प्रतिदिशा ससीम गृहीतकांसाठी मांडली आहे. मूळ
इंग्रजी बाइट्स बदललेले नाहीत.

\item \textbf{OLINC-143--OLINC-145}: क्रमवर्ती कलनाच्या
गोठवलेल्या इंग्रजी मजकुरात डाव्या बाजूच्या सूत्रांची संख्या
क्रमवर्ती आणि संबंधित संधी यांत विसंगत आहे; एका
मूल्यनात मूल्यांकनाचा निर्देशांक गाळला आहे; आणि
$n$-बाजूंच्या क्रमवर्तीच्या प्रत्येक घटकाऐवजी फक्त
पहिल्या घटकाचा निर्देशांक लिहिला आहे. मराठी आवृत्तीत
या तीन गणिती नोंदी दुरुस्त केल्या आहेत. मूळ इंग्रजी
बाइट्स बदललेले नाहीत.

\item \textbf{OLINC-146--OLINC-153}: सामान्य मोडल
तर्कशास्त्रातील जगनिरपेक्ष अनिवार्यता, परिकर्मीची
कंसजोडी, द्विपक्षीय अभिव्यंजनाचा निवडक संकेत,
पूर्तीतील जगाचा निर्देशांक, आदेशन-दृष्टांतांची विगमन
सिद्धता आणि प्रतिमानाचे क्रमित त्रिक या ठिकाणी
गोठवलेल्या इंग्रजी स्रोताच्या मर्यादित दुरुस्त्या
मराठी मजकुरात नोंदल्या आहेत. मूळ इंग्रजी बाइट्स
बदललेले नाहीत.

\item \textbf{OLINC-154--OLINC-164}: चौकट-स्तरावरील
अनुरूपता, प्रतिउदाहरणांतील प्राप्यता-अटी, पाच रूपबंधांची
व्यत्यास सिद्धता, संहततेच्या साखळीचे निर्देशांक आणि
मोडल सूत्रांचे मानक भाषांतर यांत गोठवलेल्या इंग्रजी
स्रोताच्या मर्यादित दुरुस्त्या मराठी मजकुरात नोंदल्या
आहेत. मूळ इंग्रजी बाइट्स बदललेले नाहीत.

\item \textbf{OLINC-165--OLINC-174}: मोडल प्रणालींची
सिद्धता, व्युत्पन्न नियम, प्रतिमानांद्वारे वेगळेपणा
आणि सुसंगततेचे पुरावे यांतील गोठवलेल्या इंग्रजी
स्रोताच्या मर्यादित दुरुस्त्या मराठी मजकुरात नोंदल्या
आहेत. मूळ इंग्रजी बाइट्स बदललेले नाहीत.

\item \textbf{OLINC-175--OLINC-180}: संपूर्ण सुसंगत संचांचे
पुरावे, लिंडेनबाउमच्या प्रगणनातील लांबीची अट,
मोडल निष्पत्तींतील निर्देशांक आणि \(\Sigma\) ची खूण,
तसेच सत्यता पूर्वप्रमेयातील सराव-टॅग यांतील
गोठवलेल्या इंग्रजी स्रोताच्या मर्यादित दुरुस्त्या
मराठी मजकुरात नोंदल्या आहेत. मूळ इंग्रजी बाइट्स
बदललेले नाहीत.

\item \textbf{OLINC-181--OLINC-190}: निस्यंदनांच्या
भागात तुल्यतावर्गांचे आकार, मूल्यांकनांचे प्रांत,
भागाकार-जगांच्या खुणा, \(S5\) साठी प्रतिमानांचा
वर्ग, प्रमेयाशी जुळणारा सिद्धतेचा क्रम आणि
युक्लिडी चित्रांतील आवश्यक स्व-बाण यांबाबत
गोठवलेल्या इंग्रजी स्रोताच्या मर्यादित दुरुस्त्या
मराठी मजकुरात नोंदल्या आहेत. मूळ इंग्रजी बाइट्स
बदललेले नाहीत.

\item \textbf{OLINC-191--OLINC-208}: मोडल टॅब्लोच्या
नियमांतील चिन्हे व पूर्वचिन्हे, निर्दोषता आणि
संपूर्णतेच्या पुराव्यांतील काही सूत्रे, तसेच
प्रतिदृष्टांत प्रतिमानांच्या उदाहरणांतील
चुकीच्या खुणा मराठीत दुरुस्त केल्या आहेत.
\textbf{OLINC-204} मध्ये नोंदलेली सांत
\(\Gamma\) वरून सर्वसाधारण \(\Gamma\) कडे
जाणाऱ्या स्रोत-पुराव्यातील उणीव
\textbf{SOL6-A180} मध्ये दुरुस्त केली आहे.
असांत प्रकरणासाठी सांत उपसंचांची सुसंगतता,
आधीचे लिंडेनबाउम आणि सत्यता पूर्वप्रमेय
वापरले आहेत. मूळ इंग्रजी बाइट्स बदललेले नाहीत.

\item \textbf{OLINC-209}: कालिक भाषेतील सूत्रनिर्मितीच्या
नियमात मूळ इंग्रजीत साधा \(F\) आला आहे; त्याच
विभागातील संकारकांची यादी आणि सत्यतेची व्याख्या
\(\Ftemp\) वापरतात. मराठी नियमात \(\Ftemp\)
ठेवून ही एक-चिन्हीय विसंगती दुरुस्त केली आहे.
मूळ इंग्रजी बाइट्स बदललेले नाहीत.

\item \textbf{OLINC-210}: मूळ इंग्रजी प्रस्तावनेत
\emph{Knowledge and Belief} या ग्रंथाच्या लेखकाचे
पहिले नाव ``Jaako'' असे आहे. विद्यापीठाच्या
ग्रंथसूचीनुसार योग्य स्पेलिंग ``Jaakko'' आहे;
मराठीत ``याक्को हिंटिक्का'' वापरले आहे.
\item \textbf{OLINC-211}: द्विअनुकरणाच्या ``पुढे'' आणि
``मागे'' या दोन्ही अटींमध्ये मूळ इंग्रजीत कर्त्यांचा
संच \(A\) आला आहे. प्रकरणात परिभाषित कर्तृचिन्हांचा
संच \(G\) असल्याने मराठी अटींमध्ये \(G\) वापरले आहे.
\item \textbf{OLINC-212}: सार्वजनिक घोषणा भाषेतील
संयोजकांच्या यादीत मूळ इंग्रजीत द्विसशर्त राहिला
आहे; त्याच विभागातील सूत्रनिर्मितीच्या नियमात तो
असल्याने मराठी यादीत समाविष्ट केला आहे.
\item \textbf{OLINC-213}: घोषणेच्या रिक्तपणे सत्य
असण्याच्या उदाहरणात मूळ इंग्रजीत उत्तरसूत्राचे
चिन्हांकन राहिले आहे. मराठीत व्याख्येनुसार
\([\varphi] \psi\) हे रूप वापरले आहे.
मूळ इंग्रजी बाइट्स बदललेले नाहीत.

\item \textbf{OLINC-214}: BHK उदाहरणातील फलनाच्या
फलप्रदेशात मूळ इंग्रजीत एका ठिकाणी \(C\) आले
आहे; सूत्रानुसार मराठीत \(\chi\) वापरले आहे.
\item \textbf{OLINC-215}: विकल्पयोगाच्या डाव्या
समावेशनात मूळ इंग्रजीत \(M_1\) चे प्रतिमान
\(\langle 1,M_2\rangle\) दिले आहे. योग्य
\(\langle 1,M_1\rangle\) मराठीत वापरले आहे.
\item \textbf{OLINC-216}: संयोगाच्या BHK स्पष्टीकरणात
मूळ इंग्रजीत \(\varphi_1\land \varphi_1\) आले आहे;
जोडीचा दुसरा घटक \(\varphi_2\) ची रचना असल्याने
मराठीत \(\varphi_1\land \varphi_2\) वापरले आहे.
मूळ इंग्रजी बाइट्स बदललेले नाहीत.

\item \textbf{OLINC-217}: मूळ इंग्रजीत
स्थानिक निष्पन्नतेच्या दाव्याच्या सिद्धतेत
सर्व जगांतील सत्यतेचे अधिक बळकट गृहीतक
वापरले आहे; मराठीत आधी स्थानिक दावा
सिद्ध करून मग त्यावरून सर्व जगांतील
दावा मिळवला आहे.
\item \textbf{OLINC-218}: मूळ इंग्रजीत एका
ठिकाणी मूळ संच \(X\) लाच संस्थिति म्हटले
आहे. व्याख्येनुसार \(\mathcal O\) ही संस्थिति
आणि तिच्यासह \(X\) हा सांस्थितिक अवकाश आहे;
मराठीत हा भेद ठेवला आहे.
\item \textbf{OLINC-219}: मूळ इंग्रजीत
निष्पन्नता आणि सशर्त विधानाच्या दोन
अटींमध्ये \(\subset\) चिन्ह आहे.
या चिन्हाचा अर्थ संकेतपद्धतीनुसार बदलतो;
येथे समानताही चालते हे स्पष्ट करण्यासाठी
मराठीत \(\subseteq\) वापरले आहे.
उचित उपसंचच अभिप्रेत असल्याचा जुना
दोष-दावा मागे घेतला आहे.
मूळ इंग्रजी बाइट्स बदललेले नाहीत.

\item \textbf{OLINC-220}: स्वयंसिद्धकीय निर्दोषतेच्या
गृहीतक-प्रसंगात मूळ इंग्रजीतील पूर्ति-सूत्राला
जादा \(\Gamma\) घटक आहे; मराठीत \(A_n\)
ची पूर्ति योग्य रूपात लिहिली आहे.
\item \textbf{OLINC-221--222}: नैसर्गिक
निगमनाच्या निर्दोषता-सिद्धतेत संधियोगाच्या
निष्कर्षात चुकीची सूत्रे, तसेच विकल्पयोगाच्या
प्रसंगात जगाचे स्थान आणि दोन एकघटक संच
चुकले आहेत; मराठीत संदर्भाशी जुळणारी
रूपे वापरली आहेत.
\item \textbf{OLINC-223--224}: लिंडेनबाउमच्या
रचनेतील निवडीसाठी परिमित आधीच्या
निर्देशांकांचा युक्तिवाद स्पष्ट केला आहे;
कॅनॉनिकल प्रतिमानातील एकस्वनिकतेसाठी
विगमनाचा विस्तार-लांबी हा निर्देशांक
स्पष्ट केला आहे.
\item \textbf{OLINC-225}: सत्यता पूर्वप्रमेयाच्या
प्रस्तावनेत प्रारंभीच्या \(\Delta\) ऐवजी
त्या जगाचा \(\Delta(\sigma)\) अभाज्य संच
उल्लेखिला आहे.
\item \textbf{OLINC-227}: आधीच्या मोडल
संपूर्णता प्रकरणात मूल्यांकन \(V^\Sigma\)
पुढील गद्य कंस चुकून गणितात गेला होता;
मराठीत कंस गणिताबाहेर ठेवला आहे.
\item \textbf{OLINC-226}: निर्णेयतेच्या
स्रोतसिद्धतेत जगांचे केवळ विधानचरांच्या
सत्यसंचांवर आधारित भागाकार सशर्त
सूत्राची सत्यता जपत नाही. मूळ दोन
जगे असंबद्ध असताना नव्या उपसंच-क्रमाने
त्यांत प्राप्यता निर्माण होऊ शकते.
स्रोतदोषाची नोंद जपली आहे; सध्याच्या
मराठी विभागात योग्य उपसूत्र-गाळणी व
तिची सिद्धता दिली आहे. प्रमेय बदललेले नाही.
मूळ इंग्रजी बाइट्स बदललेले नाहीत.

\item \textbf{OLINC-228--230}: अंतःप्रज्ञावादी
टॅब्लो प्रकरणातील चुकून आलेला मोडल उल्लेख,
सत्याचा प्राप्य जगांमध्ये टिकाव, आणि
सशर्त सूत्रासाठीच्या शाखांवरील सत्य-असत्य
चिन्हे मराठीत संदर्भाशी जुळवली आहेत.
\item \textbf{OLINC-231}: उदाहरणातील
दोन संधियोग-नियमांच्या कारणनिर्देशांत
ओळ ४ ऐवजी त्या सूत्राची ओळ ७ दिली आहे.
\item \textbf{OLINC-232--235}: निर्दोषतेच्या
स्रोतसिद्धतेतील प्रतिप्रतिमानाचे चिन्ह,
दोन नियमांचे निष्कर्ष-संच आणि अंतिम
तार्किक निष्पन्नता यांतील चुका
दुरुस्त केल्या आहेत. \textbf{OLINC-233}
मधील मनमानी पूर्वचिन्ह-संचाविषयीची
व्याख्यात्मक उणीव दुरुस्त केली आहे:
सर्व वापरलेल्या पूर्वचिन्ह-विस्तारांसाठी
प्राप्यता आवश्यक आहे. नवे साक्षी-पूर्वचिन्ह
किंवा त्याचा विस्तार आधी वापरलेला नसावा,
या अटीने जुने अर्थनिर्धारण जपले जाते.
इतर नियमांच्या सिद्धता मूळप्रमाणे
स्वाध्याय आहेत. मूळ इंग्रजी बाइट्स
बदललेले नाहीत.

\item \textbf{OLINC-236}: कठोर अभिव्यंजनाच्या
एक अवैध निष्पन्नता-दाव्यात मूळ इंग्रजीत
वास्तविक अभिव्यंजनाचे निषेधन आले आहे.
मराठीत कठोर अभिव्यंजनाचे निषेधन वापरले आहे.
\item \textbf{OLINC-237}: कठोर अभिव्यंजन
आवश्यकतेने सत्य किंवा असत्य ठरते या दाव्याला
S5 ची व्याप्ती मराठीत स्पष्ट केली आहे.
\item \textbf{OLINC-238}: स्टालनाकर आणि
लुईस यांना एकच अद्वितीय सर्वाधिक जवळचे जग
गृहीत धरणारी मूळ सत्यता-अट दुरुस्त केली आहे.
स्टालनाकरची निवडलेल्या जगाची अट आणि
लुईसची गोलक-अट वेगळ्या दिल्या आहेत.
लुईसच्या विश्लेषणात समसमान जवळची अनेक जगे
किंवा सर्वात जवळचे जग नसणेही शक्य आहे.
आधीची स्रोतनोंद हा भेद जपते. मूळ इंग्रजी
बाइट्स बदललेले नाहीत.

\item \textbf{OLINC-239}: गोलक प्रतिमानाच्या
सत्यता-अटीतील ‘‘एखाद्या गोलकात
अट पूर्ण झाली, तर त्यातील
सर्व लहान गोलकांतही होते’’
हा मूळ इंग्रजी दावा अति-विस्तृत
आहे. पूर्वांग सत्य असलेले जग
धारण करणाऱ्या लहान गोलकांपुरताच
तो मराठीत मर्यादित केला आहे.
\item \textbf{OLINC-240}: बाह्य अवकाशातील
आगकाडीच्या उदाहरणात मूळ गद्य
दोनदा ‘‘पेटवणे’’ म्हणते; उद्धृत
अनुमान आणि अपेक्षित प्रयोग
यांच्यानुसार मराठीत ‘‘ओढणे’’
वापरले आहे.
\item \textbf{OLINC-241}: संक्रमणकतेच्या
प्रतिमानात \(q\lif r\) सर्व
संबंधित जगांत सत्य असूनही
मूळ सूत्रात पूर्तिचिन्हावर
निषेध आहे. मराठीत तो
निषेध काढला आहे.
\item \textbf{OLINC-242}: प्रतिपरिवर्तनाच्या
प्रतिमानात मूळ इंग्रजीत
एका जगाभोवतीची गोलक-प्रणाली
संपूर्ण फलन \(O\) म्हणून,
आणि घोषित \(M_1\) प्रतिमान
पुढे \(M\) म्हणून आले आहे;
मराठीत \(O_w\) व \(M_1\)
अशी सुसंगत नावे दिली आहेत.
मूळ इंग्रजी बाइट्स बदललेले नाहीत.

\item \textbf{OLINC-243}: स्वसमावेशकतेवरील चर्चेत
रसेल संचाचे घटक वर्णन करताना मूळ इंग्रजीत
चुकून दुहेरी निषेध आला आहे. सूत्रानुसार
स्वतःचे सदस्य नसलेले संच अभिप्रेत आहेत;
मराठी गद्यात तेच स्पष्ट केले आहे.
\item \textbf{OLINC-244}: दुष्टचक्र तत्त्व
नाकारण्यानंतर मर्यादित संचरचनेची गरज येते,
असे मूळ गद्य उलटे म्हणते. उद्धृत तत्त्व
आणि पुढील युक्तिवादानुसार ते स्वीकारल्यावर
मर्यादा येते, असे मराठीत दुरुस्त केले आहे.
\end{enumerate}


\paragraph{क्रम-समरूपणाच्या पुराव्यातील मुद्रित चूक.} मूळ इंग्रजी आवृत्तीत क्रम-समरूपणाच्या एका पुराव्यात $B_{b_2}$ ला फलनाचा प्रांत म्हटले आहे; प्रदर्शित $f:A_{a_2}\to B_{b_2}$ नुसार तो सहप्रांत आहे. मराठी मजकुरात हा शब्द दुरुस्त केला आहे.

\paragraph{पुनरावर्तनातील आधारप्रकरण.} मूळ इंग्रजीत रिकाम्या प्रांताच्या फलनासाठी पुनरावर्तनाच्या व्याख्येत प्रकरण दिलेले नाही, तरी पुराव्यात त्या प्रकरणाचे मूल्य वापरले आहे. मराठी मांडणीत ते प्रकरण स्पष्ट केले आहे.
\paragraph{पायाभूततेच्या पुराव्यातील बंधित चल.} मूळ सूत्रात आधी निश्चित केलेल्या $B$ ऐवजी अपरिभाषित $b$ आले आहे; मराठी सूत्रात $B$ वापरले आहे.
\paragraph{प्रतांकावरील पुराव्यातील अट.} मूळ पुराव्यात $x\in V_\alpha$ मानून लगेच $x\notin V_\alpha$ असे म्हटले आहे. मराठी मजकुरात हा निष्कर्ष प्रतांक आणि टप्पा समान असल्याच्या प्रकरणापुरता स्पष्ट केला आहे.

\paragraph{प्रतिबिंबनाच्या पुराव्यातील सूत्रदुरुस्त्या.} मूळ इंग्रजीत साक्षीदार-टप्प्याच्या अटीत एक अतिरिक्त बंद कंस आणि $\bigcup_{m<\omega}S_n$ मध्ये निर्देशांक-विसंगती आहे; मराठी सूत्रे अनुक्रमे संतुलित कंस आणि $S_m$ वापरतात. प्रतिस्थापनाच्या प्रतिबिंबनाधारित पुराव्यातील प्रतिमा-संचाच्या अटीचा हरवलेला बंद कंसही पुरवला आहे.
\paragraph{दुर्बल प्रतिबिंबनाचे नाव आणि सापेक्षीकरण.} मूळ मजकूर दुर्बल प्रतिबिंबनाची व्याख्या करण्याआधी चुकून प्रतिस्थापनाचे दुर्बल रूप म्हणतो; मराठी मांडणीत प्रतिबिंबन म्हटले आहे. पुढील परिशिष्टात $(\exists N\,\psi(N))^M$ उलगडताना संक्रमकतेच्या अटीवरील चुकीचे $N$ अधिलेख $M$ केले आहे.
\paragraph{विभक्तीकरणाविषयी मूळ पुराव्यातील त्रुटी.} सांत विस्ताराच्या पुराव्यात आणि त्यानंतरच्या प्रश्नात प्रत्येक संक्रमक संच विभक्तीकरणाचे प्रत्येक सापेक्षीकृत उदाहरण पूर्ण करतो, असा दावा आहे. तो खोटा आहे: $X=3=\{0,1,2\}$ हा संक्रमक संच घ्या. $A=2\in X$ आणि $x=1$ निवडणारी अट घेतल्यास अपेक्षित $\{1\}\notin X$. मराठी पुराव्यात आता लघुतम संक्रमक व उपसंच-समावेशक प्रतिमान वापरले आहे; प्रश्नाच्या विधानातही आवश्यक उपसंच-समावेशकत्वाची अट दिली आहे. प्रश्नाची सिद्धता वाचकांसाठी ठेवली आहे.

\paragraph{क्रमसूचक बेरजेच्या रचनेतील तीन दुरुस्त्या.} मूळ इंग्रजीत संबंधाच्या घटकजोड्या $D$ मधून घेतल्यासारख्या लिहिल्या आहेत; त्या $D\times D$ मधून असतात. पुढे अनुक्रमक बेरीज उलगडताना दोन आधीच खूण केलेल्या घटकसंचांमध्ये पुन्हा विभक्त बेरीज लिहिली आहे; मराठीत साधे संघटन वापरले आहे. शून्याशी बेरीज करताना रिकाम्या घटकाऐवजी एकघटकी संच लिहिलेला आहे; तो रिक्त संच केला आहे.
\paragraph{प्रतांकाच्या सरावातील चिन्ह.} मूळ दुसऱ्या सरावात अपेक्षित प्रतांक आणि वरची मर्यादा यांच्यामध्ये समानतेचे चिन्ह सुटले आहे; मराठीत ते पुरवले आहे. एका निष्क्रिय टिप्पणीत विसंगत गणित-सीमाचिन्हही जोडले आहे.
\paragraph{क्रमसूचक गुणाकाराचे सीमा प्रकरण.} मूळ सार्वत्रिक सूत्रातील कठोर लघुतम ऊर्ध्वबंध शून्य डाव्या घटकासाठी चुकीचा परिणाम देतो. मराठी सूत्रात आधीच्या गुणाकारांचे संघटन वापरले आहे; त्यामुळे शून्याचे सीमा क्रमसूचक संख्येशी गुणाकार शून्य राहतो.
\paragraph{क्रमसूचक घातांकनाचा फलनसंच आणि शून्य आधार.} मूळ रचनेत आधार आणि घातांक यांचा फलनसंचातील क्रम उलटा आहे. मराठी व्याख्या आणि सरावात घातांकावरून आधाराकडे जाणारी, केवळ सांत ठिकाणी अशून्य असणारी फलने घेतली आहेत. शून्य आधार आणि $\omega$ घातांकासाठी मात्र मूळ सीमा-सूत्रात शून्य घातांकाचे मूल्य $1$ संघटित होते, तर फलनसंच रिकामा असतो. त्यामुळे मुख्य दोन्ही व्याख्या आणि सराव शून्येतर आधारापुरते केले आहेत; सरावाची सिद्धता दिलेली नाही.

\paragraph{सांततेच्या विधानातील अर्थदुरुस्ती.} मूळ इंग्रजीत $\card{A}\notin\omega$ या अटीचा अर्थ चुकून “$A$ नैसर्गिक संख्या नाही” असा दिला आहे. सांत पण नैसर्गिक संख्या नसलेले संच असतात; म्हणून मराठीत “$A$ सांत नाही” असे म्हटले आहे.
\paragraph{अनंत उत्तरवर्ती संख्येच्या पुराव्यातील मधला टप्पा.} मूळ पुरावा उत्तरवर्ती क्रमसूचक संख्या अनंत असल्यास तिची पूर्ववर्ती अनंत असते, यासाठी सांत क्रमसूचक संख्यांच्या तुल्यबलतेवरील प्रतिज्ञेचा संदर्भ देतो. तो संदर्भ एकटाच पुरेसा नाही. पूर्ववर्ती सांत असती तर तिची उत्तरवर्तीही सांत असती, हा आवश्यक टप्पा मराठीत स्पष्ट केला आहे.

\paragraph{संचांक घातांकनाची संकेतरचना.} घातसंचाच्या संचांकावरील पुराव्यात आणि घातांकनाच्या मर्यादेत मूळ स्रोत साधी वरलिखित घात-संकेतरचना वापरतो. त्या ठिकाणी मराठीत संचांक घातांकनाची स्पष्ट संकेतरचना वापरली आहे.
\paragraph{कॅनॉनिकल क्रमातील सांत प्रकरण.} आद्यखंडाच्या आकारमानाचा पुरावा मूळ स्रोत फक्त अनंत कमाल निर्देशांकासाठी स्पष्ट करतो. कमाल निर्देशांक सांत असेल तर आद्यखंड सांत असतो, हे आवश्यक प्रकरण मराठीत जोडले आहे.
\paragraph{घातांकनाच्या पुराव्यातील फलनांचे प्रकार.} वियुक्त युतीवरील फलनाचे विभाजन दोन फलनांची क्रमित जोडी देते; त्यांच्या आलेखांचा कार्तीय गुणाकार देत नाही. तसेच पुनर्गटित फलनाचा प्रत्यक्ष प्रांत कार्तीय गुणाकार आहे, त्याचा संचांक-प्रतिनिधी नाही. मराठीत हे दोन्ही प्रकार दुरुस्त केले आहेत.
\paragraph{शून्य घातांकाचा अपवाद.} अनंत संचांकाचा सांत घात स्वतः त्या संचांकाइतका असतो, या विधानात शून्य घातांक वगळावा लागतो. शून्य घाताचे मूल्य एक असते. GCH अंतर्गत पुढील घात-मर्यादेतही अनंत आधार आणि शून्येतर घातांक या आवश्यक अटी मराठीत स्पष्ट केल्या आहेत.
\paragraph{अलेफ अनुक्रमाच्या पुराव्यातील टप्पे.} प्रत्येक अनंत संचांक अलेफ अनुक्रमात येतो, या पुराव्यात लघुतम अनंत संचांकाचे आधारप्रकरण आणि अनंत संचांकांपुरते विगमन आवश्यक आहे. हे आवश्यक टप्पे मराठीत स्पष्ट केले आहेत. एकमेवतेविषयीचा वैध मूळ संक्षिप्त तर्क स्वतंत्र विस्तार न करता जपला आहे.
\paragraph{स्थिरबिंदू अनुक्रमातील वाढ.} मूळ अंतिम रचनेत प्रारंभिक संचांक आधीच बेथ-स्थिरबिंदू असेल तर पुनरावर्तन स्थिर राहते. त्यामुळे दावा केलेली काटेकोर वाढ आणि नंतरचे एकास-एक फलन मिळत नाही. मराठीत सुरुवात प्रारंभिक संचांकाच्या उत्तरवर्ती संचांकापासून केली आहे. पुढील अनुक्रमाचे शून्यवे मूल्य शून्य आहे; ते बेथ-स्थिरबिंदू नाही, हा अपवादही स्पष्ट केला आहे.


\paragraph{टार्स्की--स्कॉटच्या आक्षेपातील रिक्त संच.} तुल्यबल संचांचे सर्व प्रतिनिधी एकत्र केल्यास प्रतांकाची मर्यादा राहत नाही, हा आक्षेप अरिक्त आधारसंचासाठी आहे. रिक्त आधारसंचाचा वर्ग मात्र एकघटकी संच आहे. मराठीत हा अपवाद स्पष्ट केला आहे.
\paragraph{हार्टोग्सच्या पुराव्यातील संच आणि निर्देशांक.} मूळ पुराव्यात सुक्रमित आधारसंच हा संबंधाचा उपसंच असल्यासारखा छापला आहे; तो मूळ आधारसंचाचा उपसंच असतो. घूर्णनाशी संबंध नसलेल्या याच पुराव्यात फलनाने क्रम नेण्याच्या सूत्रात निश्चित क्रमसूचक संख्या आणि अंतर्गत निर्देशांक यांना एकच नाव दिले आहे. मराठीत वेगळी, प्रांतात बंधित नावे वापरली आहेत.
\paragraph{निवडीच्या पुराव्यातील थांबण्याची व्याख्या.} मूळ पुढील-मूल्य अटीतील असमानतेची दिशा आधीच्या निवडी पुसून टाकते. ती दुरुस्त करून, एकास-एक आणि आच्छादक फलनाचा दावा फक्त थांबण्यापूर्वीच्या अंशासाठी केला आहे. रिक्त आधारसंचाचे प्रकरण आधी सोडवले आहे; न थांबण्याचा विरोधाभास प्रत्यक्ष आधारसंचावरील हार्टोग्स निष्कर्षाने मिळतो.
\paragraph{गणनीय निवडीच्या मोजणीतील निर्देशांक.} आधीच्या उपसंचांच्या युतीकरणात मूळ स्रोत चालू निर्देशांक पुन्हा वापरतो. मराठीत युतीकरणाचा योग्य चल निर्देशांक वापरला आहे आणि शून्यवे प्रकरण वेगळे केले आहे. गणनीय युतीकरणाच्या पुराव्यात गणनीय निवड पुरेशी आहे; या विभागात उलटी सममूल्यता सिद्ध केलेली नाही, हेही स्पष्ट केले आहे.
\paragraph{निवडसंच आणि प्रतिमा.} प्रत्येक सदस्याशी एकघटकी छेद असण्याच्या अटीत बाहेरील अतिरिक्त घटकांना बंदी नाही. निवडफलनाची प्रतिमा म्हणून ओळखण्याआधी ते अतिरिक्त घटक वगळले आहेत.
\paragraph{भूमितीय विधानांचे प्रांत.} बानाख--टार्स्कीचे घनगोल येथे धन त्रिज्येचे आणि त्रिमितीय आहेत; वितालीची वर्तुळे धन त्रिज्येची आहेत. पृष्ठभागाच्या विघटनातून पृष्ठभागाच्या प्रती मिळतात, घन अंतर्भागाच्या नव्हेत. अंतरालावरील स्पर्शरेषा-फलनाच्या सूत्रातील अतिरिक्त बंद कंस काढला आहे.
\paragraph{वितालीच्या घूर्णन गटाची दुरुस्ती.} परिमेय संख्येइतक्या रेडियनांची घूर्णने पूर्ण फेरीच्या मर्यादेत घेतल्यास व्यस्त आणि संयोजनासाठी संवरण मिळत नाही. म्हणून पूर्ण फेरीच्या परिमेय भागांची घूर्णने घेतली आहेत. शून्य घूर्णनाचा व्यस्त स्वतंत्रपणे हाताळला आहे. दोन अर्ध्यांना स्वतंत्र आधार म्हणण्याऐवजी जनक संच म्हटले आहे; प्रत्येक अर्ध्यावरील कोन-दुप्पट प्रतिचित्रण एकास-एक आणि आच्छादक असते. विभाजनातील अपरिभाषित पहिल्या-अर्ध्याचे नावही योग्य गटाच्या नावाने बदलले आहे.
\paragraph{त्रिमितीय रूपरेषेतील उरलेली पुरावा-उणीव.} मूळ मजकूर प्रत्येक मुक्त गटासाठी दुप्पट विघटन सांगतो; येथे प्रत्यक्ष वापरलेले दोन जनकांवरील मुक्त गटाचे प्रकरण म्हटले आहे. पण त्रिमितीय घूर्णने त्यांच्या अक्षांवरील काही बिंदू स्थिर ठेवतात. वर्तुळावरील मुक्त क्रियेचा पुरावा त्या बिंदूंना थेट लागू होत नाही. त्यांचा स्वतंत्र विचार या रूपरेषेत दिलेला नाही; म्हणून ही पूर्ण पुरावा नसल्याची उणीव स्पष्ट ठेवली आहे.
\paragraph{मापाच्या पुराव्यातील क्षेत्र आणि चल.} माप ऋणेतर असल्याची आवश्यक अट आणि अमापनीयता ही निश्चित मापाच्या क्षेत्राविषयी असल्याचा अर्थ स्पष्ट केला आहे. शेवटच्या पुराव्यात घूर्णनाचा चल चुकून निवडलेल्या बिंदूंच्या संचात घेतला होता; तो घूर्णन गटात घेतला आहे.


\paragraph{सिद्धता-प्रकरणातील आवृत्ती आणि विभागक्रम.} मूळ इंग्रजी आवृत्तीतील नैसर्गिक भाषेचा उल्लेख या मराठी आवृत्तीसाठी स्पष्ट केला आहे. सिद्धतेचे प्रकार ताबडतोब पुढच्या विभागात येतात असे दोन उल्लेख आहेत; प्रत्यक्ष स्रोतक्रमात व्याख्यांचा विभाग आधी असल्यामुळे पुढील विभागांचा सामान्य निर्देश केला आहे. व्याख्यांच्या विभागाचा मूळ भाग-निर्देशांकही पद्धती भागाशी जुळवला आहे.
\paragraph{अस्तित्ववाची विधानातील नव्या नावाचे कार्य.} निष्कर्षात नाव नसणे एवढेच पुरेसे नाही: नव्या वस्तूविषयी अतिरिक्त, मुक्त न केलेली गृहीतके वापरता कामा नयेत. तात्पुरते सदस्यत्व गृहीतक मुक्त केल्यावर मूळ अस्तित्ववाची विधान आणि इतर अनुमत गृहीतके यांचा आधार उरतो. मूळ स्पष्टीकरणात अरिक्त संचाऐवजी बंधित चलाचे नाव आले होते; संचाचे नाव वापरले आहे.
\paragraph{संचांच्या उदाहरणांतील मुद्रणदोष.} एका संयोग-नॉनसमावेशनाच्या स्पष्टीकरणात अपरिभाषित संचाचे नाव आले होते; योग्य संयोग वापरला आहे. शोषणाच्या विस्तृत सिद्धतेत पहिला समावेशन-निष्कर्ष दोनदा दिला होता; दुसरा निष्कर्ष उलट समावेशनाने बदलला आहे. दोन उदाहरणांतील न जुळणारे कंस दुरुस्त केले आहेत.


\paragraph{नैसर्गिक संख्यांवरील विगमनाचे निर्देशांक.} मूळ विवेचनाच्या एका पायरीत $k$ हा चल असताना त्याच्या जागी $1$ ठेवण्याऐवजी दुसऱ्या चलाचे नाव आले आहे; योग्य पायरीचल वापरला आहे. फाशांच्या उदाहरणात शून्य फाशांचे प्रकरण प्रस्तावनेत असले, तरी औपचारिक विगमन $1$ पासून सुरू होते. शून्य फाशांच्या रिकाम्या बेरजेचे प्रकरण स्पष्टपणे जोडले आहे.
\paragraph{प्रबल विगमनातील रिक्त पूर्वभाग.} $k=0$ असताना $l<0$ अशी कोणतीही नैसर्गिक संख्या नसते. या रिक्त पूर्वभागातील मूळ मजकुरात बंधित $l$ ऐवजी $P(0)$ छापले आहे; विगमनाच्या प्रत्यक्ष गृहीतकानुसार $P(l)$ वापरले आहे.
\paragraph{उचित पूर्वखंडाच्या पुराव्याची व्याप्ती.} मूळ व्याख्येत रिकामी चिन्हमाला उचित पूर्वखंड ठरते, पण तीत उघडणारे आणि बंद करणारे कंस शून्य असतात. म्हणून कंसांची काटेकोर असमानता तिच्यासाठी असत्य आहे. मूळ पुराव्यातील पाचही प्रकरणे अरिक्त पूर्वखंडांचीच आहेत; मराठीत ती आवश्यक अट व्याख्या, विधान आणि पुराव्यात स्पष्ट केली आहे.


\paragraph{इतिहासातील चरित्रविषयक दुरुस्त्या.} चर्च यांच्या प्रबंधाशी संबंधित गणनीयतेचे प्रतिपादन आणि प्रथम-क्रम निर्णय समस्येवरील त्यांचे प्रमेय यांतील भेद स्पष्ट केला आहे. पेटर यांच्या जन्मनावातील उलटलेला उल्लेख, ज्युलिया रॉबिन्सन यांच्या नॅशनल अकॅडमीतील निवडीचा गणित विभागाशी संबंध, ट्यूरिंगवरील चित्रपटातील कलाकारांची नावे आणि झर्मेलो यांच्या गॉटिंगेनमधील अभ्यासस्थानाची वर्तनी तपासून दुरुस्त केली आहे.


\paragraph{सीमेची व्याख्या आणि चित्रातील निर्देशांक.} मूळ मर्यादा-व्याख्येत $x=c$ वगळलेले नाही; तसे केल्याशिवाय सामान्य सीमेची अट मिळत नाही. मराठी सूत्रात $0<|x-c|$ ही अट स्पष्ट केली आहे. निरपेक्ष मूल्याच्या आलेखात $-1$ या बिंदूचे चिन्ह दुरुस्त केले आहे. तसेच आधीच्या अवकलज-सूत्रातील बदलत्या भागाकाराचा उल्लेख स्पष्ट केला आहे. \paragraph{द्विमान विस्तार.} \textbf{SOL6-A281} मध्ये द्विमान निरूपणांची एकच निवड ठरवून एकास-एक फलनाची अट पूर्ण केली आहे. मूळ missing-value उदाहरण या प्रांतावर प्रतिमेतच येते; त्याऐवजी प्रतिमेबाहेरचे $5/6$ उदाहरण दिले आहे. ऐतिहासिक दशमान प्रयत्न आणि त्याचे आधुनिक द्विमान रूपांतर वेगळे केले आहेत. \paragraph{हिल्बर्ट वक्राच्या सिद्धतेची दुरुस्ती.} प्रांत एकक रेषा ठेवून \textbf{SOL6-A282} मध्ये प्रत्येक स्तराचे समान बंद प्राचल-अंतराल आणि जाळी-घरांचे सुसंगत refinement दिले आहे. घरकर्णाच्या मर्यादेने समसमान अभिसरण मिळते. अंतर्भूत बंद अंतराल प्रत्येक चौरस-बिंदूची प्रतिमापूर्व साक्ष देतात; प्रत्येक आदान बिंदूवरील सातत्य epsilon-delta अटीने सिद्ध केले आहे. मूळ तीन आकृत्या तशाच जपल्या आहेत.

\chapter*{संदर्भ}
\markboth{संदर्भ}{संदर्भ}

\addcontentsline{toc}{chapter}{संदर्भ}

खालील संदर्भ मूळ स्रोताच्या ग्रंथसूचीवरून घेतले आहेत. लेखकांची नावे, प्रकाशनांची मूळ शीर्षके आणि ग्रंथसूचीतील तपशील ओळखीसाठी मूळ भाषेत ठेवले आहेत. ही नोंद दुव्यांच्या सध्याच्या उपलब्धतेची पडताळणी नाही.


\par\medskip\phantomsection\label{bib:Turing1937}
Alan M. Turing (1937). \emph{On computable numbers, with an application to the ``{E}ntscheidungsproblem''}. \emph{Proceedings of the London Mathematical Society, 2nd Series}, 42, 230--265.

\par\medskip\phantomsection\label{bib:Tarski1981}
Alfred Tarski (1981). \emph{The Collected Works of {A}lfred {T}arski}. Birkh\"{a}user. Basel.

\par\medskip\phantomsection\label{bib:Church1936}
Alonzo Church (1936). \emph{A note on the {E}ntscheidungsproblem}. \emph{The Journal of Symbolic Logic}, 1, 40--41.

\par\medskip\phantomsection\label{bib:Church1936a}
Alonzo Church (1936). \emph{An Unsolvable Problem of Elementary Number Theory}. \emph{American Journal of Mathematics}, 58, 345--363.

\par\medskip\phantomsection\label{bib:Andrasfai1986}
Andr\'asfai, B\'ela (1986). \emph{R\'ozsa ({R}osa) {P}\'eter}. \emph{Periodica Polytechnica Electrical Engineering}, 30(2-3), 139--145. \url{http://www.pp.bme.hu/ee/article/view/4651}.

\par\medskip\phantomsection\label{bib:Aspray1984}
Aspray, William (1984). \emph{The {P}rinceton Mathematics Community in the 1930s: {A}lonzo {C}hurch}. Interview. \url{http://www.princeton.edu/mudd/finding_aids/mathoral/pmc05.htm}.

\par\medskip\phantomsection\label{bib:Baaz2011}
Baaz, Matthias; Papadimitriou, Christos H.; Putnam, Hilary W.; Scott, Dana S.; Harper Jr., Charles L. (2011). \emph{Kurt {G}\"odel and the Foundations of Mathematics: Horizons of Truth}. Cambridge University Press. Cambridge.

\par\medskip\phantomsection\label{bib:BanachTarski1924}
Banach, Stefan; Tarski, Alfred (1924). \emph{Sur la d\'{e}composition des ensembles de points en parties respectivement congruentes}. \emph{Fundamenta Mathematicae}, 6, 244--77.

\par\medskip\phantomsection\label{bib:Benacerraf1965}
Benacerraf, Paul (1965). \emph{What numbers could not be}. \emph{The Philosophical Review}, 74(1), 47--73.

\par\medskip\phantomsection\label{bib:Berkeley1734}
Berkeley, George (1734). \emph{The Analyst; or, a Discourse Adressed to an Infidel Mathematician}.

\par\medskip\phantomsection\label{bib:Russell1905}
Bertrand Russell (1905). \emph{On Denoting}. \emph{Mind}, 14, 479--493.

\par\medskip\phantomsection\label{bib:Russell1967}
Bertrand Russell (1967). \emph{The Autobiography of Bertrand Russell}. Allen and Unwin. London.

\par\medskip\phantomsection\label{bib:Russell1968}
Bertrand Russell (1968). \emph{The Autobiography of Bertrand Russell}. Allen and Unwin. London.

\par\medskip\phantomsection\label{bib:Russell1969}
Bertrand Russell (1969). \emph{The Autobiography of Bertrand Russell}. Allen and Unwin. London.

\par\medskip\phantomsection\label{bib:RussellND}
Bertrand Russell (n.d.). \emph{{B}ertrand {R}ussell on Smoking}. Video Interview. \url{https://www.youtube.com/watch?v=80oLTiVW_lc}.

\par\medskip\phantomsection\label{bib:Boolos1971}
Boolos, George (1971). \emph{The Iterative Conception of Set}. \emph{The Journal of Philosophy}, 68(8), 215--31.

\par\medskip\phantomsection\label{bib:Boolos1989}
Boolos, George (1989). \emph{Iteration Again}. \emph{Philosophical Topics}, 17(2), 5--21.

\par\medskip\phantomsection\label{bib:Boolos2000}
Boolos, George (2000). \emph{Must we Believe in Set Theory?}. मध्ये: \emph{Between Logic and Intuition: Essays in Honor of {Charles} {Parsons}}. पृ.~257--68. Cambridge University Press. Cambridge.

\par\medskip\phantomsection\label{bib:Burali-Forti1897}
Burali-Forti, Cesare (1897). \emph{Una questione sui numeri transfiniti}. \emph{Rendiconti del Circolo Matematico di Palermo}, 11, 154--64.

\par\medskip\phantomsection\label{bib:ButtonLT1}
Button, Tim (2021). \emph{Level Theory, Part~1: Axiomatizing the Bare Idea of a Cumulative Hierarchy of Sets}. \emph{The Bulletin of Symbolic Logic}, 27(4), 436-460.

\par\medskip\phantomsection\label{bib:Cantor1878}
Cantor, Georg (1878). \emph{Ein {Beitrag} zur {Mannigfaltigkeitslehre}}. \emph{Journal für die reine und angewandte Mathematik}, 84, 242--58.

\par\medskip\phantomsection\label{bib:Cantor1883}
Cantor, Georg (1883). \emph{Grundlagen einer allgemeinen {Mannigfaltigkeitslehre}. Ein mathematisch-philosophischer {Versuch} in der {Lehre} des {Unendlichen}}. Teubner. Leipzig.

\par\medskip\phantomsection\label{bib:Cantor1892}
Cantor, Georg (1892). \emph{{\"Uber} eine elementare {Frage} der {Mannigfaltigkeitslehre}}. \emph{Jahresbericht der deutschen Mathematiker-Vereinigung}, 1, 75--8. Teubner. Leipzig.

\par\medskip\phantomsection\label{bib:Cohen1963}
Cohen, Paul J. (1963). \emph{The independence of the continuum hypothesis}. \emph{Proceedings of the National Academy of Sciences of the United States of America}, 24, 556--557.

\par\medskip\phantomsection\label{bib:Cohen1966}
Cohen, Paul J. (1966). \emph{Set Theory and the Continuum Hypothesis}. Benjamin. Reading, MA.

\par\medskip\phantomsection\label{bib:Reid1996}
Constance Reid (1996). \emph{Julia: A Life in Mathematics}. Cambridge University Press. Cambridge. \url{https://books.google.ca/books?id=lRtSzQyHf9UC&lpg=PP1&pg=PP1#v=onepage&q&f=false}.

\par\medskip\phantomsection\label{bib:Conway2006}
Conway, John (2006). \emph{The Power of Mathematics}. मध्ये: \emph{Power}. Cambridge University Press. Cambridge. \url{http://www.cs.toronto.edu/~mackay/conway.pdf}.

\par\medskip\phantomsection\label{bib:Csicsery2016}
Csicsery, George (2016). \emph{Zala Films: {J}ulia {R}obinson and {H}ilbert's Tenth Problem}. \url{http://www.zalafilms.com/films/juliarobinson.html}.

\par\medskip\phantomsection\label{bib:Velleman2019}
Daniel J. Velleman (2019). \emph{How to Prove It: A Structured Approach}. 3rd आवृत्ती. Cambridge University Press. Cambridge.

\par\medskip\phantomsection\label{bib:Solow2013}
Daniel Solow (2013). \emph{How to Read and Do Proofs}. Wiley. Hoboken, NJ.

\par\medskip\phantomsection\label{bib:Dauben1990}
Dauben, Joseph (1990). \emph{{G}eorg {C}antor: {H}is {M}athematics and {P}hilosophy of the {I}nfinite}. Princeton University Press. Princeton.

\par\medskip\phantomsection\label{bib:Dedekind1888}
Dedekind, Richard (1888). \emph{Was sind und was sollen die {Zahlen}?}. Vieweg. Braunschweig.

\par\medskip\phantomsection\label{bib:Dick1981}
Dick, Auguste (1981). \emph{Emmy {N}oether 1882--1935}. Birkh\"auser. Boston.

\par\medskip\phantomsection\label{bib:Duncan2015}
Duncan, Arlene (2015). \emph{The {B}ertrand {R}ussell {R}esearch {C}entre}. \url{http://russell.mcmaster.ca/}.

\par\medskip\phantomsection\label{bib:Ebbinghaus2015}
Ebbinghaus, Heinz-Dieter (2015). \emph{Ernst {Z}ermelo: {A}n {A}pproach to his {L}ife and {W}ork}. Springer-Verlag. Berlin.

\par\medskip\phantomsection\label{bib:Ebbinghaus2010}
Ebbinghaus, Heinz-Dieter; Fraser, Craig G.; Kanamori, Akihiro (2010). \emph{Ernst Zermelo. Collected Works}. Springer-Verlag. Berlin.

\par\medskip\phantomsection\label{bib:Ebbinghaus2013}
Ebbinghaus, Heinz-Dieter; Kanamori, Akihiro (2013). \emph{Ernst Zermelo: Collected Works}. Springer-Verlag. Berlin.

\par\medskip\phantomsection\label{bib:Steinhart2018}
Eric Steinhart (2018). \emph{More Precisely: The Math You Need to Do Philosophy}. 2nd आवृत्ती. Broadview. Peterborough, ON.

\par\medskip\phantomsection\label{bib:Zermelo1904}
Ernst Zermelo (1904). \emph{Beweis, da{\ss} jede {M}enge wohlgeordnet werden kann}. \emph{Mathe\-ma\-tische An\-nalen}, 59, 514--516. {E}nglish translation in \hyperref[bib:Ebbinghaus2010]{Ebbinghaus 2010, pp.~115--119}.

\par\medskip\phantomsection\label{bib:Cheng2004}
Eugenia Cheng (2004). \emph{How to Write Proofs: A Quick Quide}. \url{https://eugeniacheng.com/wp-content/uploads/2017/02/cheng-proofguide.pdf}.

\par\medskip\phantomsection\label{bib:Feferman2004}
Feferman, Anita; Feferman, Solomon (2004). \emph{Alfred Tarski: Life and Logic}. Cambridge University Press. Cambridge.

\par\medskip\phantomsection\label{bib:Feferman1994}
Feferman, Solomon (1994). \emph{{J}ulia {B}owman {R}obinson 1919--1985}. \emph{Biographical Memoirs of the National Academy of Sciences}, 63, 1--28. National Academies Press. \url{http://www.nasonline.org/publications/biographical-memoirs/memoir-pdfs/robinson-julia.pdf}.

\par\medskip\phantomsection\label{bib:Godel1986}
Feferman, Solomon; Dawson Jr., John W.; Kleene, Stephen C.; Moore, Gregory H.; Solovay, Robert M.; van Heijenoort, Jean (1986). \emph{Kurt G\"odel: Collected Works. Vol. 1: Publications 1929--1936}. Oxford University Press. Oxford.

\par\medskip\phantomsection\label{bib:Godel1990}
Feferman, Solomon; Dawson Jr., John W.; Kleene, Stephen C.; Moore, Gregory H.; Solovay, Robert M.; van Heijenoort, Jean (1990). \emph{Kurt G\"odel: Collected Works. Vol. 2: Publications 1938--1974}. Oxford University Press. Oxford.

\par\medskip\phantomsection\label{bib:Fraenkel1922}
Fraenkel, Abraham (1922). \emph{\"{U}ber den {Begriff} `definit' und die {Unabh\"angigkeit} des {Auswahlaxioms}}. \emph{Sitzungsberichte der Preussischen Akadademie der Wissenschaften, Physikalisch-mathematische Klasse}, 253--257.

\par\medskip\phantomsection\label{bib:Frege1884}
Frege, Gottlob (1884). \emph{Die {Grundlagen} der {Arithmetik}: Eine logisch mathematische {Untersuchung} \"{u}ber den {Begriff} der {Zahl}}. Wilhelm Koebner. Breslau. Translation in \hyperref[bib:Frege1953]{Frege 1953}.

\par\medskip\phantomsection\label{bib:Frege1953}
Frege, Gottlob (1953). \emph{Foundations of Arithmetic}. 2nd आवृत्ती. Basil Blackwell \& Mott. Oxford.

\par\medskip\phantomsection\label{bib:Frey2015}
Frey, Holly; Wilson, Tracy V. (2015). \emph{Stuff you Missed in History Class: {E}mmy {N}oether, Mathematics Trailblazer}. Podcast audio. \url{https://www.iheart.com/podcast/stuff-you-missed-in-history-cl-21124503/episode/emmy-noether-mathematics-trailblazer-30207491/}.

\par\medskip\phantomsection\label{bib:Godel1938}
G\"odel, Kurt (1938). \emph{The consistency of the axiom of choice and the generalized continuum hypothesis}. \emph{Proceedings of the National Academy of Sciences of the United States of America}, 50, 1143--8.

\par\medskip\phantomsection\label{bib:Gentzen1935a}
Gerhard Gentzen (1935). \emph{Untersuchungen \"uber das logische {S}chlie\ss en {I}}. \emph{Mathe\-mati\-sche Zeit\-schrift}, 39, 176--210. {E}nglish translation in \hyperref[bib:Gentzen1969]{Szabo 1969}, pp.~68--131. \url{https://doi.org/10.1007/BF01201353}.

\par\medskip\phantomsection\label{bib:Gentzen1935b}
Gerhard Gentzen (1935). \emph{Untersuchungen \"uber das logische {S}chlie\ss en {II}}. \emph{Mathe\-mati\-sche Zeit\-schrift}, 39, 176--210, 405--431. {E}nglish translation in \hyperref[bib:Gentzen1969]{Szabo 1969}, pp.~68--131. \url{https://doi.org/10.1007/BF01201363}.

\par\medskip\phantomsection\label{bib:Giaquinto2007}
Giaquinto, Marcus (2007). \emph{Visual Thinking in Mathematics}. Oxford University Press. Oxford.

\par\medskip\phantomsection\label{bib:Gouvea2011}
Gouv\^{e}a, Fernando Q. (2011). \emph{Was {C}antor Surprised?}. \emph{American Mathematical Monthly}, 118(3), 198-209.

\par\medskip\phantomsection\label{bib:Grattan-Guinness1971}
Grattan-Guinness, Ivor (1971). \emph{Towards a Biography of {G}eorg {C}antor}. \emph{Annals of Science}, 27(4), 345--391.

\par\medskip\phantomsection\label{bib:Wagon2016}
Grzegorz {Tomkowicz}; Stan {Wagon} (2016). \emph{The {Banach}-{Tarski} Paradox}. Cambridge University Press. Cambridge.

\par\medskip\phantomsection\label{bib:Hartogs1915}
Hartogs, Friedrich (1915). \emph{\"{U}ber das {P}roblem der {W}ohlordnung}. \emph{Mathematische Annalen}, 76, 438--43.

\par\medskip\phantomsection\label{bib:Hausdorff1914}
Hausdorff, Felix (1914). \emph{Bemerkung über den {Inhalt} von {Punktmengen}}. \emph{Mathematische Annalen}, 75, 428--34.

\par\medskip\phantomsection\label{bib:Heijenoort1967}
Heijenoort, Jean van (1967). \emph{From {F}rege to {G}\"{o}del: A Source Book in Mathematical Logic, 1879--1931}. Harvard University Press. Cambridge, MA.

\par\medskip\phantomsection\label{bib:EndertonND}
Herbert B. Enderton (2019). \emph{{A}lonzo {C}hurch: {L}ife and {W}ork}. मध्ये: \emph{The Collected Works of {A}lonzo {C}hurch}. MIT Press. Cambridge, MA.

\par\medskip\phantomsection\label{bib:Hilbert1891}
Hilbert, David (1891). \emph{\"{U}ber die stetige {Abbildung} einer {Linie} auf ein {Fl\"{a}chenst\"{u}ck}}. \emph{Mathematische Annalen}, 38(3), 459--460.

\par\medskip\phantomsection\label{bib:EwaldSieg2013}
Hilbert, David (2013). \emph{{D}avid {H}ilbert's Lectures on the Foundations of Arithmetic and Logic 1917--1933}. Springer. Heidelberg.

\par\medskip\phantomsection\label{bib:Hodges2014}
Hodges, Andrew (2014). \emph{Alan Turing: The Enigma}. Vintage. London.

\par\medskip\phantomsection\label{bib:Hume1740}
Hume, David (1740). \emph{A Treatise of Human Nature}. London.

\par\medskip\phantomsection\label{bib:Incurvati2020}
Incurvati, Luca (2020). \emph{Conceptions of Set and the Foundations of Mathematics}. Cambridge University Press. Cambridge.

\par\medskip\phantomsection\label{bib:Irvine2015}
Irvine, Andrew David (2015). \emph{Sound Clips of {B}ertrand {R}ussell Speaking}. \url{http://plato.stanford.edu/entries/russell/russell-soundclips.html}.

\par\medskip\phantomsection\label{bib:Noether1983}
Jacobson, Nathan (1983). \emph{Emmy {N}oether: {G}esammelte {A}bhandlungen---{C}ollected {P}apers}. Springer-Verlag. Berlin.

\par\medskip\phantomsection\label{bib:Tarski1983}
John Corcoran (1983). \emph{Logic, Semantics, Metamathematics}. 2nd आवृत्ती. Hackett. Indianapolis.

\par\medskip\phantomsection\label{bib:Dawson1997}
John Dawson, Jr. (1997). \emph{Logical Dilemmas: The Life and Work of Kurt {G}\"odel}. CRC Press. Boca Raton.

\par\medskip\phantomsection\label{bib:Robinson1949}
Julia Robinson (1949). \emph{Definability and Decision Problems in Arithmetic}. \emph{The Journal of Symbolic Logic}, 14(2), 98-114. \url{http://www.jstor.org/stable/2266510}.

\par\medskip\phantomsection\label{bib:KatzKatz2012}
Katz, Karin Usadi; Katz, Mikhail G. (2012). \emph{{S}tevin Numbers and Reality}. \emph{Foundations of Science}, 17(2), 109--23.

\par\medskip\phantomsection\label{bib:Kunen1980}
Kunen, Kenneth (1980). \emph{Set Theory: An Introduction to Independence Proofs}. North Holland. New York.

\par\medskip\phantomsection\label{bib:Godel1929}
Kurt G\"odel (1929). \emph{{\"U}ber die {V}ollst\"andigkeit des {L}ogikkalk\"uls [{O}n the Completeness of the Calculus of Logic]}. Reprinted and translated in \hyperref[bib:Godel1986]{Feferman 1986}, pp.~60--101.

\par\medskip\phantomsection\label{bib:Godel1931}
Kurt G{\"{o}}del (1931). \emph{\"Uber formal unentscheidbare {S}\"atze der \emph{{P}rincipia {M}athematica} und verwandter {S}ysteme {I} [{O}n Formally Undecidable Propositions of \emph{{P}rincipia {M}athematica} and Related Systems~{I}]}. \emph{Monatshefte f\"ur Mathematik und Physik}, 38, 173--198. Reprinted and translated in \hyperref[bib:Godel1986]{Feferman 1986}, pp.~144--195. \url{https://doi.org/10.1007/BF01700692}.

\par\medskip\phantomsection\label{bib:Levy1960}
L\'{e}vy, Azriel (1960). \emph{Axiom Schemata of Strong Infinity in Axiomatic Set Theory}. \emph{Pacific Journal of Mathematics}, 10(1), 223--38.

\par\medskip\phantomsection\label{bib:LibriVoxND}
LibriVox (n.d.). \emph{Bertrand {R}ussell}. Collection of public domain audiobooks. \url{https://librivox.org/author/1508?primary_key=1508&search_category=author&search_page=1&search_form=get_results}.

\par\medskip\phantomsection\label{bib:Linnebo2010}
Linnebo, \O{}ystein (2010). \emph{Predicative and Impredicative Definitions}. \emph{Internet Encyclopedia of Philosophy}. \url{http://www.iep.utm.edu/predicat/}.

\par\medskip\phantomsection\label{bib:Linsenmayer2014}
Linsenmayer, Mark (2014). \emph{The Partially Examined Life: {G}\"odel on Math}. Podcast audio. \url{http://www.partiallyexaminedlife.com/2014/06/16/ep95-godel/}.

\par\medskip\phantomsection\label{bib:MacFarlane2015}
MacFarlane, John (2015). \emph{Alonzo {C}hurch's {JSL} Reviews}. \url{http://johnmacfarlane.net/church.html}.

\par\medskip\phantomsection\label{bib:Magnus2021}
Magnus, P. D.; Button, Tim; Loftis, J. Robert; {Thomas-Bolduc}, Aaron; Trueman, Robert; Zach, Richard (2021). \emph{Forall {\emph{x}}: {{Calgary}}. {{An}} Introduction to Formal Logic}. F21 आवृत्ती. {Open Logic Project}. {Calgary}. \url{https://forallx.openlogicproject.org/}.

\par\medskip\phantomsection\label{bib:Sautoy2014}
Marcus du Sautoy (2014). \emph{{A} Brief History of Mathematics: {G}eorg {C}antor}. Audio Recording. \url{http://www.bbc.co.uk/programmes/b00ss1j0}.

\par\medskip\phantomsection\label{bib:DavisPutnamRobinson1961}
Martin Davis; Hilary Putnam; Julia Robinson (1961). \emph{The Decision Problem for Exponential {D}iophantine Equations}. \emph{Annals of Mathematics}, 74(3), 425-436. \url{http://www.jstor.org/stable/1970289}.

\par\medskip\phantomsection\label{bib:Matijasevich1992}
Matijasevich, Yuri (1992). \emph{My Collaboration with {J}ulia {R}obinson}. \emph{The Mathematical Intelligencer}, 14(4), 38--45.

\par\medskip\phantomsection\label{bib:Menzler-Trott2007}
Menzler-Trott, Eckart (2007). \emph{Logic's Lost Genius: The Life of {G}erhard {G}entzen}. American Mathematical Society. Providence.

\par\medskip\phantomsection\label{bib:Hutchings2003}
Michael Hutchings (2003). \emph{Introduction to Mathematical Arguments}. \url{https://math.berkeley.edu/~hutching/teach/proofs.pdf}.

\par\medskip\phantomsection\label{bib:Montague1961}
Montague, Richard (1961). \emph{Semantic closure and non-finite axiomatizability {I}}. मध्ये: \emph{Infinitistic Methods: Proceedings of the Symposium on Foundations of Mathematics (Warsaw 1959)}. पृ.~45--69. Pergamon. New York.

\par\medskip\phantomsection\label{bib:Montague1965}
Montague, Richard (1965). \emph{Set theory and higher-order logic}. मध्ये: \emph{Formal systems and recursive functions}. पृ.~131--48. North-Holland. Amsterdam. Proceedings of the Eight Logic Colloquium, July 1963.

\par\medskip\phantomsection\label{bib:Imitation2014}
Morten Tyldum (2014). \emph{The Imitation Game}. Motion picture.

\par\medskip\phantomsection\label{bib:OConnorRobertson:RN}
O'Connor, John J.; Robertson, Edmund F. (2005). \emph{The real numbers: {Stevin} to {Hilbert}}. \url{http://www-history.mcs.st-and.ac.uk/HistTopics/Real_numbers_2.html}.

\par\medskip\phantomsection\label{bib:Oconnor2014}
O'Connor, John J.; Robertson, Edmund F. (2014). \emph{R\'{o}zsa {P}\'{e}ter}. \url{http://www-groups.dcs.st-and.ac.uk/~history/Biographies/Peter.html}.

\par\medskip\phantomsection\label{bib:Peter1935b}
P\'eter, R\'ozsa (1935b). \emph{Konstruktion nichtrekursiver {F}unktionen}. \emph{Mathematische Annalen}, 111, 42--60. \url{https://doi.org/10.1007/BF01472200}.

\par\medskip\phantomsection\label{bib:Peter1951}
P\'eter, R\'ozsa (1951). \emph{Rekursive Funktionen}. Akademiai Kiado. Budapest. English translation in \hyperref[bib:Peter1967]{P\'eter 1967}.

\par\medskip\phantomsection\label{bib:Peter1967}
P\'eter, R\'ozsa (1967). \emph{Recursive Functions}. Academic Press. New York.

\par\medskip\phantomsection\label{bib:Peter2010}
P\'eter, R\'ozsa (2010). \emph{Playing with Infinity}. Dover. New York. \url{https://books.google.ca/books?id=6V3wNs4uv_4C&lpg=PP1&ots=BkQZaHcR99&lr&pg=PP1#v=onepage&q&f=false}.

\par\medskip\phantomsection\label{bib:Peano1890}
Peano, Giuseppe (1890). \emph{Sur une courbe, qui remplit toute une aire plane}. \emph{Mathematische Annalen}, 36(1), 157--60.

\par\medskip\phantomsection\label{bib:Maddy1988a}
Penelope Maddy (1988). \emph{Believing the Axioms~{I}}. \emph{The Journal of Symbolic Logic}, 53(2), 481--511.

\par\medskip\phantomsection\label{bib:Maddy1988b}
Penelope Maddy (1988). \emph{Believing the Axioms~{II}}. \emph{The Journal of Symbolic Logic}, 53(3), 736--64.

\par\medskip\phantomsection\label{bib:Perimeter2015}
Perimeter Institute (2015). \emph{Emmy {N}oether: Her Life, Work, and Influence}. Video Lecture. \url{https://www.youtube.com/watch?v=tNNyAyMRsgE}.

\par\medskip\phantomsection\label{bib:Potter2004}
Potter, Michael (2004). \emph{Set Theory and its Philosophy}. Oxford University Press. Oxford.

\par\medskip\phantomsection\label{bib:Peter1935a}
P{\'e}ter, R{\'o}zsa (1935a). \emph{{\"U}ber den {Z}usammenhang der verschiedenen {B}egriffe der rekursiven {F}unktion}. \emph{Mathematische Annalen}, 110, 612--632. \url{https://doi.org/10.1007/BF01448046}.

\par\medskip\phantomsection\label{bib:Radiolab2012}
Radiolab (2012). \emph{The {T}uring Problem}. Podcast audio. \url{http://www.radiolab.org/story/193037-turing-problem/}.

\par\medskip\phantomsection\label{bib:Ramsey1925}
Ramsey, Frank Plumpton (1925). \emph{The Foundations of Mathematics}. \emph{Proceedings of the London Mathematical Society}, 25, 338--384.

\par\medskip\phantomsection\label{bib:Reid1986}
Reid, Constance (1986). \emph{The Autobiography of {J}ulia {R}obinson}. \emph{The College Mathematics Journal}, 17, 3--21.

\par\medskip\phantomsection\label{bib:Hammack2013}
Richard Hammack (2013). \emph{Book of Proof}. Virginia Commonwealth University. Richmond, VA. \url{http://www.people.vcu.edu/~rhammack/BookOfProof/BookOfProof.pdf}.

\par\medskip\phantomsection\label{bib:Heck2012}
Richard Kimberly Heck (2012). \emph{Reading Frege's Grundgesetze}. Oxford University Press. Oxford.

\par\medskip\phantomsection\label{bib:Robinson1996}
Robinson, Julia (1996). \emph{The Collected Works of Julia Robinson}. American Mathematical Society. Providence.

\par\medskip\phantomsection\label{bib:Robinson1947}
Robinson, Raphael (1947). \emph{On the decomposition of spheres}. \emph{Fundamenta Mathematicae}, 34(1), 246--60.

\par\medskip\phantomsection\label{bib:Rose2012}
Rose, Daniel (2012). \emph{A Song about {G}eorg {C}antor}. Audio Recording. \url{https://www.youtube.com/watch?v=QUP5Z4Fb5k4}.

\par\medskip\phantomsection\label{bib:Rose2010}
Rose, Nicholas J (2010). \emph{Hilbert-Type Space-Filling Curves}. \url{https://web.archive.org/web/20151010184939/http://www4.ncsu.edu/~njrose/pdfFiles/HilbertCurve.pdf}.

\par\medskip\phantomsection\label{bib:Russell1919}
Russell, Bertrand (1919). \emph{Introduction to Mathematical Philosophy}. Allen \& Unwin. London.

\par\medskip\phantomsection\label{bib:Scott1974}
Scott, Dana (1974). \emph{Axiomatizing Set Theory}. मध्ये: \emph{Axiomatic Set Theory II}. पृ.~207--14. American Mathematical Society. Proceedings of the Symposium in Pure Mathematics of the American Mathematical Society, July--August 1967.

\par\medskip\phantomsection\label{bib:Segal2014}
Segal, Sanford L. (2014). \emph{Mathematicians under the Nazis}. Princeton University Press. Princeton.

\par\medskip\phantomsection\label{bib:Shoenfield:AST}
Shoenfield, Joseph R. (1977). \emph{Axioms of Set Theory}. मध्ये: \emph{Handbook of Mathematical Logic}. पृ.~321--44. North-Holland. London.

\par\medskip\phantomsection\label{bib:Sigmund2007}
Sigmund, Karl; Dawson, John; M\"uhlberger, Kurt; Enzensberger, Hans Magnus; Kennedy, Juliette (2007). \emph{Kurt {G}\"odel: {D}as {A}lbum--{T}he {A}lbum}. \emph{The Mathematical Intelligencer}, 29(3), 73--76. Springer.

\par\medskip\phantomsection\label{bib:Skolem1922}
Skolem, Thoralf (1922). \emph{Einige {Bemerkungen} zur axiomatischen {Begr\"undung} der {Mengenlehre}}. मध्ये: \emph{Wissenschaftliche Vortr\"age gehalten auf dem f\"unften Kongress der skandanivschen Mathematiker in Helsingfors vom 4.\ bis zum 7.\ Juli 1922}. पृ.~137--52. Akademiska Bokhandeln.

\par\medskip\phantomsection\label{bib:Smith2013}
Smith, Peter (2013). \emph{An Introduction to G\"odel's Theorems}. Cambridge University Press. Cambridge.

\par\medskip\phantomsection\label{bib:Smullyan1968}
Smullyan, Raymond M. (1968). \emph{First-Order Logic}. Springer. New York, NY. Corrected reprint, New York, NY: Dover, 1995.

\par\medskip\phantomsection\label{bib:FefermanLevy1963}
Solomon {Feferman}; Azriel {Levy} (1963). \emph{Independence results in set theory by {Cohen}'s method {II}}. \emph{Notices of the American Mathematical Society}, 10, 593.

\par\medskip\phantomsection\label{bib:Sykes1992}
Sykes, Christopher (1992). \emph{{BBC} {H}orizon: The Strange Life and Death of {D}r.~{T}uring}. \url{https://www.youtube.com/watch?v=gyusnGbBSHE}.

\par\medskip\phantomsection\label{bib:Gentzen1969}
Szabo, Manfred E. (1969). \emph{The Collected Papers of {G}erhard {G}entzen}. मध्ये: \emph{The Collected Papers of {G}erhard {G}entzen}. North-Holland. Amsterdam.

\par\medskip\phantomsection\label{bib:Takeuti2003}
Takeuti, Gaisi; Passell, Nicholas; Yasugi, Mariko (2003). \emph{Memoirs of a Proof Theorist: G\"odel and Other Logicians}. World Scientific. Singapore.

\par\medskip\phantomsection\label{bib:Tamassy1994}
Tamassy, Istvan (1994). \emph{Interview with {R}\'oza {P}\'eter}. \emph{Modern Logic}, 4(3), 277--280.

\par\medskip\phantomsection\label{bib:Sandstrum2019}
Ted Sandstrum (2019). \emph{Mathematical Reasoning: Writing and Proof}. Grand Valley State University. Allendale, MI. \url{https://scholarworks.gvsu.edu/books/7/}.

\par\medskip\phantomsection\label{bib:Theelen2012}
Theelen, Andre (2012). \emph{LEGO Turing Machine}. \url{https://www.youtube.com/watch?v=FTSAiF9AHN4}.

\par\medskip\phantomsection\label{bib:Vitali1905}
Vitali, Giuseppe (1905). \emph{Sul problema della misura dei gruppi di punti di una retta}. Gamberini e Parmeggiani. Bologna.

\par\medskip\phantomsection\label{bib:VonNeumann1925}
von Neumann, John (1925). \emph{Eine {Axiomatisierung} der {Mengenlehre}}. \emph{Journal f\"{u}r die reine und angewandte {Mathematik}}, 154, 219--40.

\par\medskip\phantomsection\label{bib:Wang1990}
Wang, Hao (1990). \emph{Reflections on Kurt G{\"{o}}del}. MIT Press. Cambridge.

\par\medskip\phantomsection\label{bib:Weston2003}
Weston, Tom (2003). \emph{The {Banach}-{Tarski} Paradox}. \url{http://people.math.umass.edu/~weston/oldpapers/banach.pdf}.

\par\medskip\phantomsection\label{bib:WhiteheadRussell1910}
Whitehead, Alfred North; Russell, Bertrand (1910). \emph{Principia Mathematica}. Cambridge University Press. Cambridge.

\par\medskip\phantomsection\label{bib:Zermelo1908}
Zermelo, Ernst (1908). \emph{Untersuchungen {\"u}ber die {G}rundlagen der {M}engenlehre~{I}}. \emph{Mathematische Annalen}, 65(2), 261--281. {E}nglish translation in \hyperref[bib:Ebbinghaus2010]{Ebbinghaus 2010, pp.~189-229}.

\par\medskip\phantomsection\label{bib:Zermelo1908Untersuchungen}
Zermelo, Ernst (1908). \emph{Untersuchungen \"{u}ber die {G}rundlagen der {M}engenlehre {I}}. \emph{Mathematische Annalen}, 65, 261--81.

\par\medskip\phantomsection\label{bib:Zuckerman1973}
Zuckerman, Martin M. (1973). \emph{Formation sequences for propositional formulas}. \emph{Notre Dame Journal of Formal Logic}, 14(1), 134--138.
